% Le fonctionnement des neurones — SVTEEHB Tle TI — Cours signé OnBuch+
% Programme officiel MINESEC. Compiler avec : tectonic N05-fonctionnement-des-neurones.tex
\newcommand{\DOCMATIERE}{SVTEEHB}
\newcommand{\DOCNIVEAU}{Tle TI}
\newcommand{\DOCMODULE}{Module I : Le Monde Vivant}
\newcommand{\DOCLECON}{N05}
\newcommand{\DOCTITRE}{Le fonctionnement des neurones}
% OnBuch+ — préambule « cours signés » ENRICHI (compilé avec tectonic / XeLaTeX)
% Système visuel « Pop » d'OnBuch+ : fond crème (jamais blanc), bordures encre
% épaisses, ombres « dures » décalées, titres Archivo Black en capitales.
% Version MATHÉMATIQUES : graphes (pgfplots/tikz), boîtes pédagogiques
% (définition, propriété, méthode, attention, le savais-tu, exercice résolu,
% exercice, TP, bilan), page de garde et signature OnBuch+ ancrée en bas.
%
% Macros attendues dans main.tex AVANT \input{preamble} :
%   \DOCMATIERE  \DOCNIVEAU  \DOCMODULE  \DOCLECON (numéro)  \DOCTITRE
\documentclass[11pt]{article}

\usepackage{fontspec}
\usepackage[french]{babel}
\usepackage[a4paper,margin=2cm,top=3.4cm,bottom=2.8cm,headheight=1.4cm,footskip=1.3cm]{geometry}
\usepackage{amsmath,amssymb,amsfonts}
\usepackage{mathtools} % \xmapsto, \coloneqq... souvent utilisés par les modèles
\usepackage{cancel} % simplifications barrées (bilans de forces)
% \abs{z} (module, valeur absolue) et \norm{u} (norme), utilisés par les modèles
\providecommand{\abs}{}\let\abs\relax\DeclarePairedDelimiter{\abs}{\lvert}{\rvert}
\providecommand{\norm}{}\let\norm\relax\DeclarePairedDelimiter{\norm}{\lVert}{\rVert}
\usepackage[table]{xcolor} % [table] : fournit aussi \rowcolors (alternance de lignes)
\usepackage{pagecolor}
\usepackage{tikz}
\usetikzlibrary{arrows.meta,positioning,calc,shapes.geometric,shapes.misc,shapes.arrows,decorations.pathmorphing,decorations.pathreplacing,decorations.markings,patterns,shadows,mindmap,trees,fit,backgrounds,matrix,intersections,angles,quotes}
\usepackage{pgfplots}
\usepackage[version=4]{mhchem} % équations nucléaires \ce{^{235}_{92}U}
\usepackage[european,siunitx]{circuitikz} % schémas électriques
\pgfplotsset{compat=1.18}
\usepgfplotslibrary{fillbetween,groupplots} % aires entre courbes (calcul intégral)
\usepackage{enumitem}
\usepackage{fancyhdr}
% [most] charge toutes les bibliothèques tcolorbox (listings, minted, raster,
% documentation...) et gonfle inutilement la mémoire TeX sur un document long
% (~300 boîtes) : on ne charge que ce qui est réellement utilisé.
\usepackage[skins,breakable]{tcolorbox}
\usepackage{graphicx}
\usepackage{colortbl}
\usepackage{booktabs}
\usepackage{tabularx}
\usepackage{multirow}
\usepackage{diagbox} % case d'en-tête diagonale des tableaux à double entrée
% Colonnes extensibles centrée / gauche / droite (C, L, R), souvent utilisées par les modèles
\newcolumntype{C}{>{\centering\arraybackslash}X}
\newcolumntype{L}{>{\raggedright\arraybackslash}X}
\newcolumntype{R}{>{\raggedleft\arraybackslash}X}
\usepackage{array}
\usepackage{multicol}
\usepackage{siunitx}
% parse-numbers=false : accepte aussi \SI{2,5 \times 10^{-3}}{...}
\sisetup{locale=FR,output-decimal-marker={,},group-minimum-digits=4,parse-numbers=false}
\DeclareSIUnit\ohm{\ensuremath{\symup{\Omega}}} % U+2126 absent de la police
\DeclareSIUnit\division{div} % réglages d'oscilloscope (ms/div, V/div)
\DeclareSIUnit\tr{tr} % tours (vitesse de rotation, ex. \SI{1500}{\tr\per\minute})
\DeclareSIUnit\molar{M} % concentration molaire (ex. \SI{3}{\molar}, \SI{1}{\micro\molar})
\DeclareSIUnit\an{an} % année (le modèle confond parfois avec \year, un compteur TeX) — ex. \SI{5}{\kilo\meter\per\an}
\DeclareSIUnit\FCFA{FCFA} % franc CFA (ex. \SI{500}{\FCFA\per\kilo\gram})
\DeclareSIUnit\degreeBrix{\textdegree Brix} % teneur en sucre d'un sirop/jus (réfractométrie)
\DeclareSIUnit\month{mois} % le modèle utilise parfois \month, un compteur TeX, comme unité de durée
\DeclareSIUnit\microliter{\micro\liter} % le modèle écrit parfois \microliter en un seul token
\DeclareSIUnit\million{millions} % dénombrement (ex. \SI{15}{\million\per\mL} spermatozoïdes)
\usepackage{qrcode}
% \degree : commande fréquemment utilisée par les modèles pour un angle ou une
% température en dehors de \SI{}{} (non fournie par les paquets chargés).
\providecommand{\degree}{^{\circ}}
% \qed (fin de démonstration), utilisé par les modèles sans amsthm
\providecommand{\qed}{\hfill$\square$}
% \barre : double barre des valeurs interdites dans un tableau de variations (tkz-tab)
\providecommand{\barre}{\ensuremath{\|}}
% environnement proof (amsthm non chargé) utilisé par les modèles
\ifdefined\proof\else\newenvironment{proof}[1][Démonstration]{\par\noindent\textit{#1.}\ }{\qed\par}\fi
\usepackage{lastpage}
\usepackage{hyperref}
\hypersetup{hidelinks}

% ============================================================
% Palette « Pop » OnBuch+ — mêmes hex que la classe P (plus_ui.dart)
% ============================================================
\definecolor{poporange}{HTML}{F59321}
\definecolor{popdark}{HTML}{E07A0C}
\definecolor{popcream}{HTML}{FFF6E8}
\definecolor{popink}{HTML}{1C1714}
\definecolor{poppurple}{HTML}{7A5AE0}
\definecolor{popgreen}{HTML}{2E9E6A}
\definecolor{poppink}{HTML}{E0457B}
\definecolor{popblue}{HTML}{2F7DE1}
\definecolor{popgold}{HTML}{C98A12}
\definecolor{popmuted}{HTML}{8A7F76}
\definecolor{popline}{HTML}{EADFD2}
\definecolor{popgray}{HTML}{8A7F76}
\colorlet{popgrey}{popgray}
% Alias tolérants (le modèle nomme parfois les couleurs différemment)
\colorlet{popwhite}{white}
\colorlet{popblack}{popink}
\colorlet{popred}{poppink}
\colorlet{popyellow}{popgold}
% Teintes claires (fonds de tableaux/encadrés) — le modèle nomme parfois
% les couleurs avec un suffixe L (ex. popblueL) sans les avoir définies.
\colorlet{poporangeL}{poporange!20}
\colorlet{popdarkL}{popdark!20}
\colorlet{poppurpleL}{poppurple!20}
\colorlet{popgreenL}{popgreen!20}
\colorlet{poppinkL}{poppink!20}
\colorlet{popblueL}{popblue!20}
\colorlet{popgoldL}{popgold!20}
\colorlet{popredL}{popred!20}
\colorlet{popgrayL}{popgray!20}
% Teintes claires pour fonds de tableaux / zones
\colorlet{popblueL}{popblue!10!white}
\colorlet{poporangeL}{poporange!14!white}
\colorlet{popgreenL}{popgreen!12!white}
\colorlet{poppurpleL}{poppurple!12!white}
\colorlet{poppinkL}{poppink!12!white}
\colorlet{popgoldL}{popgold!14!white}

\pagecolor{popcream}

% ============================================================
% Typographie
% ============================================================
\defaultfontfeatures{Ligatures=TeX}
\setmainfont{PlusJakartaSans-Medium.ttf}[Path=../../../fonts/,BoldFont=PlusJakartaSans-Bold.ttf]
% Maths en sans-serif (Fira Math) avec chiffres et lettres droites de Plus
% Jakarta, pour que les formules se fondent dans le texte.
\usepackage[math-style=ISO,bold-style=ISO]{unicode-math}
\setmathfont{FiraMath-Regular.otf}[Path=../../../fonts/]
\setmathfont{PlusJakartaSans-Medium.ttf}[Path=../../../fonts/,range={up/num,up/{Latin,latin},\mathup}]
\usepackage{icomma} % virgule décimale sans espace en mode math
% Caractères mathématiques Unicode absents des polices de texte (Plus Jakarta,
% Archivo Black) : rendus par leur équivalent mathématique, en texte comme en
% formule — sinon ils s'affichent en carré vide (ex. « Arithmétique dans ℤ »).
\usepackage{newunicodechar}
\newunicodechar{ℝ}{\ensuremath{\mathbb{R}}}
\newunicodechar{ℤ}{\ensuremath{\mathbb{Z}}}
\newunicodechar{ℂ}{\ensuremath{\mathbb{C}}}
\newunicodechar{ℕ}{\ensuremath{\mathbb{N}}}
\newunicodechar{ℚ}{\ensuremath{\mathbb{Q}}}
\newunicodechar{𝔻}{\ensuremath{\mathbb{D}}}
\newunicodechar{α}{\ensuremath{\alpha}}
\newunicodechar{β}{\ensuremath{\beta}}
\newunicodechar{γ}{\ensuremath{\gamma}}
\newunicodechar{δ}{\ensuremath{\delta}}
\newunicodechar{ε}{\ensuremath{\varepsilon}}
\newunicodechar{θ}{\ensuremath{\theta}}
\newunicodechar{λ}{\ensuremath{\lambda}}
\newunicodechar{μ}{\ensuremath{\mu}}
\newunicodechar{σ}{\ensuremath{\sigma}}
\newunicodechar{τ}{\ensuremath{\tau}}
\newunicodechar{φ}{\ensuremath{\varphi}}
\newunicodechar{ψ}{\ensuremath{\psi}}
\newunicodechar{ω}{\ensuremath{\omega}}
\newunicodechar{Σ}{\ensuremath{\Sigma}}
% majuscules produites par \MakeUppercase dans les titres (α-aminés -> Α)
\newunicodechar{Α}{\ensuremath{\alpha}}
\newunicodechar{Β}{\ensuremath{\beta}}
\newunicodechar{Γ}{\ensuremath{\Gamma}}
\newunicodechar{Δ}{\ensuremath{\Delta}}
\newunicodechar{Ε}{\ensuremath{\varepsilon}}
\newunicodechar{Θ}{\ensuremath{\Theta}}
\newunicodechar{Λ}{\ensuremath{\Lambda}}
\newunicodechar{Μ}{\ensuremath{\mu}}
\newunicodechar{Τ}{\ensuremath{\tau}}
\newunicodechar{Φ}{\ensuremath{\Phi}}
\newunicodechar{Ψ}{\ensuremath{\Psi}}
\newunicodechar{Ω}{\ensuremath{\Omega}}
\newunicodechar{↦}{\ensuremath{\mapsto}}
\newunicodechar{∈}{\ensuremath{\in}}
\newunicodechar{∉}{\ensuremath{\notin}}
\newunicodechar{∧}{\ensuremath{\wedge}}
\newunicodechar{⇔}{\ensuremath{\Leftrightarrow}}
\newunicodechar{⟺}{\ensuremath{\Longleftrightarrow}}
\newunicodechar{⇒}{\ensuremath{\Rightarrow}}
\newunicodechar{⟹}{\ensuremath{\Longrightarrow}}
\newunicodechar{∘}{\ensuremath{\circ}}
\newunicodechar{∪}{\ensuremath{\cup}}
\newunicodechar{∩}{\ensuremath{\cap}}
\newunicodechar{≡}{\ensuremath{\equiv}}
\newunicodechar{⊂}{\ensuremath{\subset}}
\newunicodechar{⊥}{\ensuremath{\perp}}
\newunicodechar{∃}{\ensuremath{\exists}}
\newunicodechar{∀}{\ensuremath{\forall}}
\newunicodechar{∛}{\ensuremath{\sqrt[3]{\,}}}
\newunicodechar{∜}{\ensuremath{\sqrt[4]{\,}}}
\newunicodechar{⃗}{\ensuremath{{}^{\to}}}
\newunicodechar{ⁿ}{\ifmmode^{n}\else\textsuperscript{\ensuremath{n}}\fi}
\newunicodechar{ˣ}{\ifmmode^{x}\else\textsuperscript{\ensuremath{x}}\fi}
\newunicodechar{ᵇ}{\ifmmode^{b}\else\textsuperscript{\ensuremath{b}}\fi}
\newunicodechar{ʳ}{\ifmmode^{r}\else\textsuperscript{\ensuremath{r}}\fi}
\newunicodechar{ᵘ}{\ifmmode^{u}\else\textsuperscript{\ensuremath{u}}\fi}
\newunicodechar{ʸ}{\ifmmode^{y}\else\textsuperscript{\ensuremath{y}}\fi}
\newunicodechar{ᵃ}{\ifmmode^{a}\else\textsuperscript{\ensuremath{a}}\fi}
\newunicodechar{ᵉ}{\ifmmode^{e}\else\textsuperscript{\ensuremath{e}}\fi}
\newunicodechar{ᵏ}{\ifmmode^{k}\else\textsuperscript{\ensuremath{k}}\fi}
\newunicodechar{ᵖ}{\ifmmode^{p}\else\textsuperscript{\ensuremath{p}}\fi}
\newunicodechar{ᴺ}{\ifmmode^{N}\else\textsuperscript{\ensuremath{N}}\fi}
\newunicodechar{ᶜ}{\ifmmode^{c}\else\textsuperscript{\ensuremath{c}}\fi}
\newunicodechar{ʰ}{\ifmmode^{h}\else\textsuperscript{\ensuremath{h}}\fi}
\newunicodechar{ᵠ}{\ifmmode^{q}\else\textsuperscript{\ensuremath{q}}\fi}
\newunicodechar{ₙ}{\ifmmode_{n}\else\textsubscript{\ensuremath{n}}\fi}
\newunicodechar{ₐ}{\ifmmode_{a}\else\textsubscript{\ensuremath{a}}\fi}
\newunicodechar{ᵢ}{\ifmmode_{i}\else\textsubscript{\ensuremath{i}}\fi}
\newunicodechar{ᵣ}{\ifmmode_{r}\else\textsubscript{\ensuremath{r}}\fi}
\newunicodechar{ₖ}{\ifmmode_{k}\else\textsubscript{\ensuremath{k}}\fi}
\newunicodechar{ᵦ}{\ifmmode_{\beta}\else\textsubscript{\ensuremath{\beta}}\fi}
\newunicodechar{ₓ}{\ifmmode_{x}\else\textsubscript{\ensuremath{x}}\fi}
\newfontfamily\popdisplay{ArchivoBlack-Regular.ttf}[Path=../../../fonts/]
\newfontfamily\poplogo{SpaceGrotesk-Bold.ttf}[Path=../../../fonts/]
\newfontfamily\popsemi{PlusJakartaSans-SemiBold.ttf}[Path=../../../fonts/]
\newfontfamily\popxbold{PlusJakartaSans-ExtraBold.ttf}[Path=../../../fonts/]
\newfontfamily\popxboldls{PlusJakartaSans-ExtraBold.ttf}[Path=../../../fonts/,LetterSpace=3.0]

\setlist[enumerate]{leftmargin=1.7em,itemsep=5pt,topsep=4pt}
\setlist[itemize]{leftmargin=1.7em,itemsep=4pt,topsep=4pt,label={\color{poporange}\textbullet}}
\setlength{\parindent}{0pt}
\setlength{\parskip}{5pt}
\renewcommand{\arraystretch}{1.25}

% pgfplots : style Pop par défaut
\pgfplotsset{
  every axis/.append style={axis line style={popink,line width=1pt},tick style={popink},
    grid style={popline,line width=0.5pt},label style={font=\small},tick label style={font=\footnotesize}},
  popcurve/.style={poporange,line width=1.8pt,no markers,smooth},
}
% Même style disponible dans un \draw TikZ simple (hors pgfplots)
\tikzset{popcurve/.style={poporange,line width=1.8pt}}
\tikzset{popgrid/.style={popline,very thin}}
\colorlet{popcurve}{poporange} % le nom du style est parfois utilisé comme couleur

% ============================================================
% Pill / squiggle
% ============================================================
\newcommand{\poppill}[3]{%
  \tikz[baseline=-0.55ex]\node[draw=popink,line width=1.2pt,rounded corners=2.6pt,%
    fill=#2,inner xsep=6.5pt,inner ysep=3pt]{%
    \popxboldls\fontsize{7}{7}\selectfont\color{#3}\MakeUppercase{#1}};%
}
\newcommand{\popsquiggle}[1]{%
  \tikz[baseline=-0.4ex]{
    \draw[#1,line width=2.6pt,line cap=round]
      plot [smooth] coordinates {(0,0.09) (0.34,0.01) (0.68,0.21) (1.02,0.01) (1.36,0.10)};
  }%
}
% Mot-clé surligné (marqueur orange sous le texte)
\newcommand{\cle}[1]{{\popxbold\color{popdark}#1}}

% ============================================================
% En-tête / pied
% ============================================================
\pagestyle{fancy}
\fancyhf{}
\renewcommand{\headrulewidth}{0pt}
\renewcommand{\footrulewidth}{0pt}
\lhead{\raisebox{-0.15ex}{\poplogo\fontsize{13}{13}\selectfont\color{popink}OnBuch\color{poporange}+}}
\rhead{\poppill{\DOCMATIERE\ \textbullet\ \DOCNIVEAU\ \textbullet\ Leçon \DOCLECON}{white}{popink}}
\lfoot{\popxbold\fontsize{7.6}{7.6}\selectfont\color{popmuted}\MakeUppercase{Cours signé OnBuch+}\ \textbullet\ {\popsemi\DOCTITRE}}
\rfoot{\tikz[baseline=-0.6ex]\node[rounded corners=8.5pt,fill=popink,minimum height=17pt,minimum width=17pt,inner xsep=5pt,inner sep=0]{\popxbold\fontsize{8}{8}\selectfont\color{popcream}\thepage\,/\,\pageref*{LastPage}};}

% ============================================================
% Titres
% ============================================================
\newcommand{\chaptitre}[2]{%
  \begin{tcolorbox}[enhanced,colback=poporange,colframe=popink,boxrule=2.6pt,arc=11pt,
      drop shadow=popink,left=15pt,right=15pt,top=12pt,bottom=11pt,before skip=3pt,after skip=18pt]
    \poppill{#2}{popink}{popcream}\par\vspace{9pt}
    {\raggedright\hyphenpenalty=10000\exhyphenpenalty=10000\popdisplay\fontsize{22}{24}\selectfont\color{popcream}\MakeUppercase{#1}\par}
  \end{tcolorbox}
}
% Section numérotée : pastille orange avec le numéro + titre Archivo
\newcounter{popsec}
\newcommand{\coursec}[1]{%
  \stepcounter{popsec}\setcounter{popsub}{0}%
  \par\vspace{16pt}\noindent
  \begin{minipage}{\linewidth}
  \tikz[baseline=-0.7ex]\node[draw=popink,line width=1.6pt,fill=poporange,rounded corners=4pt,minimum size=22pt,inner sep=0]{\popdisplay\fontsize{12}{12}\selectfont\color{popcream}\thepopsec};%
  \hspace{8pt}{\popdisplay\fontsize{15}{17}\selectfont\color{popink}\MakeUppercase{#1}}\par
  \vspace{4pt}\noindent{\color{poporange}\rule{36pt}{3.6pt}}
  \end{minipage}\par\nopagebreak\vspace{8pt}
}
\newcounter{popsub}
\newcommand{\courssub}[1]{%
  \stepcounter{popsub}%
  \par\vspace{9pt}\noindent{\popxbold\fontsize{11.5}{13}\selectfont\color{popink}{\color{poporange}\thepopsec.\thepopsub}\hspace{6pt}#1}\par\nopagebreak\vspace{4pt}
}

% ============================================================
% Boîtes pédagogiques — même recette : fond blanc, bordure épaisse colorée,
% ombre dure encre, badge plein. Titre optionnel [#1].
% ============================================================
\tcbset{popbase/.style={breakable,enhanced,colback=white,boxrule=2.3pt,arc=9pt,
  shadow={3pt}{-3pt}{0pt}{popink},left=14pt,right=14pt,top=11pt,bottom=11pt,before skip=11pt,after skip=15pt,
  fontupper=\popsemi\fontsize{10.6}{14.5}\selectfont\color{popink},
  fontlower=\popsemi\fontsize{10.6}{14.5}\selectfont\color{popink}}}
% Coche dessinée (le glyphe ✓ n'existe pas dans les polices du document)
\newcommand{\popcheck}[1]{\tikz[baseline=-0.5ex]\draw[#1,line width=1.6pt,line cap=round,line join=round](-0.13,0.0)--(-0.04,-0.09)--(0.13,0.1);}
\providecommand{\coche}{\popcheck{popgreen}}
\providecommand{\resultat}[1]{\par\smallskip\noindent\fcolorbox{popgreen}{popgreenL}{\parbox{\dimexpr\linewidth-2\fboxsep-2\fboxrule\relax}{\textbf{Résultat.} #1}}\par\smallskip}
\newcommand{\popboxhead}[3]{\poppill{#1}{#2}{white}\ifx\relax#3\relax\else\hspace{6pt}{\popxbold\fontsize{10.6}{12}\selectfont\color{popink}#3}\fi\par\vspace{7pt}}

\newtcolorbox{aretenir}[1][]{popbase,colframe=popgreen,before upper={\popboxhead{À retenir}{popgreen}{#1}}}
\newtcolorbox{exemplebox}[1][]{popbase,colframe=poporange,before upper={\popboxhead{Exemple}{poporange}{#1}}}
\newtcolorbox{definition}[1][]{popbase,colframe=popblue,before upper={\popboxhead{Définition}{popblue}{#1}}}
\newtcolorbox{propriete}[1][]{popbase,colframe=poppurple,before upper={\popboxhead{Propriété}{poppurple}{#1}}}
\newtcolorbox{methode}[1][]{popbase,colframe=popgold,before upper={\popboxhead{Méthode}{popgold}{#1}}}
\newtcolorbox{attention}[1][]{popbase,colframe=poppink,before upper={\popboxhead{Attention}{poppink}{#1}}}
\newtcolorbox{experience}[1][]{popbase,colframe=popblue,colback=popblueL,before upper={\popboxhead{Expérience}{popblue}{#1}}}
% « Le savais-tu ? » : carte violette pleine (contexte, culture, Cameroun)
\newtcolorbox{savaistu}[1][]{popbase,colback=poppurple,colframe=popink,
  fontupper=\popsemi\fontsize{10.4}{14.2}\selectfont\color{white},
  before upper={\poppill{Le savais-tu\,?}{white}{poppurple}\ifx\relax#1\relax\else\hspace{6pt}{\popxbold\color{white}#1}\fi\par\vspace{7pt}}}
% Exercice résolu : énoncé en haut, \tcblower puis la solution sur fond crème
\newcounter{popexo}\newcounter{popexr}
\newtcolorbox{exoresolu}[1][]{popbase,colframe=poporange,colbacklower=popcream,
  segmentation style={popink,line width=1.2pt,dashed},
  before upper={\stepcounter{popexr}\popboxhead{Exercice résolu \thepopexr}{poporange}{#1}},
  before lower={\poppill{Solution}{popink}{popcream}\par\vspace{6pt}}}
% Exercice (non résolu en place) — la correction est dans la partie Corrigés
\newtcolorbox{exercice}[1][]{popbase,colframe=popink,
  before upper={\stepcounter{popexo}\popboxhead{Exercice \thepopexo}{popink}{#1}}}
% Corrigé
\newtcolorbox{corrige}[1][]{popbase,colframe=popgreen,colback=popgreenL,
  before upper={\popboxhead{Corrigé}{popgreen}{#1}}}
% Objectifs / prérequis / bilan
\newtcolorbox{objectifs}{popbase,colframe=popink,colback=poporangeL,
  before upper={\popboxhead{Objectifs de la leçon}{popink}{}}}
\newtcolorbox{prerequis}{popbase,colframe=popink,colback=popblueL,
  before upper={\popboxhead{Prérequis}{popblue}{}}}
\newtcolorbox{bilan}{popbase,colframe=popink,colback=white,boxrule=2.8pt,
  before upper={\popboxhead{Fiche bilan}{poporange}{L'essentiel à connaître pour l'examen}}}
% Figure encadrée avec légende
\newtcolorbox{popfigure}[1][]{enhanced,breakable=false,colback=white,colframe=popline,boxrule=1.4pt,arc=8pt,
  left=10pt,right=10pt,top=10pt,bottom=8pt,before skip=10pt,after skip=12pt,halign=center,
  lower separated=false,halign lower=center,
  fontlower=\popsemi\fontsize{9}{11.5}\selectfont\color{popmuted},
  #1}
\newcommand{\legende}[1]{\par\vspace{6pt}{\popsemi\fontsize{9}{11.5}\selectfont\color{popmuted}#1}}

% ============================================================
% Signature OnBuch+
% ============================================================
\newcommand{\poplogoinline}[1]{{\poplogo\fontsize{#1}{#1}\selectfont\color{popink}OnBuch\color{poporange}+}}
\newcommand{\signatureonbuch}{%
  \begin{tcolorbox}[enhanced,colback=popink,colframe=popink,boxrule=2.2pt,arc=10pt,
      drop shadow=poporange,left=14pt,right=14pt,top=10pt,bottom=10pt,before skip=0pt,after skip=0pt]
    \tikz[baseline=-0.7ex]\node[circle,fill=popgreen,draw=popcream,line width=1.3pt,minimum size=22pt,inner sep=0]{\popcheck{white}};%
    \hspace{9pt}\begin{minipage}[c]{0.62\linewidth}
      {\popxbold\fontsize{12}{13}\selectfont\color{popcream}Signé\ \textbullet\ {\poplogo OnBuch\color{poporange}+}}\par\vspace{2pt}
      {\raggedright\popsemi\fontsize{8}{10.5}\selectfont\color{popline}\MakeUppercase{Équipe pédagogique OnBuch+}\\\MakeUppercase{Conforme au programme officiel MINESEC}\par}
    \end{minipage}\hfill
    {\poplogo\fontsize{22}{22}\selectfont\color{popcream}OnBuch\color{poporange}+}
  \end{tcolorbox}%
}

% ============================================================
% Page de garde
% #1 titre · #2 matière · #3 niveau · #4 module · #5 n° leçon · #6 durée ·
% #7 description / objectifs (texte ou itemize)
% ============================================================
\newcommand{\pagegarde}[7]{%
  \thispagestyle{empty}
  \newgeometry{a4paper,margin=1.7cm,top=1.6cm,bottom=1.4cm}%
  \begin{tikzpicture}[remember picture,overlay]
    \fill[poporange!18] ([xshift=-3.2cm,yshift=-3.6cm]current page.north east) circle (4.6cm);
    \fill[poppurple!14] ([xshift=2.4cm,yshift=7.6cm]current page.south west) circle (3.8cm);
  \end{tikzpicture}%
  {\poplogo\fontsize{30}{30}\selectfont\color{popink}OnBuch\color{poporange}+}\hfill
  \raisebox{0.6ex}{\poppill{Cours signé \textbullet\ Premium}{popink}{popcream}}\par
  \vspace{1pt}\popsquiggle{poppurple}\par
  \vspace{8pt}
  \begin{tcolorbox}[enhanced,colback=poporange,colframe=popink,boxrule=2.8pt,arc=13pt,
      drop shadow=popink,left=18pt,right=18pt,top=12pt,bottom=13pt,after skip=10pt]
    \poppill{#2 \textbullet\ #3}{popink}{popcream}\hspace{5pt}\poppill{#4}{white}{popink}\par\vspace{8pt}
    {\popdisplay\fontsize{14}{14}\selectfont\color{popink}LEÇON #5}\par\vspace{4pt}
    {\raggedright\hyphenpenalty=10000\exhyphenpenalty=10000\popdisplay\fontsize{25}{27}\selectfont\color{popcream}\MakeUppercase{#1}\par}
    \vspace{8pt}
    {\popxbold\fontsize{9.5}{9.5}\selectfont\color{popink}\MakeUppercase{Durée conseillée : #6}}
  \end{tcolorbox}
  \begin{tcolorbox}[enhanced,colback=white,colframe=popink,boxrule=2.2pt,arc=10pt,
      drop shadow=popink,left=16pt,right=16pt,top=10pt,bottom=10pt,after skip=8pt]
    \poppill{Dans cette leçon}{poppurple}{white}\par\vspace{5pt}
    \popsemi\fontsize{9.3}{12.6}\selectfont\color{popink}#7
  \end{tcolorbox}
  \noindent\begin{minipage}[t]{0.487\linewidth}
    \begin{tcolorbox}[enhanced,colback=popgreen,colframe=popink,boxrule=2.2pt,arc=10pt,
        drop shadow=popink,left=11pt,right=11pt,top=10pt,bottom=10pt,height=3.1cm,valign=center]
      \tcbox[colback=white,colframe=popink,boxrule=1.3pt,arc=5pt,boxsep=0pt,nobeforeafter,
        left=3pt,right=3pt,top=3pt,bottom=3pt,valign=center]{\qrcode[height=1.9cm]{https://play.google.com/store/apps/details?id=com.luvvix.onbuch&hl=fr}}%
      \hspace{8pt}\begin{minipage}[c]{0.5\linewidth}
        {\popdisplay\fontsize{11}{12.5}\selectfont\color{white}\MakeUppercase{Télécharge OnBuch}}\par\vspace{4pt}
        {\raggedright\popsemi\fontsize{8.4}{11}\selectfont\color{white}Gratuit \textbullet\ Google Play\\Tuteur IA, annales, concours, résultats\par}
      \end{minipage}
    \end{tcolorbox}
  \end{minipage}\hfill
  \begin{minipage}[t]{0.487\linewidth}
    \begin{tcolorbox}[enhanced,colback=poppurple,colframe=popink,boxrule=2.2pt,arc=10pt,
        drop shadow=popink,left=11pt,right=11pt,top=10pt,bottom=10pt,height=3.1cm,valign=center]
      \tikz[baseline=-0.7ex]\node[circle,fill=white,draw=popink,line width=1.3pt,minimum size=1.6cm,inner sep=0]{\popdisplay\fontsize{24}{24}\selectfont\color{poppurple}+};%
      \hspace{8pt}\begin{minipage}[c]{0.6\linewidth}
        {\popdisplay\fontsize{11}{12.5}\selectfont\color{white}\MakeUppercase{Passe à OnBuch+}}\par\vspace{4pt}
        {\raggedright\popsemi\fontsize{8.4}{11}\selectfont\color{white}Tous les cours signés, Tuteur IA illimité, fiches PDF — onglet OnBuch+ de l'app\par}
      \end{minipage}
    \end{tcolorbox}
  \end{minipage}
  \vspace{8pt}
  \popcontenu
  \vfill
  \signatureonbuch
  \restoregeometry
  \newpage
}

% Rangée « Ce cours contient » (page de garde)
\newcommand{\poptile}[3]{% #1 couleur #2 gros texte #3 légende
  \begin{tcolorbox}[enhanced,colback=#1,colframe=popink,boxrule=2pt,arc=9pt,shadow={2.5pt}{-2.5pt}{0pt}{popink},
      left=6pt,right=6pt,top=9pt,bottom=9pt,halign=center,valign=center,height=1.75cm,width=\linewidth,nobeforeafter]
    {\popdisplay\fontsize{12.5}{13}\selectfont\color{white}\MakeUppercase{#2}}\par\vspace{3pt}
    {\centering\popsemi\fontsize{7.8}{9.5}\selectfont\color{white}#3\par}
  \end{tcolorbox}}
\newcommand{\popcontenu}{%
  \par\noindent{\popxboldls\fontsize{7.5}{7.5}\selectfont\color{popmuted}CE COURS SIGNÉ CONTIENT}\par\vspace{7pt}
  \noindent\begin{minipage}[t]{0.235\linewidth}\poptile{popblue}{Cours}{complet, illustré, pas à pas}\end{minipage}\hfill
  \begin{minipage}[t]{0.235\linewidth}\poptile{poporange}{Méthodes}{et exercices résolus}\end{minipage}\hfill
  \begin{minipage}[t]{0.235\linewidth}\poptile{poppink}{Exercices}{type examen + corrigés}\end{minipage}\hfill
  \begin{minipage}[t]{0.235\linewidth}\poptile{popgreen}{Fiche bilan}{l'essentiel à retenir}\end{minipage}\par}

% Sommaire
\newtcolorbox{sommaire}{enhanced,colback=white,colframe=popink,boxrule=2.2pt,arc=10pt,
  drop shadow=popink,left=18pt,right=18pt,top=13pt,bottom=13pt,before skip=6pt,after skip=20pt,
  before upper={\poppill{Sommaire}{poppurple}{white}\par\vspace{9pt}\popsemi\fontsize{10.6}{14.5}\selectfont\color{popink}}}

% Signature de fin de cours — poussée tout en bas de la dernière page
\newcommand{\findecours}{%
  \par\vfill\nopagebreak
  \begin{center}\popsquiggle{poporange}\end{center}\vspace{4pt}
  \signatureonbuch
}
\AtBeginDocument{\def\llbracket{\mathopen{[\mkern-3.5mu[}}\def\rrbracket{\mathclose{]\mkern-3.5mu]}}} % intervalles d'entiers (glyphe ⟦ absent de la police)
% Glyphes absents de Fira Math : redéfinis à partir de glyphes disponibles
\AtBeginDocument{%
  \def\setminus{\mathbin{\backslash}}\let\smallsetminus\setminus
  \def\vdots{\mathord{\vbox{\baselineskip3.2pt\lineskiplimit0pt\kern4pt\hbox{.}\hbox{.}\hbox{.}}}}%
  \def\ddots{\mathinner{\mkern1mu\raise6.5pt\hbox{.}\mkern2mu\raise3.6pt\hbox{.}\mkern2mu\raise0.7pt\hbox{.}\mkern1mu}}%
  \def\triangle{\mathord{\scalebox{0.9}{$\symup{\Delta}$}}}%
  \def\nabla{\mathord{\raisebox{\depth}{\scalebox{1}[-1]{$\symup{\Delta}$}}}}%
  \def\checkmark{\popcheck{popdark}}%
  \def\star{\mathbin{\ast}}\def\bigstar{\mathord{\ast}}%
  \def\diamond{\mathbin{\scalebox{0.6}{$\Diamond$}}}%
  \def\bigcup{\mathop{\mathchoice{\raisebox{-0.3em}{\scalebox{1.7}{$\cup$}}}{\scalebox{1.25}{$\cup$}}{\cup}{\cup}}\displaylimits}%
  \def\bigcap{\mathop{\mathchoice{\raisebox{-0.3em}{\scalebox{1.7}{$\cap$}}}{\scalebox{1.25}{$\cap$}}{\cap}{\cap}}\displaylimits}%
  \def\lesseqqgtr{\mathrel{\lessgtr}}\def\bigoplus{\mathop{\oplus}\displaylimits}%
}
\providecommand{\ssi}{si et seulement si} % abréviation fréquente
\AtBeginDocument{\providecommand{\wideparen}{\overparen}} % arc de cercle

%% FIGFIX-BEGIN v4 — correctifs globaux des figures (docs/onbuch-generation/tools/figfix)
\usepackage{adjustbox}\usepackage{etoolbox}
% toute figure TikZ/pgfplots plus large que la ligne est réduite à la largeur disponible
\BeforeBeginEnvironment{tikzpicture}{\begin{adjustbox}{max width=\linewidth}}
\AfterEndEnvironment{tikzpicture}{\end{adjustbox}}
% légendes pgfplots lisibles par-dessus les courbes
\pgfplotsset{every axis/.append style={legend style={fill=white,fill opacity=0.92,text opacity=1,draw=black!15,inner sep=3pt}}}
% étiquettes d'arêtes (syntaxe edge["..."]) avec fond blanc
\tikzset{every edge quotes/.append style={fill=white,inner sep=1.2pt}}
%% FIGFIX-END
\begin{document}

\pagegarde{Le fonctionnement des neurones}{SVTEEHB}{Tle TI}{Module I : Le Monde Vivant}{N05}{5 h}{Cette leçon étudie les mécanismes électriques et chimiques qui permettent aux neurones de générer, conduire et transmettre le message nerveux. Elle s'appuie sur l'exploitation d'expériences historiques et de courbes d'enregistrement pour établir les notions de potentiel de repos, de potentiel d'action et de transmission synaptique.
\vspace{4pt}
\begin{multicols}{2}\small
\begin{itemize}
\item À la fin de cette leçon, je sais observer et identifier les centres nerveux d'un mammifère
\item À la fin de cette leçon, je sais schématiser un neurone et légender son schéma
\item À la fin de cette leçon, je sais décrire le protocole de mesure des potentiels de repos et d'action
\item À la fin de cette leçon, je sais représenter et interpréter la courbe du potentiel d'action
\item À la fin de cette leçon, je sais expliquer la conduction du message nerveux le long d'une fibre nerveuse et les facteurs de variation de sa vitesse
\item À la fin de cette leçon, je sais définir une synapse et distinguer les différents types de synapses
\end{itemize}\end{multicols}}

\begin{sommaire}
\begin{multicols}{2}
\begin{enumerate}
\item Le tissu nerveux et le neurone (rappels)
\item Le potentiel de repos
\item Le potentiel d'action
\item La conduction du message nerveux et sa vitesse
\item La synapse : types et fonctionnement de la synapse chimique
\item Effets de quelques substances sur la transmission synaptique
\item Activité d'intégration
\item Méthodes et exercices résolus
\item Exercices
\item Corrigés des exercices
\item Fiche bilan
\end{enumerate}
\end{multicols}
\end{sommaire}

\begin{objectifs}
\begin{itemize}
\item À la fin de cette leçon, je sais observer et identifier les centres nerveux d'un mammifère
\item À la fin de cette leçon, je sais schématiser un neurone et légender son schéma
\item À la fin de cette leçon, je sais décrire le protocole de mesure des potentiels de repos et d'action
\item À la fin de cette leçon, je sais représenter et interpréter la courbe du potentiel d'action
\item À la fin de cette leçon, je sais expliquer la conduction du message nerveux le long d'une fibre nerveuse et les facteurs de variation de sa vitesse
\item À la fin de cette leçon, je sais définir une synapse et distinguer les différents types de synapses
\item À la fin de cette leçon, je sais schématiser une synapse à transmission chimique et expliquer son fonctionnement
\item À la fin de cette leçon, je sais expliquer les effets de quelques substances (curare, atropine, drogues) sur la transmission synaptique
\end{itemize}
\end{objectifs}

\begin{prerequis}
\begin{itemize}
\item Structure générale du système nerveux (cerveau, moelle épinière, nerfs) vue en classe de Première
\item Notion de tissu et d'organisation cellulaire (Tle, début du Module I)
\item Structure de la membrane cellulaire (bicouche phospholipidique, protéines membranaires)
\item Notions d'électricité de base : différence de potentiel, courant électrique (Physique)
\item Réflexes et arc réflexe (notions de Seconde/Première)
\end{itemize}
\end{prerequis}

\newpage

\coursec{Le tissu nerveux et le neurone (rappels)}

Lors d'un match de football au stade de Yaoundé, un joueur reçoit un coup violent à la jambe : en une fraction de seconde, il retire sa jambe et crie de douleur. Comment l'information a-t-elle voyagé si vite de son pied jusqu'à son cerveau ? Pourquoi certains anesthésiants utilisés à l'hôpital bloquent-ils cette douleur, et comment certaines drogues perturbent-elles le fonctionnement de nos nerfs ? Comprendre le fonctionnement des neurones permet de répondre à ces questions du quotidien.

\courssub{Observation des centres nerveux : substance grise et substance blanche}

\begin{experience}[Observation macroscopique d'un centre nerveux (cerveau, cervelet, moelle épinière de mouton ou de bœuf)]
\textbf{Matériel :} Moelle épinière fraîche de mouton (abattoir de Ngaoundéré ou Douala), lame, pince, loupe binoculaire, bac de dissection, eau physiologique.
\textbf{Protocole :}
\begin{enumerate}
    \item Placer la moelle épinière dans le bac de dissection, face dorsale vers le haut.
    \item À l'aide de la lame, réaliser une coupe transversale nette au niveau de la région lombaire.
    \item Observer la coupe à la loupe binoculaire (grossissement $\times 10$ à $\times 40$).
    \item Dessiner ce que vous observez en distinguant les zones de couleur différente.
\end{enumerate}
\textbf{Observations :} La coupe transversale montre une région centrale en forme de papillon (ou de H), de couleur \textbf{gris foncé}, entourée d'une région plus large, de couleur \textbf{blanc nacré}.
\textbf{Interprétation :} La région centrale, riche en corps cellulaires de neurones et en cellules gliales, constitue la \textbf{substance grise}. La région périphérique, riche en axones gainés de myéline (lipides blancs), constitue la \textbf{substance blanche}.
\textbf{Conclusion :} Les centres nerveux (encéphale, moelle épinière) sont organisés en substance grise (corps cellulaires, zones de traitement de l'information) et substance blanche (fibres nerveuses, zones de conduction de l'information).
\legende{Coupe transversale de la moelle épinière d'un mammifère : organisation en substance grise (centrale) et substance blanche (périphérique).}
\end{experience}

\begin{popfigure}
\begin{tikzpicture}[scale=1.2]
% White matter background
\fill[popblue!10] (-3,-2.5) rectangle (3,2.5);
% Grey matter (butterfly shape)
\fill[popgray!60] (-0.5,-2) .. controls (-1.5,-1) and (-1.5,1) .. (-0.5,0) .. controls (-1.5,-1) .. (-0.5,-2); % Left horn
\fill[popgray!60] (0.5,-2) .. controls (1.5,-1) and (1.5,1) .. (0.5,0) .. controls (1.5,-1) .. (0.5,-2); % Right horn
\fill[popgray!60] (-0.5,-0.2) rectangle (0.5,0.2); % Central commissure
% Central canal
\fill[popblue!50] (-0.1,-0.1) rectangle (0.1,0.1);
% Dorsal root ganglion (left)
\draw[fill=poporange!30, thick] (-2.5,1.5) ellipse (0.4 and 0.3);
% Ventral root (right)
\draw[thick, ->] (2.5, -1.5) -- (3.2, -1.5);
% Labels
\node[align=center] at (0, 2.2) {\textbf{Substance blanche} (axones myélinisés)};
\node[align=center, text=popdark] at (-1.5, 0) {\textbf{Corne dorsale}};
\node[align=center, text=popdark] at (1.5, 0) {\textbf{Corne ventrale}};
\node[align=center, text=popdark] at (0, -1.5) {\textbf{Commissure grise}};
\node[align=center, text=popblue] at (0, 0) {\textbf{Canal central}};
\node[align=center] at (-2.5, 2.1) {Ganglion rachidien};
\node[align=center] at (3.3, -1.5) {Racine ventrale};
% Arrows for roots
\draw[thick, ->] (-2.5, 1.2) -- (-1.5, 0.5);
\draw[thick, ->] (-1.5, -0.5) -- (-2.5, -1.2);
\end{tikzpicture}
\legende{Coupe transversale schématique de la moelle épinière d'un mammifère. La substance grise (corps cellulaires, noyaux) forme un papillon central autour du canal central. La substance blanche (faisceaux d'axones myélinisés) l'entoure. Les racines dorsales (sensitives) et ventrales (motrices) émergent latéralement.}
\end{popfigure}

\begin{definition}[Tissu nerveux]
Le \textbf{tissu nerveux} est un tissu animal spécialisé dans la réception, l'intégration et la transmission de l'information. Il est constitué de deux types cellulaires principaux :
\begin{itemize}
    \item Les \textbf{neurones} (ou cellules nerveuses) : cellules excitées, unités structurales et fonctionnelles, assurant la conduction de l'influx nerveux.
    \item Les \textbf{cellules gliales} (ou neuroglie) : cellules non excitables, plus nombreuses que les neurones (ratio $\approx 1,5:1$ chez l'humain), assurant le support structural, la nutrition, l'isolement électrique (myélinisation), la défense immunitaire et l'homéostasie ionique.
\end{itemize}
\end{definition}

\begin{savaistu}[Les cellules gliales, des partenaires indispensables]
Au Cameroun, comme partout, la recherche sur les gliomes (tumeurs des cellules gliales) mobilise les services de neurochirurgie des hôpitaux centraux (Yaoundé, Douala). Les astrocytes régulent la concentration en $K^+$ et recaptent les neurotransmetteurs ; les oligodendrocytes (SNC) et les cellules de Schwann (SNP) fabriquent la myéline ; la microglie assure la phagocytose. Sans glie, pas de conduction rapide ni de survie neuronale !
\end{savaistu}

\courssub{Structure d'un neurone : l'unité morphofonctionnelle}

\begin{definition}[Neurone]
Le \textbf{neurone} est une cellule nerveuse hautement différenciée, spécialisée dans la réception, la conduction et la transmission du message nerveux (influx nerveux). Il est \emph{amitotique} (ne se divise plus chez l'adulte) et présente une polarité morphologique et fonctionnelle marquée.
\end{definition}

\begin{popfigure}
\begin{tikzpicture}[scale=0.9, every node/.style={font=\small}]
% Soma (cell body)
\draw[thick, fill=poporange!20] (0,0) ellipse (1.2 and 0.9);
% Nucleus
\draw[thick, fill=poporange!50] (0,0) circle (0.35);
\node at (0,0) {\textbf{Noyau}};
% Nucleolus
\draw[fill=popdark] (0.1,0.1) circle (0.08);
\node at (0.25,0.25) {Nucléole};

% Dendrites
\draw[thick, popgreen] (-1.2, 0.5) .. controls (-2, 1.2) and (-2.5, 1.5) .. (-3, 1.8);
\draw[thick, popgreen] (-1.2, 0.5) .. controls (-1.8, 0.8) and (-2.2, 0.5) .. (-2.5, 0.2);
\draw[thick, popgreen] (-1.2, 0.2) .. controls (-2, 0) and (-2.5, -0.3) .. (-3, -0.5);
\draw[thick, popgreen] (-1.2, -0.2) .. controls (-1.8, -0.5) and (-2.2, -0.8) .. (-2.5, -1.2);
\node[popgreen, align=left] at (-3.5, 1.5) {Dendrites\\ (zone de réception)};

% Axon Hillock
\draw[thick, popdark] (1.2, 0) -- (1.8, 0);
\node at (1.5, -0.4) {Cône d'émergence};

% Axon with Myelin Sheath (Schwann cells)
% Myelin segments
\foreach \x in {2.0, 3.5, 5.0, 6.5, 8.0} {
    \draw[thick, fill=popblue!30] (\x, -0.4) rectangle (\x+1.2, 0.4);
    \draw[thick, popblue] (\x, -0.4) -- (\x, 0.4); % Left edge
    \draw[thick, popblue] (\x+1.2, -0.4) -- (\x+1.2, 0.4); % Right edge (Node)
    \node[popblue, rotate=90] at (\x+0.6, 0.7) {Cellule de Schwann};
}
% Nodes of Ranvier labels
\node[poppurple] at (3.2, -0.8) {Nœud de Ranvier};
\draw[poppurple, <->] (3.2, -0.6) -- (3.2, -0.45);
\node[poppurple] at (6.2, -0.8) {Nœud de Ranvier};
\draw[poppurple, <->] (6.2, -0.6) -- (6.2, -0.45);

% Axoplasm inside
\draw[popdark, dashed] (1.8, 0) -- (9.5, 0);

% Synaptic terminals
\draw[thick, fill=poppink!30] (9.5, 0) .. controls (10, 0.5) and (10.5, 0.5) .. (11, 0.2);
\draw[thick, fill=poppink!30] (9.5, 0) .. controls (10, -0.3) and (10.5, -0.3) .. (11, -0.2);
\draw[thick, fill=poppink!30] (9.5, 0) .. controls (10.2, 0) and (10.5, 0.2) .. (10.8, 0);
\node[poppink, align=left] at (11.5, 0.5) {Terminaisons\\ synaptiques\\ (boutons terminaux)};

% Legend arrows
\draw[->, thick, popdark] (-0.5, -1.2) -- (-0.5, -1.8) node[below, align=center] {Corps cellulaire\\(Soma / Péricaryon)};
\draw[->, thick, popdark] (1.2, -1.2) -- (1.2, -1.8) node[below] {Axone (Cylindraxe)};
\end{tikzpicture}
\legende{Schéma d'un neurone myélinisé type (motoneurone spinal). 1 : Corps cellulaire (soma) avec noyau et nucléole ; 2 : Dendrites (arborescentes, réception) ; 3 : Cône d'émergence (zone de déclenchement) ; 4 : Axone unique, long, gainé de myéline ; 5 : Cellules de Schwann formant la gaine de myéline ; 6 : Étranglements de Ranvier (zones nues) ; 7 : Terminaisons synaptiques (émission).}
\end{popfigure}

\begin{methode}[Schématiser un neurone (savoir-faire exigé au Bac)]
\begin{enumerate}
    \item \textbf{Corps cellulaire (soma) :} Dessiner une ellipse ou un polygone arrondi. Y placer le \textbf{noyau} (cercle central) et le \textbf{nucléole} (point central).
    \item \textbf{Dendrites :} Dessiner des prolongements courts, nombreux, très ramifiés (arborisation) partant du soma. Les légender \emph{« Dendrites : réception du message »}.
    \item \textbf{Axone (cylindraxe) :} Dessiner \textbf{un seul} prolongement long, de calibre constant, partant du cône d'émergence. Indiquer la direction de propagation (flèche).
    \item \textbf{Gaine de myéline :} Représenter des segments espacés (gainages de cellules de Schwann) le long de l'axone. Hachurer ou colorier différemment.
    \item \textbf{Nœuds de Ranvier :} Laisser des espaces nus entre les segments de myéline. Les légender.
    \item \textbf{Terminaisons synaptiques :} À l'extrémité de l'axone, dessiner des ramifications se terminant par des boutons (vésicules synaptiques suggérées). Légender \emph{« Terminaisons synaptiques : transmission »}.
    \item \textbf{Légendes :} Utiliser des traits de rappel (pas de flèches pour les légendes), écriture horizontale, vocabulaire précis.
\end{enumerate}
\end{methode}

\begin{attention}[Erreur fréquente : Nerf vs Neurone]
\begin{itemize}
    \item \textbf{Neurone} = \textbf{cellule} unique (échelle microscopique, $\approx 10$ à $100\,\mu\text{m}$ pour le soma, axone jusqu'à $1\,\text{m}$).
    \item \textbf{Fibre nerveuse} = \textbf{axone} + gaine de myéline (si myélinisé) + neurlemme (membrane de Schwann).
    \item \textbf{Nerf} = \textbf{organe} macroscopique (visible à l'œil nu), cordon blanchâtre constitué de \textbf{faisceaux de fibres nerveuses} (axones) enveloppés de tissu conjonctif (épineur, périneur, endoneur). Un nerf contient des milliers de fibres (sensitives et/ou motrices).
\end{itemize}
\emph{Ne jamais écrire : « Le nerf transmet l'influx » mais « La fibre nerveuse (axone) conduit l'influx ; le nerf est le câble qui les regroupe ».}
\end{attention}

\courssub{Classification fonctionnelle des neurones}

Selon la direction de l'information par rapport au système nerveux central (SNC), on distingue trois types fonctionnels :

\begin{center}
\begin{tabularx}{\linewidth}{|l|X|X|X|}
\hline
\rowcolor{popblueL}
\textbf{Type} & \textbf{Neurone sensitif (afferent)} & \textbf{Neurone moteur (efferent)} & \textbf{Interneurone (association)} \\
\hline
\textbf{Localisation du soma} & Ganglion rachidien (hors SNC) & Corne ventrale moelle / noyaux tronçon (dans SNC) & Substance grise (SNC uniquement) \\
\hline
\textbf{Morphologie} & Unipolaire (pseudo-unipolaire) : un tronc unique se divisant en deux (périphérique et central) & Multipolaire : 1 axone long, nombreuses dendrites & Multipolaire : axone court, dendrites ramifiées localement \\
\hline
\textbf{Fonction} & Conduit l'influx \textbf{vers} le SNC (récepteurs $\to$ SNC) & Conduit l'influx \textbf{depuis} le SNC vers l'effecteur (muscle, glande) & Intègre et relaie l'information \textbf{à l'intérieur} du SNC (99\% des neurones) \\
\hline
\textbf{Exemple} & Neurone de la peau (douleur, toucher) $\to$ moelle & Motoneurone alpha $\to$ fibre musculaire squelettique & Neurone de Renshaw (inhibition récurrente) \\
\hline
\end{tabularx}
\end{center}

\begin{exemplebox}[Le réflexe myotatique (exemple camerounais : genou du médecin)]
Au dispensaire, le médecin tape le tendon rotulien avec son marteau. L'étirement brutal du muscle quadriceps active les \textbf{neurones sensitifs} (fuscules neuromusculaires) $\to$ influx vers la moelle (racine dorsale) $\to$ synapse sur \textbf{motoneurone} (corne ventrale) $\to$ contraction du quadriceps (extension de la jambe) \textbf{ET} synapse sur \textbf{interneurone inhibiteur} $\to$ inhibition du motoneurone du muscle antagoniste (ischio-jambiers). C'est un arc réflexe \textbf{monosynaptique} (une seule synapse sensitif-moteur) pour la voie excitatrice.
\end{exemplebox}

\courssub{La gaine de myéline : isolation et rapidité}

\begin{propriete}[Formation de la gaine de myéline]
\begin{itemize}
    \item \textbf{Système Nerveux Périphérique (SNP) :} Les \textbf{cellules de Schwann} s'enroulent en spirale autour de l'axone. Chaque cellule de Schwann myélinise \textbf{un seul segment} d'un \textbf{seul axone}. La membrane plasmique de la cellule de Schwann (riche en lipides) forme la myéline ; son cytoplasme et son noyau restent à l'extérieur : la \textbf{neurlemme} (gaine de Schwann), essentielle pour la régénération après section.
    \item \textbf{Système Nerveux Central (SNC) :} Les \textbf{oligodendrocytes} étendent des prolongements qui s'enroulent autour de \textbf{plusieurs axones} (jusqu'à 30-50). Pas de neurlemme $\Rightarrow$ régénération quasi nulle dans le SNC.
\end{itemize}
\end{propriete}

\begin{popfigure}
\begin{tikzpicture}[scale=1.1]
% Axon core
\draw[thick, popdark, fill=popcream] (0,0) rectangle (10,0.6);
\node at (5, -0.5) {Axone (axoplasme)};

% Schwann cells wrapping (simplified side view)
\foreach \x in {0.5, 2.0, 3.5, 5.0, 6.5, 8.0} {
    % Myelin layers (concentric rectangles)
    \foreach \layer in {0, 0.1, 0.2, 0.3} {
        \draw[popblue!60, thick] (\x-\layer, -\layer) rectangle (\x+1.2+\layer, 0.6+\layer);
    }
    % Schwann cell nucleus (outside)
    \draw[fill=poporange!50] (\x+0.6, 1.2) circle (0.15);
    \node[font=\tiny] at (\x+0.6, 1.5) {Noyau SC};
    % Neurilemma label
    \node[font=\tiny, popblue] at (\x+0.6, -0.8) {Neurlemme};
}
% Nodes of Ranvier
\foreach \x in {1.7, 3.2, 4.7, 6.2, 7.7} {
    \draw[poppurple, very thick] (\x, -0.1) -- (\x, 0.7);
    \node[poppurple, font=\small, rotate=90] at (\x, 1.0) {Nœud de Ranvier};
}
% Ion channels at node
\draw[popred, fill=popred!20] (1.7, 0.3) rectangle (1.75, 0.5);
\draw[popred, fill=popred!20] (1.7, 0.1) rectangle (1.75, 0.3);
\node[popred, font=\tiny, rotate=90] at (1.6, 0.4) {Canaux $Na^+$/$K^+$};
\end{tikzpicture}
\legende{Organisation de la gaine de myéline dans le SNP (cellules de Schwann). La myéline (bleu) est formée par l'enroulement membranaire. Les nœuds de Ranvier (zones violettes) concentrent les canaux ioniques voltage-dépendants. La neurlemme (membrane externe de la cellule de Schwann) assure la continuité et la régénération.}
\end{popfigure}

\begin{aretenir}[Notions clés : Fibre nerveuse et Nerf]
\begin{itemize}
    \item \textbf{Fibre nerveuse myélinisée} = Axone + Gaine de myéline (segments) + Nœuds de Ranvier + Neurlemme (SNP). Diamètre total : $1$ à $20\,\mu\text{m}$. Conduction saltatoire (rapide, $\approx 100\,\text{m/s}$).
    \item \textbf{Fibre nerveuse amyélinisée} = Axone + cellules de Schwann non enroulées (gainage simple) + Neurlemme. Diamètre : $0,2$ à $1\,\mu\text{m}$. Conduction continue (lente, $\approx 1\,\text{m/s}$).
    \item \textbf{Nerf} = Faisceau de fibres nerveuses (sensitives, motrices, végétatives) + tissu conjonctif (épineur, périneur, endoneur) + vaisseaux sanguins (vasa vasorum). Structure macroscopique, blanchâtre (si myélinisé) ou grisâtre (si amyélinisé).
\end{itemize}
\end{aretenir}

\courssub{Exercice résolu : Identification sur document microscopique}

\begin{exoresolu}[Analyse d'une coupe transversale de nerf sciatique de grenouille (coloration osmique : myéline en noir)]
\textbf{Énoncé :} On observe au microscope optique ($\times 400$) une coupe transversale de nerf sciatique de grenouille. Le document montre des structures circulaires foncées (anneaux noirs) de diamètres variés, séparées par un tissu conjonctif clair. Au centre de certains anneaux, on distingue un point clair. À l'extérieur des anneaux, des noyaux aplatis sont visibles.
\begin{enumerate}
    \item Identifier les structures observées (A : anneau noir, B : point clair central, C : noyaux aplatis externes, D : tissu clair entre les faisceaux).
    \item Préciser le type de fibres nerveuses présentes. Justifier.
    \item Expliquer l'intérêt de la coloration osmique (fixateur de lipides).
\end{enumerate}
\tcblower
\textbf{Solution détaillée :}
\begin{enumerate}
    \item \textbf{A} = Gaine de myéline (bicouches lipidiques colorées en noir par l'osmium). \textbf{B} = Axone (axoplasme, pauvre en lipides, donc clair). \textbf{C} = Noyaux des cellules de Schwann (neurlemme). \textbf{D} = Tissu conjonctif (endoneur/périneur) non coloré.
    \item Ce sont des \textbf{fibres nerveuses myélinisées}. Justification : présence d'une gaine de myéline épaisse (anneau noir) interrompue par des nœuds de Ranvier (non visibles en coupe transversale unique mais inférée par la structure segmentaire) et présence de la neurlemme (noyaux externes).
    \item L'osmium (tétraoxyde d'osmium) fixe spécifiquement les \textbf{lipides} (doubles liaisons insaturées). La myéline étant riche en lipides ($\approx 70-80\%$ de la masse sèche), elle se colore intensément en noir, permettant de visualiser parfaitement l'épaisseur de la gaine et le calibre de l'axone.
\end{enumerate}
\end{exoresolu}

\begin{exercice}[Entraînement : Du neurone au nerf]
Un étudiant observe une préparation microscopique d'un nerf vague de rat. Il dessine ce qu'il voit : de gros faisceaux circulaires séparés par du tissu conjonctif. Dans chaque faisceau, des structures rondes claires (diamètre $\approx 5\,\mu\text{m}$) sont entourées d'un anneau fin et foncé. Il note la présence de noyaux aplatis en périphérie des anneaux.
\begin{enumerate}
    \item Ce document correspond-il à une coupe \textbf{longitudinale} ou \textbf{transversale} ? Justifier.
    \item Les structures rondes claires correspondent-elles aux noyaux des neurones ? Pourquoi ?
    \item Quelle est l'origine embryologique des cellules dont les noyaux sont aplatis en périphérie ?
    \item Ce nerf est-il sensitif, moteur ou mixte ? Comment le savoir avec certitude ?
\end{enumerate}
\emph{Indications : Rappeler l'organisation du tissu conjonctif nerveux (épineur, périneur, endoneur) et l'origine de la crête neurale.}
\end{exercice}

\begin{corrige}[Exercice n°1 : Entraînement : Du neurone au nerf]
\begin{enumerate}
    \item \textbf{Coupe transversale.} Justification : Les axones sont coupés perpendiculairement à leur grand axe, apparaissant comme des sections circulaires (rondes). En coupe longitudinale, on verrait des tubes longs parallèles.
    \item \textbf{Non.} Les structures claires sont les \textbf{axones} (axoplasme). Les noyaux des neurones (somata) se trouvent dans les ganglions (hors SNC) ou dans la substance grise (SNC), \textbf{jamais} dans le nerf périphérique lui-même (qui ne contient que des axones).
    \item Ces noyaux appartiennent aux \textbf{cellules de Schwann}. Elles sont issues de la \textbf{crête neurale} (ectoderme neural), tout comme les neurones sensitifs et les cellules gliales du SNP.
    \item \textbf{Impossible à dire avec certitude sur ce document seul.} Un nerf périphérique typique (comme le vague) est \textbf{mixte} (contient des fibres sensitives afférentes et des fibres motrices efférentes, dont parasympathiques). Pour le prouver, il faut soit une traçage rétrograde, soit une stimulation électrique sélective, soit l'observation de corps cellulaires moteurs (ganglions parasympathiques) et sensitifs (ganglion du vague) distincts.
\end{enumerate}
\end{corrige}

\begin{savaistu}[Record de longueur : le nerf sciatique]
Le plus long axone du corps humain (et des mammifères) appartient à un \textbf{motoneurone alpha} dont le corps cellulaire est dans la moelle épinière (segment lombaire L4-S1) et dont la terminaison innerve le muscle extenseur de l'orteil. Chez un adulte de $1,80\,\text{m}$, cet axone mesure \textbf{presque $1\,\text{m}$ de long} ! Pourtant, son diamètre n'est que de $\approx 10-20\,\mu\text{m}$. Le transport axonal (antérograde : mitochondries, vésicules ; rétrograde : facteurs neurotrophiques, virus comme le tétanos ou la rage) met des jours pour parcourir cette distance. C'est pourquoi une blessure au bas du dos (hernie discale L5-S1) peut provoquer une paralysie du pied (pied bot) sans lésion directe de la jambe.
\end{savaistu}

\coursec{Le potentiel de repos}

Après avoir rappelé la structure du tissu nerveux et l'anatomie du neurone (corps cellulaire, dendrites, axone, terminaisons synaptiques), nous nous posons maintenant une question fondamentale : \emph{que se passe-t-il au niveau électrique dans un neurone qui ne reçoit aucune stimulation ?} Est-il électriquement neutre, ou existe-t-il une activité électrique intrinsèque ? C'est l'objet de cette section : comprendre l'origine et la nature du \cle{potentiel de repos}, fondement de toute l'excitabilité neuronale.

\courssub{Problématique : une activité électrique au repos ?}

L'observation au microscope optique ou électronique d'un neurone au repos ne révèle aucun mouvement, aucun changement morphologique apparent. Pourtant, dès les années 1930-1940, les physiologistes suspectent l'existence d'une différence de potentiel électrique de part et d'autre de la membrane plasmique. L'hypothèse directrice est la suivante : \emph{la membrane du neurone au repos est sélectivement perméable à certains ions, ce qui crée une séparation de charges électriques et donc une tension électrique transmembranaire.} Pour valider cette hypothèse, il faut mesurer cette tension. Mais comment mesurer une différence de potentiel à l'intérieur d'une cellule microscopique sans la détruire ?

\courssub{Protocole de mesure du potentiel de repos (savoir-faire exigé)}

La mesure directe du potentiel transmembranaire exige une technique délicate : la \cle{microélectrode intracellulaire}. Le choix du matériel biologique est crucial : on utilise l'\cle{axone géant du calmar} (\emph{Loligo pealeii} ou \emph{Dosidicus gigas}), dont le diamètre peut atteindre \SI{1}{\milli\metre} (contre \SI{1}{\micro\metre} à \SI{20}{\micro\metre} pour un axone mammalien), permettant l'insertion d'une microélectrode sans rupture immédiate de la membrane.

\begin{methode}[Protocole de mesure du potentiel de repos]
\textbf{Objectif :} Mettre en évidence et mesurer la différence de potentiel électrique entre l'intérieur et l'extérieur d'un neurone non stimulé.

\textbf{Matériel biologique :} Axone géant de calmar (diamètre $\approx$ \SI{0,5}{--}\SI{1}{mm}) maintenu dans un bain de solution saline (eau de mer artificielle) à température contrôlée.

\textbf{Matériel instrumental :}
\begin{enumerate}
    \item \textbf{Deux électrodes :}
    \begin{itemize}
        \item Une \textbf{microélectrode de verre} (pointe effilée au micron, remplie d'une solution conductrice, ex. \ce{KCl} \SI{3}{M}), montée sur un micromanipulateur pour une insertion précise dans l'axone.
        \item Une \textbf{électrode de référence} (électrode à calomel saturé ou \ce{Ag/AgCl}) plongée dans le bain extracellulaire.
    \end{itemize}
    \item Un \textbf{voltmètre à très haute impédance d'entrée} (ou un oscilloscope à balayage lent) pour mesurer la tension sans prélever de courant significatif.
    \item Un système d'acquisition de données (ordinateur) pour l'affichage et l'enregistrement.
\end{enumerate}

\textbf{Protocole opératoire :}
\begin{enumerate}
    \item Placer l'axone dans la chambre d'enregistrement perfusée par la solution saline.
    \item Connecter l'électrode de référence au bain (définie comme le \textbf{zéro du potentiel}, $V_{\text{ext}} = 0$ mV).
    \item Approcher la microélectrode de la membrane à l'aide du micromanipulateur sous contrôle microscopique.
    \item Pénétrer brutalement mais proprement la membrane plasmique (souvent aidé par une légère surpression ou une vibration ultrasonique \emph{« buzz »}).
    \item Observer la déviation de la trace sur l'oscilloscope / voltmètre.
    \item Enregistrer la valeur stabilisée de la tension.
    \item Retirer la microélectrode pour vérifier le retour à zéro (témoin de l'intégrité de la mesure).
\end{enumerate}

\textbf{Grandeur mesurée :} Différence de potentiel $V_m = V_{\text{int}} - V_{\text{ext}}$ exprimée en \textbf{millivolts (mV)}.
\end{methode}

\begin{experience}[Mesure du potentiel de repos sur l'axone géant de calmar]
\textbf{Observation :} Avant la pénétration, l'appareil indique \SI{0}{mV} (les deux électrodes sont dans le même milieu extracellulaire). Au moment où la pointe de la microélectrode franchit la membrane plasmique, la trace chute brutalement et se stabilise à une valeur négative constante, typiquement autour de \textbf{\SI{-70}{mV}}. Lorsque la microélectrode est retirée, la trace remonte à \SI{0}{mV}.

\textbf{Interprétation :} L'intérieur de l'axone est chargé négativement par rapport à l'extérieur. Cette différence de potentiel stable, mesurée en l'absence de toute stimulation, est le \textbf{potentiel de repos} (noté $V_{\text{repos}}$ ou $E_m$).

\textbf{Conclusion :} Le neurone au repos n'est pas électriquement neutre ; il est \textbf{polarisé}. La membrane agit comme un condensateur chargé, l'intérieur étant le pôle négatif et l'extérieur le pôle positif.
\end{experience}

\begin{popfigure}
\begin{tikzpicture}[scale=0.9, every node/.style={font=\small}]
    % Bain extracellulaire
    \draw[popblue!30, fill=popblue!5, thick] (0,0) rectangle (12,5);
    \node[popblue, align=center] at (6,4.5) {Bain de solution saline (eau de mer artificielle)};
    \node[popmuted, align=center] at (6,0.3) {Milieu extracellulaire : [Na$^+$] élevé, [K$^+$] faible};
    
    % Axone géant
    \draw[popdark, fill=popcream, thick] (3,1.5) rectangle (9,3.5);
    \node[popdark, align=center] at (6,2.5) {Axone géant de calmar (diam. $\approx$ 1 mm)};
    \node[popmuted, align=center] at (6,1.8) {Milieu intracellulaire : [K$^+$] élevé, [Na$^+$] faible, Protéines$^-$};
    
    % Microélectrode intracellulaire
    \draw[poporange, very thick, -{Stealth[length=3mm]}] (6,5.5) -- (6,2.5);
    \draw[poporange, fill=poporange!20] (5.85,5.5) rectangle (6.15,6.2) node[midway, rotate=90] {Pointe $\approx$ 0.5--1 $\mu$m};
    \node[poporange, align=center, anchor=south] at (6,6.3) {Microélectrode de verre (KCl 3M)};
    \node[poporange, align=center, anchor=north] at (6,5.5) {Micromanipulateur};
    
    % Électrode de référence
    \draw[popgreen, very thick, -{Stealth[length=3mm]}] (10.5,2.5) -- (12.5,2.5);
    \draw[popgreen, fill=popgreen!20] (12.5,2.3) rectangle (13.2,2.7);
    \node[popgreen, align=center, anchor=west] at (13.3,2.5) {Électrode de référence (Calomel/\\Ag-AgCl) : \textbf{0 mV}};
    
    % Voltmetre / Oscilloscope
    \draw[poppurple, thick, fill=poppurple!10, rounded corners] (6,6.5) rectangle (11,8.5);
    \node[poppurple] at (8.5,7.8) {Voltmètre / Oscilloscope};
    \node[poppurple] at (8.5,7.3) {Haute impédance ($>10^{12}\ \Omega$)};
    \draw[poppurple, thick] (8.5,6.8) circle (0.6);
    \node[poppurple, font=\Large\bfseries] at (8.5,6.8) {\SI{-70}{mV}};
    
    % Connexions
    \draw[poporange, thick] (6,6.2) -- (6,6.5);
    \draw[popgreen, thick] (10.5,2.5) -- (10.5,6.5);
    \draw[popgreen, thick] (10.5,6.5) -- (9.5,6.5);
    \draw[poporange, thick] (6,6.5) -- (7.5,6.5);
    
    % Flèches ions (schématique)
    \draw[popblue, ->, thick] (1.5,2.5) -- (2.8,2.5) node[midway, above, font=\tiny] {Na$^+$};
    \draw[popred, <-, thick] (9.2,2.5) -- (10.5,2.5) node[midway, above, font=\tiny] {K$^+$};
\end{tikzpicture}
\legende{Dispositif expérimental classique (Hodgkin \& Huxley, 1939) pour la mesure du potentiel de repos. L'axone géant de calmar est empalé par une microélectrode de verre reliée à un voltmètre à haute impédance. L'électrode de référence dans le bain définit le zéro du potentiel. La mesure stabilisée affiche environ -70 mV.}
\end{popfigure}

\courssub{Résultat expérimental : le potentiel de repos}

\begin{definition}[Potentiel de repos]
Le \textbf{potentiel de repos} ($V_{\text{repos}}$ ou $E_m$) est la \textbf{différence de potentiel électrique} qui existe de part et d'autre de la membrane plasmique d'une cellule excitable (neurone, fibre musculaire) \textbf{non stimulée}. Il est conventionnellement exprimé en millivolts (mV) par rapport au milieu extracellulaire pris comme référence (0 mV). Sa valeur typique pour un neurone mammalien est d'environ \textbf{\SI{-70}{mV}} (entre \SI{-60}{mV} et \SI{-90}{mV} selon les types cellulaires).
\end{definition}

\begin{aretenir}[Caractéristiques du potentiel de repos]
\begin{itemize}
    \item \textbf{Valeur :} Environ \SI{-70}{mV} (intérieur négatif).
    \item \textbf{Signe :} Le signe « moins » indique que le potentiel \emph{intérieur} ($V_{\text{int}}$) est \textbf{inférieur} au potentiel \emph{extérieur} ($V_{\text{ext}}$). $V_m = V_{\text{int}} - V_{\text{ext}} = -70$ mV.
    \item \textbf{Stabilité :} Il est constant dans le temps tant que la cellule n'est pas stimulée et que son environnement ionique ne change pas.
    \item \textbf{Universalite :} Toutes les cellules animales vivantes possèdent un potentiel de repos (souvent entre -50 et -90 mV), mais seules les cellules \emph{excitable} (neurones, muscles) savent le modifier rapidement pour générer un potentiel d'action.
\end{itemize}
\end{aretenir}

\courssub{Interprétation : l'origine ionique du potentiel de repos}

Comment une simple membrane lipidique peut-elle maintenir une telle différence de potentiel ? L'expérience décisive consiste à modifier la composition ionique du milieu extracellulaire et à observer l'effet sur $V_{\text{repos}}$. Les résultats (Hodgkin, Huxley, Katz, 1949) montrent que :
\begin{enumerate}
    \item Le potentiel de repos dépend \textbf{principalement} de la concentration en ions potassium (\ce{K+}) extracellulaire.
    \item Il dépend \textbf{peu} de la concentration en ions sodium (\ce{Na+}) extracellulaire (au repos).
\end{enumerate}
Cela oriente vers l'hypothèse d'une \textbf{perméabilité sélective} de la membrane au repos : elle est \textbf{beaucoup plus perméable aux ions \ce{K+} qu'aux ions \ce{Na+}}.

\begin{propriete}[Répartition ionique transmembranaire au repos (neurone mammalien typique)]
\begin{center}
\begin{tabularx}{\linewidth}{|l|X|X|X|}\hline
\rowcolor{popblueL} \textbf{Ion / Espèce} & \textbf{Concentration intracellulaire} & \textbf{Concentration extracellulaire} & \textbf{Gradient de concentration} \\ \hline
\ce{Na+} (Sodium) & \SIrange{10}{15}{mM} & \SI{145}{mM} & \textbf{Extracellulaire $\gg$ Intracellulaire} ($\times 10$) \\ \hline
\ce{K+} (Potassium) & \SI{140}{mM} & \SI{4}{mM} & \textbf{Intracellulaire $\gg$ Extracellulaire} ($\times 35$) \\ \hline
\ce{Cl-} (Chlore) & \SIrange{4}{30}{mM} & \SI{120}{mM} & Extracellulaire $>$ Intracellulaire \\ \hline
\ce{A-} (Protéines/anions organiques) & \SI{100}{mM} (équiv.) & $\approx 0$ & \textbf{Quasi exclusivement intracellulaire} \\ \hline
\end{tabularx}
\end{center}
\textit{Source : Données physiologiques classiques (ex. neurone moteur mammalien). Les valeurs varient selon le type cellulaire.}
\end{propriete}

\begin{attention}[Le piège du signe et des gradients]
\begin{itemize}
    \item \textbf{Ne pas confondre} le \textbf{signe du potentiel} (négatif à l'intérieur) et le \textbf{sens des gradients de concentration}.
    \item \ce{K+} : Concentré \emph{à l'intérieur} $\rightarrow$ Tendance à \textbf{sortir} (gradient chimique sortant).
    \item \ce{Na+} : Concentré \emph{à l'extérieur} $\rightarrow$ Tendance à \textbf{entrer} (gradient chimique entrant).
    \item \ce{A-} (protéines) : Piégées \emph{à l'intérieur} $\rightarrow$ Contribuent à la charge négative fixe intracellulaire.
\end{itemize}
Au repos, la membrane est \textbf{perméable sélectivement au K+}. Les ions \ce{K+} sortent donc le long de leur gradient de concentration (diffusion). En sortant, ils emportent des charges positives vers l'extérieur, laissant à l'intérieur un excès de charges négatives (protéines \ce{A-} et quelques \ce{Cl-} non compensés). Cela crée le potentiel négatif intérieur.
\end{attention}

\begin{exemplebox}[Valeurs chiffrées clés à retenir pour le Bac]
\begin{itemize}
    \item Potentiel de repos typique : \textbf{\SI{-70}{mV}}.
    \item Ratio de concentrations \ce{Na+} (Ext/Int) : \textbf{$\approx 10$}.
    \item Ratio de concentrations \ce{K+} (Int/Ext) : \textbf{$\approx 35$}.
    \item Perméabilité relative au repos : $P_{\ce{K+}} / P_{\ce{Na+}} \approx \textbf{25 à 50}$ (la membrane est 25 à 50 fois plus perméable au K+ qu'au Na+).
    \item Équilibre de Nernst pour le K+ ($E_{\ce{K}}$) : \SI{-90}{mV} (calculé via $E = \frac{RT}{zF} \ln \frac{[K^+]_{\text{ext}}}{[K^+]_{\text{int}}}$). Le $V_{\text{repos}}$ (\SI{-70}{mV}) est \emph{proche} de $E_{\ce{K}}$ mais \emph{légèrement plus positif} car la faible perméabilité au \ce{Na+} (entrée de charges +) « tire » le potentiel vers le haut.
\end{itemize}
\end{exemplebox}

\courssub{Les deux piliers du maintien du potentiel de repos}

Le potentiel de repos est une situation d'\textbf{équilibre dynamique (stationnaire)} et non un équilibre thermodynamique (qui impliquerait l'absence de flux nets et la mort cellulaire). Deux mécanismes couplés le maintiennent :

\begin{popfigure}
\begin{tikzpicture}[scale=0.95, every node/.style={font=\small}]
    % Membrane lipidique
    \draw[popdark, fill=popcream, thick] (0,2) rectangle (12,3);
    \node[popdark] at (6,2.5) {Bicouche phospholipidique (Membrane plasmique)};
    \node[popmuted, font=\tiny] at (6,1.6) {Intracellulaire (Bas) \hspace{5cm} Extracellulaire (Haut)};
    
    % Ions K+ Intracellulaire (Gros)
    \foreach \x in {1, 2.5, 4, 5.5, 7, 8.5, 10} {
        \draw[popred, fill=popred!20, thick] (\x, 1.2) circle (0.25) node {\textbf{K$^+$}};
    }
    % Ions Na+ Extracellulaire (Gros)
    \foreach \x in {1.5, 3, 4.5, 6, 7.5, 9, 10.5} {
        \draw[popblue, fill=popblue!20, thick] (\x, 3.8) circle (0.25) node {\textbf{Na$^+$}};
    }
    % Ions K+ Extracellulaire (Peu)
    \foreach \x in {2, 5, 8, 11} {
        \draw[popred, fill=popred!20, thick] (\x, 3.8) circle (0.2) node {\scriptsize\textbf{K$^+$}};
    }
    % Ions Na+ Intracellulaire (Peu)
    \foreach \x in {3.5, 6.5, 9.5} {
        \draw[popblue, fill=popblue!20, thick] (\x, 1.2) circle (0.2) node {\scriptsize\textbf{Na$^+$}};
    }
    % Protéines A- (Intracellulaire)
    \foreach \x in {1.8, 4.8, 7.8, 10.8} {
        \draw[poppurple, fill=poppurple!20, thick] (\x, 1.2) ellipse (0.35 and 0.15) node {\scriptsize\textbf{A$^-$}};
    }
    
    % Canaux K+ fuite (Leak) - Ouverts au repos
    \draw[popgreen, very thick] (3, 2) -- (3, 3);
    \draw[popgreen, fill=popgreen!30] (2.7, 2.3) rectangle (3.3, 2.7);
    \node[popgreen, anchor=south, font=\tiny] at (3, 2.7) {Canal K$^+$ « fuite » (ouvert)};
    
    % Canaux Na+ fuite - Très peu / Fermés au repos
    \draw[poporange, very thick, dashed] (7, 2) -- (7, 3);
    \draw[poporange, fill=poporange!30, dashed] (6.7, 2.3) rectangle (7.3, 2.7);
    \node[poporange, anchor=north, font=\tiny] at (7, 2.3) {Canal Na$^+$ (fermé/peu perméable)};
    
    % Pompe Na+/K+ ATPase
    \draw[popgold, very thick, fill=popgold!50] (10, 2) rectangle (11, 3);
    \node[popdark, font=\tiny, align=center] at (10.5, 2.5) {Pompe \\ Na$^+$/K$^+$ \\ ATPase};
    % Flèches pompe
    \draw[popblue, ->, thick] (10.5, 1.5) -- (10.5, 2) node[midway, right, font=\tiny] {3 Na$^+$ OUT};
    \draw[popred, ->, thick] (10.5, 3) -- (10.5, 3.5) node[midway, right, font=\tiny] {2 K$^+$ IN};
    \draw[poppurple, ->, thick] (10.5, 2.5) -- (9.3, 2.5) node[midway, above, font=\tiny] {ATP $\rightarrow$ ADP + Pi};
    
    % Flux passifs K+ (Sortie gros)
    \draw[popred, ->, very thick] (3, 2.5) -- (3, 3.5) node[midway, right] {\textbf{Flux K$^+$ sortant (fuite)}};
    % Flux passifs Na+ (Entrée petit)
    \draw[popblue, ->, thick] (7, 3.5) -- (7, 2.5) node[midway, left] {Flux Na$^+$ entrant (fuite)};
    
    % Légende bas
    \node[align=left, font=\scriptsize] at (1, 0.5) {
        \textbf{Légende :} \textcolor{popred}{$\bullet$} K$^+$ (Conc. Int $\gg$ Ext) \quad
        \textcolor{popblue}{$\bullet$} Na$^+$ (Conc. Ext $\gg$ Int) \quad
        \textcolor{poppurple}{$\bullet$} A$^-$ (Protéines, Int seulement) \quad
        \textcolor{popgreen}{$\blacksquare$} Canal K$^+$ fuite (Ouvert) \quad
        \textcolor{poporange}{$\blacksquare$} Canal Na$^+$ (Fermé) \quad
        \textcolor{popgold}{$\blacksquare$} Pompe Na$^+$/K$^+$ (Active)
    };
\end{tikzpicture}
\legende{Schéma de la membrane au repos. La perméabilité sélective (canaux K$^+$ fuite ouverts, canaux Na$^+$ fermés) laisse le K$^+$ sortir le long de son gradient de concentration, créant la charge négative intérieure. La pompe Na$^+$/K$^+$ ATPase (3 Na$^+$ sortis / 2 K$^+$ rentrés / ATP hydrolysé) maintient les gradients de concentration à long terme et contribue directement (électrogénicité) à la négativité intérieure.}
\end{popfigure}

\begin{enumerate}
    \item \textbf{La perméabilité sélective passive (Canaux « fuite » ou \emph{leak channels}) :} Au repos, la membrane possède des canaux ioniques constitutivement ouverts, \textbf{sélectifs pour le K$^+$}. Le K$^+$ diffuse vers l'extérieur (son gradient chimique). Cette sortie de charges positives polarise la membrane négativement à l'intérieur. Ce flux passif tend à amener $V_m$ vers le potentiel d'équilibre de Nernst pour le K$^+$ ($E_{\ce{K}} \approx \SI{-90}{mV}$).
    \item \textbf{Le transport actif : la pompe \ce{Na+}/\ce{K+}-ATPase :} Si seul le flux passif existait, les gradients de concentration s'effondreraient (le K$^+$ sortirait jusqu'à équilibre, le Na$^+$ entrerait). La pompe \ce{Na+}/\ce{K+}-ATPase, \textbf{énergivore (hydrolyse 1 ATP par cycle)}, expulse \textbf{3 Na$^+$} vers l'extérieur et rentre \textbf{2 K$^+$} vers l'intérieur.
    \begin{itemize}
        \item Elle \textbf{maintient les gradients de concentration} (haut [K$^+$]$_{\text{int}}$, haut [Na$^+$]$_{\text{ext}}$) à long terme.
        \item Elle est \textbf{électrogène} : elle sort 3 charges positives et n'en rentre que 2 $\rightarrow$ elle contribue \textbf{directement} à l'hyperpolarisation (rend l'intérieur plus négatif de quelques mV).
    \end{itemize}
\end{enumerate}

\begin{aretenir}[Bilan mécanistique du potentiel de repos]
Le potentiel de repos (\SI{-70}{mV}) résulte de l'association de deux facteurs :
\begin{enumerate}
    \item \textbf{Facteur principal (passif) :} La \textbf{haute perméabilité au K$^+$} (canaux fuite) qui laisse le K$^+$ sortir, chargeant l'intérieur négativement. $V_m$ se rapproche de $E_{\ce{K}}$ (\SI{-90}{mV}).
    \item \textbf{Facteur correcteur (actif) :} La \textbf{pompe \ce{Na+}/\ce{K+}-ATPase} qui maintient les gradients [K$^+$] et [Na$^+$] et dont l'électrogénicité (3 Na$^+$ out / 2 K$^+$ in) tire $V_m$ vers des valeurs plus négatives.
    \item \textbf{Facteur déstabilisateur mineur :} La \textbf{faible perméabilité au Na$^+$} qui laisse entrer un peu de Na$^+$, tirant $V_m$ vers des valeurs plus positives (loin de $E_{\ce{K}}$).
\end{enumerate}
\emph{Note programme :} Au niveau Terminale, on retient le \textbf{résultat} (répartition ionique, rôle de la pompe, perméabilité K$^+$ > Na$^+$) sans démonstration biophysique complète (équation de Goldman-Hodgkin-Katz non au programme).
\end{aretenir}

\courssub{Le neurone au repos : une fibre polarisée}

\begin{definition}[Polarisation / Fibre polarisée]
On dit qu'une membrane (ou une fibre nerveuse) est \textbf{polarisée} lorsqu'il existe une différence de potentiel non nulle entre ses deux faces. Au repos, le neurone est \textbf{polarisé à \SI{-70}{mV}} (intérieur négatif). C'est cet état « chargé » (comme un condensateur) qui rend la cellule \textbf{excitable} : elle dispose d'une énergie potentielle électrique prête à être libérée sous forme de potentiel d'action si un stimulus suffit à ouvrir les canaux voltage-dépendants.
\end{definition}

\begin{savaistu}[Hodgkin, Huxley et l'axone géant : une aventure Nobel]
C'est sur l'\textbf{axone géant du calmar} (\emph{Loligo pealeii}) que les physiologistes britanniques \textbf{Alan Hodgkin} et \textbf{Andrew Huxley} ont réalisé, en 1939 puis après-guerre, les premières mesures intracellulaires fiables du potentiel de repos et d'action. Le diamètre exceptionnel de cet axone (jusqu'à \SI{1}{mm}, visible à l'œil nu) a permis l'insertion de microélectrodes en verre. Ils ont élucidé la base ionique de l'excitabilité (rôle du Na$^+$ et du K$^+$), modélisé mathématiquement le potentiel d'action (équations de Hodgkin-Huxley) et reçu le \textbf{Prix Nobel de Physiologie ou Médecine en 1963} (partagé avec John Eccles pour les synapses). Leur travail reste la pierre angulaire de la neurophysiologie moderne.
\end{savaistu}

\begin{exoresolu}[Analyse d'un protocole de mesure du potentiel de repos]
\textbf{Énoncé :} Un élève réalise une mesure de potentiel de repos sur un axone géant de calmar. Il note les étapes suivantes dans son cahier de manipulations :
\begin{enumerate}
    \item « Je place l'électrode de référence dans le bain. »
    \item « J'insère la microélectrode dans l'axone. »
    \item « Je lis \SI{-68}{mV} sur l'écran. »
    \item « Je retire la microélectrode, l'écran affiche \SI{0}{mV}. »
\end{enumerate}
\textbf{Questions :}
\begin{enumerate}
    \item Compléter le protocole en précisant le matériel biologique, la nature des électrodes, l'appareil de mesure et la grandeur mesurée avec son unité.
    \item Interpréter le signe négatif de la mesure.
    \item Expliquer pourquoi la tension revient à \SI{0}{mV} quand on retire la microélectrode.
\end{enumerate}

\tcblower
\textbf{Solution détaillée :}

\textbf{1. Protocole complet et rigoureux (attendu au Bac) :}
\begin{itemize}
    \item \textbf{Matériel biologique :} Axone géant de calmar (diamètre $\approx$ \SI{1}{mm}) maintenu dans un bain d'eau de mer artificielle (solution saline isotonique).
    \item \textbf{Électrodes :} Une microélectrode de verre (pointe $\approx$ \SI{0,5}{\micro\metre}, remplie de \ce{KCl} \SI{3}{M}) montée sur micromanipulateur pour l'enregistrement intracellulaire ; une électrode de référence (calomel saturé ou \ce{Ag/AgCl}) plongée dans le bain extracellulaire.
    \item \textbf{Appareil de mesure :} Voltmètre à très haute impédance d'entrée ($> 10^{12}\ \Omega$) ou oscilloscope, pour éviter tout courant de fuite qui fausserait la mesure.
    \item \textbf{Grandeur mesurée :} Différence de potentiel transmembranaire $V_m = V_{\text{int}} - V_{\text{ext}}$, exprimée en \textbf{millivolts (mV)}. L'électrode de référence définit le zéro ($V_{\text{ext}} = 0$ mV).
    \item \textbf{Résultat :} $V_m = \mathbf{-68\ mV}$ (valeur compatible avec la littérature : \SI{-70}{mV} $\pm$ \SI{10}{mV}).
\end{itemize}

\textbf{2. Interprétation du signe négatif :}
La valeur mesurée est $V_m = V_{\text{int}} - V_{\text{ext}} = \SI{-68}{mV}$. Comme $V_{\text{ext}}$ est la référence (0 mV), cela signifie que \textbf{le potentiel électrique à l'intérieur de l'axone est de \SI{-68}{mV} par rapport à l'extérieur}. L'intérieur est \textbf{négatif} relativement à l'extérieur. Cela traduit un excès de charges négatives (anions protéiques \ce{A-} piégés + déséquilibre de distribution des ions diffusibles K$^+$/Na$^+$/Cl$^-$) du côté intracellulaire.

\textbf{3. Retour à zéro à la retrait :}
Lorsque la microélectrode est retirée, sa pointe se retrouve à nouveau dans le milieu extracellulaire (le bain). Les deux électrodes (microélectrode et référence) baignent alors dans le \textbf{même milieu équipotentiel}. La différence de potentiel entre elles est donc nulle (\SI{0}{mV}). Ce retour à zéro valide l'absence de polarisation parasite des électrodes (potentiel de jonction liquide stable) et confirme que la valeur \SI{-68}{mV} mesurée précédemment provenait bien d'une différence de potentiel \emph{transmembranaire}.
\end{exoresolu}

\begin{exercice}[Entraînement : Le potentiel de repos]
\textbf{1.} Définir le potentiel de repos d'un neurone. Donner sa valeur typique et son unité.
\textbf{2.} Citer les deux mécanismes principaux qui permettent son maintien.
\textbf{3.} Compléter le tableau ci-dessous pour un neurone mammalien typique au repos :

\begin{center}
\begin{tabularx}{\linewidth}{|l|X|X|X|}\hline
\rowcolor{popblueL} \textbf{Ion} & \textbf{[Intracellulaire] approx.} & \textbf{[Extracellulaire] approx.} & \textbf{Sens du gradient chimique} \\ \hline
\ce{Na+} & \dots & \dots & \dots \\ \hline
\ce{K+} & \dots & \dots & \dots \\ \hline
\ce{Cl-} & \dots & \dots & \dots \\ \hline
Protéines (\ce{A-}) & \dots & \dots & \dots \\ \hline
\end{tabularx}
\end{center}

\textbf{4.} Expliquer pourquoi le potentiel de repos (\SI{-70}{mV}) est moins négatif que le potentiel d'équilibre de Nernst pour le potassium ($E_{\ce{K}} \approx \SI{-90}{mV}$).
\textbf{5.} Prédire l'effet sur le potentiel de repos d'une inhibition totale de la pompe \ce{Na+}/\ce{K+}-ATPase (ex. par l'ouabaïne) à court terme (minutes) et à long terme (heures). Justifier.
\end{exercice}

\begin{corrige}[Exercice n\textsuperscript{er} : Le potentiel de repos]
\textbf{1.} Le potentiel de repos est la différence de potentiel électrique existant de part et d'autre de la membrane plasmique d'un neurone non stimulé. Valeur typique : \textbf{\SI{-70}{mV}} (intérieur négatif).
\textbf{2.}
\begin{itemize}
    \item \textbf{Perméabilité sélective passive :} Canaux K$^+$ « fuite » ouverts $\rightarrow$ sortie de K$^+$ le long de son gradient $\rightarrow$ charge négative intérieure.
    \item \textbf{Transport actif :} Pompe \ce{Na+}/\ce{K+}-ATPase (3 Na$^+$ out / 2 K$^+$ in / ATP) $\rightarrow$ maintient les gradients de concentration + effet électrogène direct (hyperpolarisant).
\end{itemize}
\textbf{3.}
\begin{center}
\begin{tabularx}{\linewidth}{|l|X|X|X|}\hline
\rowcolor{popblueL} \textbf{Ion} & \textbf{[Intracellulaire] approx.} & \textbf{[Extracellulaire] approx.} & \textbf{Sens du gradient chimique} \\ \hline
\ce{Na+} & \SI{10}{--}\SI{15}{mM} & \SI{145}{mM} & Extracellulaire $\rightarrow$ Intracellulaire (Entrée) \\ \hline
\ce{K+} & \SI{140}{mM} & \SI{4}{mM} & Intracellulaire $\rightarrow$ Extracellulaire (Sortie) \\ \hline
\ce{Cl-} & \SI{4}{--}\SI{30}{mM} & \SI{120}{mM} & Extracellulaire $\rightarrow$ Intracellulaire (Entrée) \\ \hline
Protéines (\ce{A-}) & \SI{\approx 100}{mM} (eq.) & $\approx 0$ & Piégées Intracellulaire \\ \hline
\end{tabularx}
\end{center}
\textbf{4.} $V_{\text{repos}}$ (\SI{-70}{mV}) $> E_{\ce{K}}$ (\SI{-90}{mV}) car la membrane au repos, bien que très perméable au K$^+$, possède une \textbf{faible perméabilité résiduelle au Na$^+$}. L'entrée passive de quelques Na$^+$ (charges positives) dépolarise légèrement la membrane par rapport à l'équilibre parfait K$^+$. La pompe \ce{Na+}/\ce{K+} (électrogène) s'oppose à cette dépolarisation mais ne l'annule pas totalement.
\textbf{5.}
\begin{itemize}
    \item \textbf{Court terme (min) :} $V_{\text{repos}}$ \textbf{diminue légèrement (dépolarisation de quelques mV)}. Perte de la contribution électrogène directe de la pompe (3 charges + sorties vs 2 rentrées). Les gradients [K$^+$] et [Na$^+$] n'ont pas encore le temps de s'effondrer significativement.
    \item \textbf{Long terme (heures) :} $V_{\text{repos}}$ \textbf{s'effondre vers \SI{0}{mV} (dépolarisation massive)}. Sans pompe, les fuites passives (K$^+$ sort, Na$^+$ entre) dissipent les gradients de concentration. [K$^+$]$_{\text{int}}$ chute, [Na$^+$]$_{\text{int}}$ monte $\rightarrow$ $E_{\ce{K}}$ devient moins négatif $\rightarrow$ $V_{\text{repos}}$ suit. La cellule meurt (perte de l'excitabilité, œdème cellulaire par entrée d'eau osmotique).
\end{itemize}
\end{corrige}

\coursec{Le potentiel d'action}

Dans la section précédente, nous avons établi que tout neurone au repos maintient une différence de potentiel transmembranaire stable, le \cle{potentiel de repos}, d'environ $-70$ mV (face interne négative par rapport à la face externe). Cet état électrique résulte de la répartition inégale des ions $\ce{Na+}$ et $\ce{K+}$ et du travail incessant de la pompe $\ce{Na+}/\ce{K+}$-ATPase. Mais un neurone ne sert pas à rester au repos : sa fonction est de \emph{recevoir} et de \emph{transmettre} des informations. Que devient alors le potentiel membranaire lorsque le neurone est stimulé ? C'est toute la question de cette section, et sa réponse constitue l'un des savoirs les plus classiques — et les plus rentables — du programme de Terminale.

\courssub{Problème et dispositif expérimental}

\textbf{Problème.} Lorsqu'un neurone reçoit une stimulation (pression sur la peau, molécule odorante, message d'un neurone voisin), son potentiel membranaire reste-t-il inchangé, ou subit-il une modification ? Et si modification il y a, quelle est son ampleur, sa durée, ses causes ?

\begin{experience}[Stimulation électrique de l'axone géant de calmar]
\textbf{Matériel.} On utilise l'\cle{axone géant} du calmar (\emph{Loligo}), une fibre nerveuse dont le diamètre peut atteindre \SI{1}{mm}, ce qui permet d'y introduire facilement des microélectrodes. Le montage comprend :
\begin{itemize}
\item une \textbf{microélectrode de mesure} enfoncée à l'intérieur de l'axone, reliée à un oscilloscope cathodique (la référence étant une électrode placée dans le milieu extracellulaire) ;
\item deux \textbf{électrodes de stimulation} délivrant des chocs électriques brefs (durée inférieure à \SI{1}{ms}) d'intensité réglable.
\end{itemize}
\textbf{Protocole.} On enregistre d'abord le potentiel de repos ($-70$ mV environ). On applique ensuite des stimulations d'intensité croissante et on observe l'écran de l'oscilloscope.

\textbf{Observations.}
\begin{itemize}
\item Pour des stimulations \textbf{faibles}, l'écran n'affiche aucune variation notable : le potentiel reste voisin de $-70$ mV.
\item À partir d'une certaine intensité de stimulation, appelée \cle{seuil}, l'oscilloscope enregistre une variation brève et spectaculaire : le potentiel s'élève brutalement, \textbf{dépasse zéro} et atteint environ $+30$ mV, puis redescend, devient même transitoirement plus négatif que le potentiel de repos (environ $-80$ mV), avant de revenir à $-70$ mV. L'ensemble dure à peine 1 à 2 ms.
\item Si l'on augmente encore l'intensité de la stimulation au-delà du seuil, la forme et l'amplitude de la variation enregistrée \textbf{ne changent pas}.
\end{itemize}
\textbf{Interprétation.} La stimulation efficace déclenche une réponse électrique stéréotypée : une inversion transitoire de la polarité membranaire. Le fait que cette réponse soit identique quelle que soit l'intensité du stimulus supraliminaire suggère qu'elle obéit à une loi de type « tout ou rien ».
\end{experience}

\courssub{Définition et courbe du potentiel d'action}

\begin{definition}[Potentiel d'action]
Le \cle{potentiel d'action} (PA) est une \textbf{variation brève et transitoire du potentiel membranaire}, caractérisée par une \textbf{inversion de polarité} (la face interne devient positive, environ $+30$ mV), déclenchée par une \textbf{stimulation efficace} (d'intensité au moins égale au seuil). Il est suivi d'un retour spontané au potentiel de repos.
\end{definition}

\begin{exemplebox}[Ordres de grandeur à connaître pour le Bac]
\begin{itemize}
\item Amplitude totale : environ \SI{100}{mV} (de $-70$ mV à $+30$ mV) ;
\item Durée totale : environ 1 à 2 ms ;
\item Potentiel seuil : environ $-55$ mV (valeur de référence pour les neurones de mammifères) ;
\item Hyperpolarisation : jusqu'à environ $-80$ mV.
\end{itemize}
Ces valeurs sont variables selon les types de neurones et la température, mais les ordres de grandeur ci-dessus doivent être maîtrisés pour l'épreuve.
\end{exemplebox}

\begin{popfigure}
\begin{center}
\begin{tikzpicture}
\begin{axis}[
  width=0.95\linewidth, height=7.2cm,
  xlabel={Temps (ms)}, ylabel={Potentiel membranaire (mV)},
  xmin=-0.3, xmax=6.5, ymin=-95, ymax=55,
  xtick={0,1,2,3,4,5,6},
  ytick={-80,-70,-55,0,30},
  grid=both, grid style={popline!40},
  axis line style={popdark},
  label style={popdark},
  clip=true
]
% Potentiel de repos
\addplot[popcurve, domain=-0.3:1, samples=20]{-70};
% Dépolarisation
\addplot[popcurve, domain=1:1.9, samples=60]{-70 + 100/(1+exp(-9*(x-1.45)))};
% Repolarisation
\addplot[popcurve, domain=1.9:3.6, samples=80]{30 - 110/(1+exp(-6*(x-2.6)))};
% Hyperpolarisation puis retour au repos
\addplot[popcurve, domain=3.6:6.5, samples=60]{-70 - 10*exp(-(x-3.6)/0.9)};
% Lignes de référence (étiquettes à droite pour éviter le rognage)
\draw[dashed, popmuted] (axis cs:-0.3,-70) -- (axis cs:6.5,-70) node[pos=0.98, anchor=west, popdark, font=\small]{repos $-70$ mV};
\draw[dashed, poporange] (axis cs:-0.3,-55) -- (axis cs:6.5,-55) node[pos=0.98, anchor=west, popdark, font=\small]{seuil $-55$ mV};
\draw[dashed, popmuted] (axis cs:-0.3,0) -- (axis cs:6.5,0);
\draw[dashed, popgreen] (axis cs:-0.3,30) -- (axis cs:6.5,30) node[pos=0.98, anchor=west, popdark, font=\small]{$+30$ mV};
% Annotations des phases
\node[popblue, font=\small\bfseries] at (axis cs:1.35,15) {dépolarisation};
\node[popgreen, font=\small\bfseries] at (axis cs:2.35,42) {pic};
\node[poppurple, font=\small\bfseries] at (axis cs:3.1,5) {repolarisation};
\node[poporange, font=\small\bfseries] at (axis cs:4.6,-90) {hyperpolarisation};
\node[popdark, font=\small] at (axis cs:6,-62) {retour au repos};
% Stimulus
\draw[-{Stealth}, poppink, thick] (axis cs:1,-95) -- (axis cs:1,-78);
\node[poppink, font=\small] at (axis cs:1.55,-92) {stimulus};
\end{axis}
\end{tikzpicture}
\end{center}
\legende{Courbe du potentiel d'action enregistrée à l'oscilloscope sur l'axone géant de calmar. En abscisse : le temps (ms) ; en ordonnée : le potentiel membranaire (mV). On distingue le potentiel de repos ($-70$ mV), le seuil (environ $-55$ mV), la dépolarisation rapide jusqu'au pic ($+30$ mV), la repolarisation, l'hyperpolarisation transitoire, puis le retour au repos.}
\end{popfigure}

\begin{methode}[Analyser et interpréter une courbe de potentiel d'action]
Face à une courbe de PA au Baccalauréat, procède toujours dans le même ordre :
\begin{enumerate}
\item \textbf{Identifier la valeur du potentiel de repos} (plateau initial, environ $-70$ mV) : la membrane est polarisée, face interne négative.
\item \textbf{Repérer le seuil} : valeur du potentiel (environ $-55$ mV) à partir de laquelle la variation s'emballe.
\item \textbf{Repérer la dépolarisation} : montée rapide du potentiel, franchissement de zéro, pic à $+30$ mV environ.
\item \textbf{Repérer la repolarisation} : descente rapide du potentiel vers les valeurs négatives.
\item \textbf{Repérer l'hyperpolarisation} éventuelle : le potentiel devient transitoirement plus négatif que le repos (environ $-80$ mV).
\item \textbf{Associer chaque phase aux mouvements ioniques} correspondants (voir ci-dessous) : c'est l'étape d'interprétation, la plus valorisée dans la notation.
\end{enumerate}
\end{methode}

\courssub{Interprétation ionique des phases du potentiel d'action}

La membrane du neurone possède, outre les canaux « passifs » responsables du potentiel de repos, des \cle{canaux ioniques voltage-dépendants} : des canaux $\ce{Na+}$ et des canaux $\ce{K+}$ qui s'ouvrent ou se ferment selon la valeur du potentiel membranaire. Ce sont eux qui expliquent la courbe.

\textbf{1. La dépolarisation : entrée massive de $\ce{Na+}$.} Lorsque la stimulation amène le potentiel au seuil (environ $-55$ mV), les canaux $\ce{Na+}$ voltage-dépendants s'ouvrent brutalement. Comme le $\ce{Na+}$ est beaucoup plus concentré à l'extérieur et que la face interne est négative, les ions $\ce{Na+}$ (chargés positivement) \emph{se précipitent vers l'intérieur}, entraînés à la fois par le gradient de concentration et par le gradient électrique. L'entrée de charges positives rend l'intérieur de la membrane de moins en moins négatif, ce qui ouvre encore plus de canaux $\ce{Na+}$ : c'est un phénomène auto-amplificateur. Le potentiel franchit zéro et culmine vers $+30$ mV : c'est l'\textbf{inversion de polarité}.

\textbf{2. La repolarisation : fermeture des canaux $\ce{Na+}$ et sortie de $\ce{K+}$.} Au sommet du pic, deux événements se produisent : les canaux $\ce{Na+}$ s'inactivent (ils se ferment et deviennent temporairement incapables de se rouvrir), et les canaux $\ce{K+}$ voltage-dépendants, plus lents, s'ouvrent largement. Les ions $\ce{K+}$, très concentrés à l'intérieur, \emph{sortent} alors massivement. Cette perte de charges positives rend la face interne à nouveau négative : le potentiel redescend.

\textbf{3. L'hyperpolarisation.} Les canaux $\ce{K+}$ se ferment avec un certain retard : le $\ce{K+}$ continue de sortir quelques instants de trop. Le potentiel devient alors transitoirement plus négatif que le potentiel de repos (environ $-80$ mV).

\textbf{4. Le retour au repos.} Les canaux $\ce{K+}$ finissent par se fermer ; la pompe $\ce{Na+}/\ce{K+}$-ATPase et les courants de fuite rétablissent le potentiel de repos à $-70$ mV. Notons que les quantités d'ions déplacées lors d'un PA sont infimes : les gradients de concentration ne sont pas modifiés de façon mesurable, et la pompe n'a qu'un rôle d'entretien à long terme.

\begin{popfigure}
\begin{center}
\begin{tikzpicture}[font=\small]
% Trois états de membrane
\foreach \x/\titre/\etat in {0/{Repos ($-70$ mV)}/repos, 5.2/{Dépolarisation ($+30$ mV)}/depol, 10.4/{Repolarisation}/repol}{
  % milieu extérieur / intérieur
  \fill[popblueL] (\x-1.7,0.6) rectangle (\x+1.7,2.2);
  \fill[popcream] (\x-1.7,-2.2) rectangle (\x+1.7,-0.6);
  \draw[popdark, very thick] (\x-1.7,0.6) -- (\x+1.7,0.6);
  \draw[popdark, very thick] (\x-1.7,-0.6) -- (\x+1.7,-0.6);
  \node[popdark] at (\x,1.9) {extérieur};
  \node[popdark] at (\x,-1.9) {intérieur};
  \node[popdark, font=\small\bfseries] at (\x,2.7) {\titre};
  % canal Na+ (gauche)
  \draw[popgreen, thick, fill=popgreenL] (\x-1.1,0.6) rectangle (\x-0.5,-0.6);
  \node[popdark, font=\scriptsize] at (\x-0.8,-1.1) {canal $\ce{Na+}$};
  % canal K+ (droite)
  \draw[poppurple, thick, fill=poppinkL] (\x+0.5,0.6) rectangle (\x+1.1,-0.6);
  \node[popdark, font=\scriptsize] at (\x+0.8,-1.1) {canal $\ce{K+}$};
}
% Repos : canaux fermés (croix)
\draw[popink, very thick] (-1.1,0.15) -- (-0.5,-0.15);
\draw[popink, very thick] (-1.1,-0.15) -- (-0.5,0.15);
\draw[popink, very thick] (0.5,0.15) -- (1.1,-0.15);
\draw[popink, very thick] (0.5,-0.15) -- (1.1,0.15);
\node[popink, font=\scriptsize] at (0,-2.9) {canaux voltage-dépendants fermés};
% Dépolarisation : Na+ entre
\draw[-{Stealth}, popgreen, very thick] (4.4,1.5) -- (4.4,-1.4);
\node[popgreen!60!popdark, font=\scriptsize\bfseries] at (3.4,0.4) {entrée de $\ce{Na+}$};
\draw[popink, very thick] (10.9,0.15) -- (11.5,-0.15);
\draw[popink, very thick] (10.9,-0.15) -- (11.5,0.15);
% Repolarisation : K+ sort
\draw[-{Stealth}, poppurple, very thick] (11.2,-1.5) -- (11.2,1.4);
\node[poppurple, font=\scriptsize\bfseries] at (12.3,0.4) {sortie de $\ce{K+}$};
\draw[popink, very thick] (9.3,0.15) -- (9.9,-0.15);
\draw[popink, very thick] (9.3,-0.15) -- (9.9,0.15);
\node[popink, font=\scriptsize] at (9.6,-2.9) {canal $\ce{Na+}$ inactivé};
% Flèches de transition
\draw[-{Stealth}, popdark, thick] (2.1,0) -- (3.1,0);
\draw[-{Stealth}, popdark, thick] (7.3,0) -- (8.3,0);
\end{tikzpicture}
\end{center}
\legende{Mouvements ioniques au cours du potentiel d'action. Au repos, les canaux voltage-dépendants sont fermés. Lors de la dépolarisation, l'ouverture des canaux $\ce{Na+}$ provoque une entrée massive d'ions $\ce{Na+}$. Lors de la repolarisation, les canaux $\ce{Na+}$ s'inactivent et l'ouverture des canaux $\ce{K+}$ permet une sortie massive d'ions $\ce{K+}$.}
\end{popfigure}

\courssub{La loi du tout ou rien et la période réfractaire}

Reprenons les résultats de l'expérience du calmar : en dessous du seuil, aucune réponse ; au-dessus, une réponse toujours identique. Cette observation fondamentale se généralise.

\begin{propriete}[Loi du tout ou rien (admise)]
Le potentiel d'action obéit à la \cle{loi du tout ou rien} :
\begin{itemize}
\item si l'intensité de la stimulation est \textbf{inférieure au seuil}, il n'y a \textbf{pas de potentiel d'action} ;
\item si l'intensité est \textbf{égale ou supérieure au seuil}, un potentiel d'action se déclenche, avec \textbf{toujours la même amplitude} (environ \SI{100}{mV}) et la même durée, quelle que soit l'intensité du stimulus.
\end{itemize}
Cette propriété est admise au programme ; sa justification repose sur le caractère auto-amplificateur de l'ouverture des canaux $\ce{Na+}$ voltage-dépendants.
\end{propriete}

Une conséquence importante découle de l'inactivation des canaux $\ce{Na+}$ : juste après un PA, la membrane est \textbf{inexcitable} pendant environ 1 ms. C'est la \cle{période réfractaire} (savoir admis) : pendant ce laps de temps, aucune stimulation, même très intense, ne peut déclencher un nouveau PA. Il en résulte que les potentiels d'action ne peuvent pas se succéder à une fréquence infinie : le message nerveux est constitué de PA \textbf{discrets et discontinus}, séparés les uns des autres, comme des salves d'impulsions.

\begin{attention}[Le piège classique du Baccalauréat]
Erreur fréquente : écrire qu'« un stimulus plus fort produit un potentiel d'action plus grand ». C'est \textbf{faux}. En vertu de la loi du tout ou rien, l'amplitude du PA est constante. Un stimulus plus intense se traduit par une \textbf{fréquence plus élevée} de potentiels d'action (plus de PA par seconde), jamais par des PA plus amples. L'information sur l'intensité d'une sensation (une lumière plus vive, une douleur plus forte) est donc codée en \cle{fréquence}, non en amplitude.
\end{attention}

\begin{aretenir}[L'essentiel sur le codage du message nerveux]
Le message nerveux est un \textbf{signal électrique} constitué d'une succession de potentiels d'action tous identiques. Il est \textbf{codé en fréquence} : plus la stimulation est intense, plus la fréquence des PA est élevée (jusqu'à quelques centaines de PA par seconde, limite imposée par la période réfractaire). Ni l'amplitude ni la durée des PA ne transportent d'information.
\end{aretenir}

\begin{exoresolu}[Lecture d'un enregistrement oscilloscopique]
\textbf{Énoncé.} Sur l'axone géant d'un calmar maintenu à \SI{20}{\celsius}, on enregistre le potentiel membranaire. Au repos, l'oscilloscope indique $-70$ mV. On applique trois stimulations successives : $S_1$ d'intensité 2 unités arbitraires (u.a.), $S_2$ de 5 u.a. et $S_3$ de 10 u.a. Résultats : $S_1$ ne provoque aucune variation ; $S_2$ et $S_3$ provoquent chacune un PA culminant à $+30$ mV et durant environ 2 ms.
\begin{enumerate}
\item Calculer l'amplitude du potentiel d'action enregistré.
\item Que représente l'intensité seuil dans cette expérience ? Encadrer sa valeur.
\item Pourquoi $S_3$, deux fois plus intense que $S_2$, ne produit-elle pas un PA plus ample ?
\end{enumerate}
\tcblower
\textbf{Solution détaillée.}
\begin{enumerate}
\item L'amplitude est la différence entre le pic et le potentiel de repos :
\[ A = (+30) - (-70) = 100 \text{ mV}. \]
\item L'intensité seuil est l'intensité minimale de stimulation capable de déclencher un PA. Ici, $S_1$ (2 u.a.) est inefficace et $S_2$ (5 u.a.) est efficace : le seuil est donc compris entre 2 et 5 u.a. (bornes exclues pour la valeur exacte).
\item Parce que le PA obéit à la loi du tout ou rien : dès que le seuil est atteint, la réponse est maximale et stéréotypée (amplitude constante d'environ \SI{100}{mV}). L'augmentation de l'intensité du stimulus ne modifie ni l'amplitude ni la durée du PA ; in vivo, elle se traduirait par une augmentation de la fréquence des PA.
\end{enumerate}
\end{exoresolu}

\begin{savaistu}[Des calmars aux hôpitaux camerounais]
C'est grâce à l'axone géant du calmar que Alan Hodgkin et Andrew Huxley ont élucidé, dans les années 1950, le mécanisme ionique du potentiel d'action — travaux couronnés par le prix Nobel de physiologie ou médecine en 1963. Aujourd'hui, la mesure des signaux électriques du système nerveux est un outil diagnostique quotidien : l'électroencéphalogramme (EEG), pratiqué dans les hôpitaux de référence de Yaoundé et de Douala, enregistre l'activité électrique globale de millions de neurones et aide au diagnostic de l'épilepsie, une affection neurologique fréquente au Cameroun.
\end{savaistu}

En résumé, la stimulation efficace d'un neurone déclenche un potentiel d'action : inversion brève de la polarité membranaire due à l'entrée de $\ce{Na+}$ puis à la sortie de $\ce{K+}$, obéissant à la loi du tout ou rien et codé en fréquence. Reste à comprendre comment ce signal, né en un point du neurone, se propage le long de la fibre nerveuse : ce sera l'objet de la section suivante.

\coursec{La conduction du message nerveux et sa vitesse}

Dans la section précédente, nous avons vu comment un neurone, stimulé au-dessus de son seuil, génère en un point de sa membrane un \cle{potentiel d'action} : une inversion brève et stéréotypée du potentiel de membrane, de $-\SI{70}{mV}$ à environ $+\SI{30}{mV}$, suivie d'une repolarisation. Mais un potentiel d'action qui naît au niveau du corps cellulaire ou du cône d'émergence ne servirait à rien s'il restait sur place. Or le mot qu'il porte peut devoir parcourir des distances considérables : l'axone d'un motoneurone qui commande les muscles du pied mesure près d'un mètre chez l'adulte. Se pose alors le problème central de cette section : \textbf{comment le potentiel d'action se déplace-t-il le long de la fibre nerveuse, et à quelle vitesse ?}

\courssub{Le problème de la propagation : une observation pour démarrer}

\begin{experience}[Une observation qui pose problème]
Pince-t-on la pointe du pied d'une grenouille décérébrée : la patte se rétracte presque instantanément. Le message douloureux a donc voyagé le long du nerf sciatique, sur plusieurs centimètres, en une fraction de seconde. Deux hypothèses s'offrent à nous :
\begin{itemize}
\item \textbf{Hypothèse 1 :} le message nerveux est un courant électrique qui circule dans l'axone comme dans un fil de cuivre.
\item \textbf{Hypothèse 2 :} le message nerveux est un événement membranaire qui se \emph{régénère} de proche en proche le long de la fibre.
\end{itemize}
Pour trancher, enregistrons le potentiel d'action en deux points éloignés d'un même axone. Si l'hypothèse 1 est correcte, l'amplitude du signal doit \emph{diminuer} avec la distance (comme la tension diminue le long d'un fil résistif). Si l'hypothèse 2 est correcte, l'amplitude doit rester \emph{identique} quel que soit le point d'enregistrement.
\end{experience}

\textbf{Résultat expérimental :} le potentiel d'action enregistré à \SI{10}{cm} du point de stimulation a exactement la même amplitude (environ $\SI{100}{mV}$ de variation totale) que celui enregistré à \SI{1}{cm}. Seul change le \emph{retard} avec lequel il apparaît. L'hypothèse 1 est donc rejetée : le message nerveux ne s'affaiblit pas avec la distance, il se \textbf{régénère}.

\begin{definition}[Influx nerveux]
On appelle \cle{influx nerveux} le \textbf{potentiel d'action qui se propage} le long d'une fibre nerveuse, de proche en proche, sans perte d'amplitude. La propagation de l'influx nerveux constitue la \cle{conduction} du message nerveux.
\end{definition}

\textbf{Le mécanisme de la propagation : les courants locaux}\par

\textbf{Comment un potentiel d'action en naît un autre}\par

Considérons une région de l'axone en train de produire un potentiel d'action : sa membrane est \emph{dépolarisée}, c'est-à-dire que son intérieur est devenu positif ($+\SI{30}{mV}$) par rapport à l'extérieur, suite à l'ouverture des canaux \ce{Na+} voltage-dépendants. Juste à côté, la région voisine est encore au repos : intérieur négatif ($-\SI{70}{mV}$), extérieur positif.

Cette différence de potentiel entre deux régions adjacentes de la membrane provoque, de part et d'autre de celle-ci, de petits déplacements de charges électriques appelés \cle{courants locaux} :
\begin{itemize}
\item à l'\textbf{intérieur} de l'axone, les charges positives s'écoulent de la zone active (positive) vers la zone au repos (négative) ;
\item à l'\textbf{extérieur}, le courant circule en sens inverse, de la zone au repos vers la zone active.
\end{itemize}

Le courant local qui atteint la région voisine y fait entrer des charges positives : le potentiel de membrane de cette région s'élève donc de $-\SI{70}{mV}$ vers des valeurs moins négatives. C'est une \textbf{dépolarisation}. Lorsque celle-ci atteint le \cle{seuil} (environ $-\SI{55}{mV}$), les canaux \ce{Na+} voltage-dépendants de cette région s'ouvrent à leur tour : un \textbf{nouveau potentiel d'action naît}. Celui-ci dépolarise la région suivante, et ainsi de suite.

\begin{propriete}[Propagation de proche en proche]
Le potentiel d'action créé en un point de la fibre nerveuse engendre des \cle{courants locaux} qui dépolarisent la région voisine jusqu'au seuil ; celle-ci génère à son tour un potentiel d'action de même amplitude. Le message nerveux se propage ainsi \textbf{de proche en proche, sans affaiblissement}, chaque segment de membrane régénérant intégralement le signal.
\end{propriete}

\begin{attention}[Le message nerveux n'est pas un courant dans un fil]
Erreur classique de rédaction au Baccalauréat : écrire que « le courant nerveux circule dans l'axone comme dans un fil électrique ». C'est faux. Dans un fil, le signal s'atténue avec la distance ; dans l'axone, le potentiel d'action est \textbf{régénéré} à chaque segment de membrane grâce à l'ouverture de nouveaux canaux \ce{Na+}. Les courants locaux ne servent qu'à \emph{déclencher} le PA voisin, pas à transporter le message. Image utile : une rangée de dominos — chaque domino qui tombe fait tomber le suivant ; l'énergie vient de chaque domino, pas du premier.
\end{attention}

\begin{popfigure}
\begin{center}
\begin{tikzpicture}[scale=1.0, every node/.style={font=\small}]
% Axone : cylindre stylisé
\draw[thick, popblue] (0,0.8) -- (12,0.8);
\draw[thick, popblue] (0,-0.8) -- (12,-0.8);
\draw[thick, popblue] (0,0.8) arc (90:270:0.15 and 0.8);
\draw[thick, popblue] (12,0.8) arc (90:-90:0.15 and 0.8);
% Zone active (PA)
\fill[poporange!35] (3.5,-0.8) rectangle (5.5,0.8);
\node[popdark] at (4.5,0) {\textbf{PA}};
% Charges extérieures
\node[popink] at (1.2,1.15) {$+$};
\node[popink] at (2.2,1.15) {$+$};
\node[popink] at (3.0,1.15) {$+$};
\node[popink] at (4.0,1.15) {$-$};
\node[popink] at (4.5,1.15) {$-$};
\node[popink] at (5.0,1.15) {$-$};
\node[popink] at (6.0,1.15) {$+$};
\node[popink] at (7.0,1.15) {$+$};
\node[popink] at (8.0,1.15) {$+$};
\node[popink] at (9.0,1.15) {$+$};
\node[popink] at (10.0,1.15) {$+$};
\node[popink] at (11.0,1.15) {$+$};
% Charges intérieures
\node[popink] at (1.2,-1.15) {$-$};
\node[popink] at (2.2,-1.15) {$-$};
\node[popink] at (3.0,-1.15) {$-$};
\node[popink] at (4.0,-1.15) {$+$};
\node[popink] at (4.5,-1.15) {$+$};
\node[popink] at (5.0,-1.15) {$+$};
\node[popink] at (6.0,-1.15) {$-$};
\node[popink] at (7.0,-1.15) {$-$};
\node[popink] at (8.0,-1.15) {$-$};
\node[popink] at (9.0,-1.15) {$-$};
\node[popink] at (10.0,-1.15) {$-$};
\node[popink] at (11.0,-1.15) {$-$};
% Courants locaux (flèches)
\draw[-{Stealth[length=3mm]}, very thick, popgreen] (5.3,0.35) .. controls (6.3,0.55) .. (7.3,0.35);
\draw[-{Stealth[length=3mm]}, very thick, popgreen] (7.3,-0.35) .. controls (6.3,-0.55) .. (5.3,-0.35);
\draw[-{Stealth[length=3mm]}, very thick, popgreen] (3.7,0.35) .. controls (2.7,0.55) .. (1.7,0.35);
\draw[-{Stealth[length=3mm]}, very thick, popgreen] (1.7,-0.35) .. controls (2.7,-0.55) .. (3.7,-0.35);
% Sens de propagation
\draw[-{Stealth[length=4mm]}, ultra thick, poppurple] (8.0,-1.7) -- (11.0,-1.7);
\node[poppurple, font=\small\bfseries] at (9.5,-2.1) {Sens de propagation de l'influx};
% Annotations
\node[popdark, font=\small] at (4.5,1.7) {Zone active (dépolarisée)};
\node[popdark, font=\small] at (9.0,1.7) {Zone au repos};
\node[popgreen, font=\small] at (6.3,1.05) {courants locaux};
\end{tikzpicture}
\end{center}
\legende{Propagation du potentiel d'action le long d'un axone. La zone active (orange, membrane dépolarisée : intérieur $+$, extérieur $-$) crée des courants locaux (flèches vertes) qui dépolarisent la région voisine encore au repos ; celle-ci génère à son tour un PA. Le signal se régénère ainsi de proche en proche sans s'affaiblir.}
\end{popfigure}

\textbf{Pourquoi la propagation est-elle unidirectionnelle \emph{in vivo} ?}\par

Les courants locaux partent en réalité des \emph{deux côtés} de la zone active. Pourtant, dans l'organisme, l'influx nerveux se propage toujours dans un seul sens : du corps cellulaire vers l'extrémité de l'axone (conduction \cle{orthodromique}). Pourquoi ne revient-il pas en arrière ?

La réponse tient à la \cle{période réfractaire}. Immédiatement après le passage d'un potentiel d'action, la membrane traversée est temporairement \textbf{inexcitable} :
\begin{itemize}
\item pendant la \textbf{période réfractaire absolue} (environ $\SI{1}{ms}$), les canaux \ce{Na+} sont inactivés : aucun nouveau PA ne peut naître, quelle que soit l'intensité de la stimulation ;
\item pendant la \textbf{période réfractaire relative} (quelques millisecondes), seule une stimulation supraliminaire peut déclencher un PA.
\end{itemize}
Ainsi, la région que l'influx vient de quitter est réfractaire : le courant local qui revient vers elle ne peut y déclencher aucun PA. Seule la région située \emph{en avant}, encore au repos et donc excitable, répond. La propagation est donc \textbf{unidirectionnelle}.

\begin{attention}[Distinction à ne pas confondre]
Sur une fibre nerveuse \emph{isolée} stimulée en son milieu (expérience de laboratoire), l'influx se propage dans \emph{les deux sens} à partir du point stimulé : la conduction est alors bidirectionnelle. La propagation unidirectionnelle n'existe que \emph{in vivo}, parce que la stimulation naturelle naît toujours à une extrémité du neurone et que la période réfractaire interdit tout retour en arrière.
\end{attention}

\courssub{La mesure de la vitesse de conduction}

\textbf{Le dispositif expérimental}\par

\begin{experience}[Mesure de la vitesse de conduction sur le nerf sciatique de grenouille]
\textbf{Matériel :} un nerf sciatique de grenouille maintenu en vie dans une cuve à nerfs (liquide physiologique), un stimulateur électrique, deux paires d'électrodes d'enregistrement $E_1$ et $E_2$ reliées à un oscilloscope à deux voies, une règle millimétrée.

\textbf{Protocole :}
\begin{enumerate}
\item On place le nerf sur les électrodes ; on mesure la distance $d$ séparant $E_1$ de $E_2$ (par exemple $d = \SI{30}{mm}$).
\item On applique une stimulation efficace à l'extrémité proximale du nerf.
\item On enregistre simultanément les potentiels d'action captés par $E_1$ puis par $E_2$.
\item On mesure sur l'écran le décalage temporel $\Delta t$ entre les deux PA (par exemple $\Delta t = \SI{1}{ms}$).
\end{enumerate}

\textbf{Observations :} les deux PA enregistrés ont la \emph{même amplitude} (confirmation de la régénération du signal), mais celui de $E_2$ apparaît avec un \emph{retard} $\Delta t$ par rapport à celui de $E_1$, retard d'autant plus grand que la distance $d$ est grande.

\textbf{Interprétation :} le temps $\Delta t$ correspond à la durée mise par l'influx pour parcourir la distance $d$. La vitesse de conduction est donc :
\[ v = \frac{d}{\Delta t} = \frac{\SI{30}{mm}}{\SI{1}{ms}} = \SI{30}{m.s^{-1}}. \]
\end{experience}

\begin{methode}[Calculer une vitesse de conduction]
Pour calculer la vitesse de conduction d'un influx nerveux :
\begin{enumerate}
\item Repérer sur le document la \textbf{distance} $d$ entre les deux électrodes d'enregistrement (ou entre le point de stimulation et l'électrode, si le temps de latence est donné).
\item Mesurer sur l'oscillogramme le \textbf{décalage temporel} $\Delta t$ entre les deux potentiels d'action, en tenant compte de l'échelle de temps (balayage).
\item Appliquer la relation $v = \dfrac{d}{\Delta t}$, en convertissant les unités : distance en mètres, temps en secondes, vitesse en $\si{m.s^{-1}}$.
\item Vérifier la cohérence du résultat : une vitesse de conduction nerveuse est comprise entre moins de $\SI{1}{m.s^{-1}}$ et environ $\SI{120}{m.s^{-1}}$.
\end{enumerate}
\end{methode}

\begin{exoresolu}[Vitesse de conduction dans le nerf sciatique]
On stimule le nerf sciatique d'une grenouille et l'on enregistre les PA à l'aide de deux électrodes $E_1$ et $E_2$ distantes de $d = \SI{45}{mm}$. Sur l'oscillogramme, le PA enregistré par $E_2$ apparaît $\SI{1,5}{ms}$ après celui de $E_1$. Calculer la vitesse de conduction de l'influx dans ce nerf.
\tcblower
On applique la relation $v = \dfrac{d}{\Delta t}$ avec $d = \SI{45}{mm} = \SI{45e-3}{m}$ et $\Delta t = \SI{1,5}{ms} = \SI{1,5e-3}{s}$ :
\[ v = \frac{\num{45e-3}}{\num{1,5e-3}} = \SI{30}{m.s^{-1}}. \]
La vitesse de conduction de l'influx nerveux dans ce nerf est de $\SI{30}{m.s^{-1}}$, valeur typique d'un nerf de grenouille contenant des fibres myélinisées de diamètre moyen.
\end{exoresolu}

\textbf{Des vitesses très variables selon les fibres}\par

En répétant la mesure sur des fibres nerveuses de types différents, on obtient des vitesses extrêmement variables : de \textbf{moins de $\SI{1}{m.s^{-1}}$} pour certaines fibres fines (fibres de la douleur lente, fibres du système nerveux végétatif) à \textbf{plus de $\SI{100}{m.s^{-1}}$} pour les grosses fibres motrices ou sensitives des mammifères. D'où viennent ces différences ? Deux facteurs expliquent l'essentiel de cette variation : le \textbf{diamètre de la fibre} et la \textbf{présence de myéline}.

\courssub{Les facteurs de variation de la vitesse de conduction}

\textbf{Premier facteur : le diamètre de la fibre}\par

À structure égale, \textbf{plus le diamètre de la fibre nerveuse est grand, plus la conduction est rapide}. L'explication est physique : l'intérieur de l'axone (axoplasme) oppose une résistance au passage des courants locaux, comme un tuyau étroit s'oppose à l'écoulement de l'eau. Plus l'axone est gros, plus cette résistance interne est faible ; les courants locaux s'étendent alors plus loin le long de la fibre, dépolarisent plus rapidement les régions voisines, et le PA se régénère plus vite.

\begin{exemplebox}[Ordres de grandeur chez les fibres amyéliniques]
Chez le calmar, dont l'axone géant (diamètre voisin de $\SI{500}{\micro m}$) a servi de modèle historique à Hodgkin et Huxley, la conduction atteint environ $\SI{25}{m.s^{-1}}$ malgré l'absence de myéline : c'est le très gros diamètre qui assure la rapidité. Chez l'Homme, une fibre amyélinique fine (diamètre inférieur à $\SI{1}{\micro m}$) conduit à peine à $\num{0,5}$ à $\SI{2}{m.s^{-1}}$.
\end{exemplebox}

\textbf{Deuxième facteur : la gaine de myéline et la conduction saltatoire}\par

Chez les vertébrés, la plupart des fibres nerveuses sont entourées d'une \cle{gaine de myéline}, manchon riche en lipides formé par l'enroulement de cellules gliales (cellules de Schwann dans le système nerveux périphérique, oligodendrocytes dans le système nerveux central). Cette gaine est \emph{interrompue} à intervalles réguliers (tous les $0,2$ à $\SI{2}{mm}$ environ) par des étranglements appelés \cle{nœuds de Ranvier}, où la membrane de l'axone est à nu et très riche en canaux \ce{Na+} voltage-dépendants.

La myéline est un excellent \textbf{isolant électrique} : entre deux nœuds, les courants locaux ne peuvent pas traverser la membrane et se propagent donc très loin et très vite à l'intérieur de l'axone, sans fuite. En revanche, les échanges ioniques (et donc la régénération du PA) ne sont possibles \emph{qu'au niveau des nœuds de Ranvier}. Il en résulte un mode de conduction remarquable :

\begin{propriete}[Conduction saltatoire]
Dans une fibre myélinisée, le potentiel d'action n'est régénéré qu'au niveau des \cle{nœuds de Ranvier} : les courants locaux couplent directement un nœud au nœud suivant, dont ils provoquent la dépolarisation. L'influx semble ainsi « sauter » de nœud en nœud : c'est la \cle{conduction saltatoire} (du latin \emph{saltare}, sauter). Elle est beaucoup plus rapide et moins coûteuse en énergie que la conduction continue des fibres amyéliniques.
\end{propriete}

\begin{popfigure}
\begin{center}
\begin{tikzpicture}[scale=1.0, every node/.style={font=\small}]
% Axone central
\draw[thick, popblue] (0,0.45) -- (13,0.45);
\draw[thick, popblue] (0,-0.45) -- (13,-0.45);
% Manchons de myéline
\fill[popgold!45] (0.3,-0.75) rectangle (3.4,0.75);
\fill[popgold!45] (4.6,-0.75) rectangle (7.7,0.75);
\fill[popgold!45] (8.9,-0.75) rectangle (12.0,0.75);
\draw[thick, popgold] (0.3,-0.75) rectangle (3.4,0.75);
\draw[thick, popgold] (4.6,-0.75) rectangle (7.7,0.75);
\draw[thick, popgold] (8.9,-0.75) rectangle (12.0,0.75);
% Noeuds de Ranvier
\fill[poporange!60] (3.4,-0.5) rectangle (4.6,0.5);
\fill[poporange!60] (7.7,-0.5) rectangle (8.9,0.5);
\fill[poporange!60] (12.0,-0.5) rectangle (12.7,0.5);
\node[popdark, font=\small] at (4.0,1.25) {Nœud de Ranvier};
\node[popdark, font=\small] at (8.3,1.25) {Nœud de Ranvier};
\node[popdark, font=\small] at (1.85,-1.15) {gaine de myéline};
\node[popdark, font=\small] at (6.15,-1.15) {gaine de myéline};
% Sauts de l'influx
\draw[-{Stealth[length=3.5mm]}, very thick, poppurple] (1.85,0) .. controls (3.0,1.9) and (5.0,1.9) .. (6.15,0);
\draw[-{Stealth[length=3.5mm]}, very thick, poppurple] (6.15,0) .. controls (7.3,1.9) and (9.3,1.9) .. (10.45,0);
\draw[-{Stealth[length=3.5mm]}, very thick, poppurple] (10.45,0) .. controls (11.2,1.9) and (12.2,1.9) .. (12.9,0.3);
% PA aux noeuds
\node[popink, font=\small\bfseries] at (4.0,-0.05) {PA};
\node[popink, font=\small\bfseries] at (8.3,-0.05) {PA};
\node[popink, font=\small\bfseries] at (12.35,-0.05) {PA};
% Sens de propagation
\draw[-{Stealth[length=4mm]}, ultra thick, popgreen] (1.0,-1.8) -- (4.5,-1.8);
\node[popgreen, font=\small\bfseries] at (2.75,-2.2) {Sens de propagation};
\end{tikzpicture}
\end{center}
\legende{Conduction saltatoire dans une fibre myélinisée. La gaine de myéline (jaune) isole l'axone entre les nœuds de Ranvier (orange) ; le potentiel d'action n'est régénéré qu'aux nœuds, et les courants locaux (flèches violettes) couplent directement un nœud au suivant : l'influx « saute » de nœud en nœud, d'où une vitesse de conduction très élevée.}
\end{popfigure}

\textbf{Comparaison chiffrée : fibres amyéliniques et myélinisées}\par

Les données expérimentales montrent que, à diamètre égal, une fibre myélinisée conduit \textbf{beaucoup plus vite} qu'une fibre amyélinique, et que la vitesse croît avec le diamètre dans les deux cas. Pour les fibres myélinisées des mammifères, la vitesse de conduction (en $\si{m.s^{-1}}$) est approximativement proportionnelle au diamètre total de la fibre (en $\si{\micro m}$), de l'ordre de $6$ fois le diamètre.

\begin{center}
\begin{tabularx}{\linewidth}{|l|X|X|}
\hline
\rowcolor{popblueL} \textbf{Type de fibre} & \textbf{Diamètre} & \textbf{Vitesse de conduction} \\
\hline
Fibre amyélinique fine (douleur lente, système végétatif) & inférieur à $\SI{1}{\micro m}$ & $\num{0,5}$ à $\SI{2}{m.s^{-1}}$ \\
\hline
Fibre myélinisée de diamètre moyen (sensibilité tactile) & $5$ à $\SI{10}{\micro m}$ & $30$ à $\SI{60}{m.s^{-1}}$ \\
\hline
Grosse fibre myélinisée (motoneurones, sensibilité musculaire) & $12$ à $\SI{20}{\micro m}$ & $80$ à $\SI{120}{m.s^{-1}}$ \\
\hline
\end{tabularx}
\end{center}

\begin{popfigure}
\begin{center}
\begin{tikzpicture}
\begin{axis}[
width=0.85\linewidth, height=7cm,
xlabel={Diamètre de la fibre ($\si{\micro m}$)},
ylabel={Vitesse de conduction ($\si{m.s^{-1}}$)},
xmin=0, xmax=22, ymin=0, ymax=130,
legend pos=north west,
legend style={font=\small},
grid=major, grid style={popline!40},
axis lines=left,
]
\addplot[popcurve, domain=1:20, samples=100, popblue, very thick] {6*x};
\addlegendentry{Fibres myélinisées}
\addplot[popcurve, domain=0.2:2, samples=100, poporange, very thick] {1.4*x};
\addlegendentry{Fibres amyéliniques}
\end{axis}
\end{tikzpicture}
\end{center}
\legende{Vitesse de conduction en fonction du diamètre de la fibre nerveuse. La vitesse croît avec le diamètre dans les deux cas, mais à diamètre égal la fibre myélinisée (bleu) conduit beaucoup plus vite que la fibre amyélinique (orange) : la myéline multiplie la vitesse par un facteur qui peut atteindre $50$ à $100$.}
\end{popfigure}

\begin{aretenir}[L'essentiel de la section]
\begin{itemize}
\item L'\cle{influx nerveux} est le potentiel d'action qui se propage le long de la fibre nerveuse.
\item La propagation se fait \textbf{de proche en proche} : les courants locaux issus de la zone active dépolarisent la région voisine jusqu'au seuil, qui génère à son tour un PA de même amplitude. Le signal ne s'affaiblit donc jamais.
\item \emph{In vivo}, la propagation est \textbf{unidirectionnelle} grâce à la période réfractaire, qui rend inexcitable la région que l'influx vient de quitter.
\item La vitesse de conduction se mesure par $v = \dfrac{d}{\Delta t}$ ; elle varie de moins de $\SI{1}{m.s^{-1}}$ à plus de $\SI{100}{m.s^{-1}}$.
\item Elle augmente avec le \textbf{diamètre} de la fibre et surtout avec la \textbf{myélinisation} : dans les fibres myélinisées, la conduction est \textbf{saltatoire}, le PA « sautant » de nœud de Ranvier en nœud de Ranvier.
\end{itemize}
\end{aretenir}

\begin{savaistu}[La sclérose en plaques : quand la myéline disparaît]
La \cle{sclérose en plaques} est une maladie auto-immune au cours de laquelle le système immunitaire du patient détruit la gaine de myéline des fibres nerveuses du système nerveux central. Privées de leur isolant, les fibres perdent la conduction saltatoire : l'influx est ralenti, déformé, parfois totalement bloqué. Il en résulte des troubles variés — paralysies, troubles de la sensibilité, de la vision ou de l'équilibre — selon les régions démyélinisées. Cette maladie illustre de façon frappante le rôle fonctionnel de la myéline : sans elle, la rapidité et la fidélité de la communication nerveuse s'effondrent. Des cas sont diagnostiqués dans les hôpitaux camerounais, où l'imagerie médicale (IRM) permet de visualiser les plaques de démyélinisation.
\end{savaistu}

\begin{exercice}[Exploitation d'un oscillogramme]
Sur un nerf sciatique de grenouille, deux électrodes d'enregistrement sont placées à $\SI{60}{mm}$ l'une de l'autre. Après stimulation, l'oscillogramme montre deux PA séparés par un décalage de $\SI{2}{ms}$. (1) Calculer la vitesse de conduction de l'influx. (2) L'amplitude du second PA est égale à celle du premier : que peut-on en conclure sur le mode de propagation ? (3) Ce nerf contient-il vraisemblablement des fibres myélinisées ? Justifier.
\end{exercice}

\begin{corrige}[Exercice : Exploitation d'un oscillogramme]
(1) $v = \dfrac{d}{\Delta t} = \dfrac{\num{60e-3}}{\num{2e-3}} = \SI{30}{m.s^{-1}}$. (2) L'égalité des amplitudes montre que le PA se régénère intégralement de proche en proche : la propagation est active et non passive, le signal ne s'affaiblit pas avec la distance. (3) Oui : une vitesse de $\SI{30}{m.s^{-1}}$ est incompatible avec des fibres amyéliniques fines ($\num{0,5}$ à $\SI{2}{m.s^{-1}}$) ; elle traduit la présence de fibres myélinisées à conduction saltatoire.
\end{corrige}

Ainsi, le neurone n'est pas seulement capable de générer un potentiel d'action : il le fait voyager, vite et sans perte, jusqu'à son extrémité. Mais que devient ce message lorsqu'il arrive au bout de l'axone, face à un autre neurone ou à une cellule effectrice ? C'est toute la question de la \textbf{synapse}, que nous aborderons dans la section suivante.

\coursec{La synapse : types et fonctionnement de la synapse chimique}

Dans la section précédente, nous avons suivi le message nerveux le long d'une fibre : le potentiel d'action se propage de proche en proche, plus ou moins vite selon le diamètre de la fibre et la présence de myéline. Mais un neurone ne fonctionne jamais seul : l'influx arrivé au bout de l'axone doit être \emph{transmis} à une autre cellule — un autre neurone, une cellule musculaire ou une cellule glandulaire. Se pose alors une question fondamentale : le message nerveux passe-t-il \cle{directement}, sous forme électrique, d'un neurone à l'autre, comme le courant le long d'un fil ? Ou bien existe-t-il une étape intermédiaire ? C'est toute la question de la \cle{synapse}, que nous allons résoudre en suivant la démarche des chercheurs qui l'ont élucidée.

\courssub{Le problème : continuité ou discontinuité entre les neurones ?}

\textbf{Observation.} À la fin du XIX\textsuperscript{e} siècle, l'histologiste espagnol Santiago Ramón y Cajal, utilisant la coloration de Golgi, observe au microscope que les neurones du tissu nerveux ne se touchent pas réellement : entre l'extrémité d'un axone et la cellule voisine subsiste un très fin \emph{intervalle}. Le tissu nerveux serait donc fait de cellules \emph{discontinues}.

\textbf{Problème.} Si les neurones ne sont pas en continuité, comment le message nerveux franchit-il l'intervalle qui les sépare ? Deux hypothèses s'affrontent au début du XX\textsuperscript{e} siècle :
\begin{itemize}
\item \textbf{Hypothèse électrique} : le potentiel d'action se transmet directement, par couplage électrique, d'une cellule à l'autre ;
\item \textbf{Hypothèse chimique} : le neurone libère une \emph{substance chimique} qui diffuse jusqu'à la cellule suivante et l'excite.
\end{itemize}

C'est une expérience célèbre, aussi simple qu'élégante, qui va trancher.

\courssub{L'expérience d'Otto Loewi (1921) : la preuve d'une transmission chimique}

\begin{experience}[L'expérience des deux cœurs de grenouille (Otto Loewi, 1921)]
\textbf{Principe.} Le cœur de grenouille continue de battre plusieurs heures hors de l'organisme s'il est maintenu dans du liquide de Ringer (solution physiologique). Son rythme est ralenti par la stimulation du \cle{nerf vague} (nerf pneumogastrique, X\textsuperscript{e} paire de nerfs crâniens).

\textbf{Protocole.}
\begin{enumerate}
\item Un premier cœur de grenouille (cœur A) est placé dans une cuve contenant du liquide de Ringer, avec son nerf vague laissé en place ;
\item On stimule électriquement le nerf vague du cœur A : on observe un \textbf{ralentissement} des battements du cœur A ;
\item On recueille alors le liquide ayant baigné le cœur A et on le transfère sur un second cœur (cœur B), \emph{dénervé} (dont on a sectionné le nerf vague) ;
\item On observe que le cœur B \textbf{ralentit à son tour}, alors qu'il n'a reçu aucune stimulation électrique.
\end{enumerate}

\textbf{Observations.} Le cœur B ralentit sans avoir été stimulé : le seul lien entre les deux cœurs est le liquide transféré.

\textbf{Interprétation.} La stimulation du nerf vague a provoqué la \emph{libération d'une substance chimique} par les terminaisons nerveuses dans le liquide de Ringer. Cette substance, transférée avec le liquide, a agi sur le cœur B. Loewi l'appela \emph{Vagusstoff} (« substance du vague ») ; il s'agit de l'\cle{acétylcholine}.

\textbf{Conclusion.} La transmission du message nerveux d'un neurone vers sa cellule cible est de nature \textbf{chimique} : elle fait intervenir une substance médiatrice, le neurotransmetteur. Otto Loewi reçut le prix Nobel de physiologie ou médecine en 1936 pour cette découverte (partagé avec Henry Dale).
\end{experience}

\begin{popfigure}
\begin{center}
\begin{tikzpicture}[scale=0.9, font=\small]
% Cuve A
\draw[thick, popblue] (-1.2,0) -- (-1.2,2.6) -- (2.6,2.6) -- (2.6,0);
\draw[popblue!40, fill=popblue!15] (-1.1,0.05) rectangle (2.5,1.7);
\node[popblue] at (0.7,2.85) {\textbf{Cœur A (innervé)}};
% Coeur A
\draw[thick, poppink, fill=poppinkL] (0.7,0.9) ellipse (0.55 and 0.7);
\draw[thick, poppink] (0.7,1.6) .. controls (0.5,2.1) and (0.9,2.1) .. (0.7,1.6);
% nerf vague
\draw[thick, popgold] (0.7,1.55) -- (0.7,3.4);
\node[popgold, right] at (0.75,3.3) {nerf vague};
% électrode
\draw[thick, popdark] (1.6,3.6) -- (0.85,3.15);
\node[popdark, above right] at (1.55,3.55) {électrode (stimulation)};
% flèche transfert
\draw[-{Stealth[length=3mm]}, very thick, poporange] (2.9,1.0) -- (4.6,1.0);
\node[poporange, align=center] at (3.75,0.55) {transfert du\\liquide de Ringer};
% Cuve B
\draw[thick, popblue] (5.0,0) -- (5.0,2.6) -- (8.8,2.6) -- (8.8,0);
\draw[popblue!40, fill=popblue!15] (5.1,0.05) rectangle (8.7,1.7);
\node[popblue] at (6.9,2.85) {\textbf{Cœur B (dénervé)}};
\draw[thick, poppink, fill=poppinkL] (6.9,0.9) ellipse (0.55 and 0.7);
\draw[thick, poppink] (6.9,1.6) .. controls (6.7,2.1) and (7.1,2.1) .. (6.9,1.6);
% nerf sectionné
\draw[thick, popgold, dashed] (6.9,1.55) -- (6.9,2.3);
\draw[thick, popdark] (6.6,2.35) -- (7.2,2.55);
\node[popmuted, right] at (7.0,2.45) {nerf sectionné};
% résultats
\node[align=center, popdark] at (0.7,-0.6) {stimulation $\Rightarrow$\\ralentissement du cœur A};
\node[align=center, popdark] at (6.9,-0.6) {aucune stimulation, mais\\ralentissement du cœur B};
\end{tikzpicture}
\end{center}
\legende{Dispositif de l'expérience d'Otto Loewi (1921). La stimulation du nerf vague du cœur A libère une substance (l'acétylcholine) dans le liquide de Ringer ; transféré sur le cœur B dénervé, ce liquide ralentit ses battements : la transmission nerveuse est de nature chimique.}
\end{popfigure}

\begin{aretenir}[Ce qu'il faut retenir de l'expérience de Loewi]
Le passage du message d'un neurone à sa cellule cible n'est pas un simple phénomène électrique : il met en jeu la \cle{libération d'une substance chimique} (neurotransmetteur) qui diffuse et agit sur la cellule réceptrice. C'est la naissance du concept de \emph{transmission synaptique chimique}.
\end{aretenir}

\courssub{Définition et diversité des synapses}

\begin{definition}[Synapse]
Une \cle{synapse} est une \textbf{zone de contact fonctionnel} assurant la transmission du message nerveux d'un neurone (élément \emph{présynaptique}) vers une autre cellule (élément \emph{postsynaptique}) : un autre neurone, une cellule musculaire ou une cellule glandulaire. Les deux membranes ne se touchent pas : elles sont séparées par un très fin espace, la \emph{fente synaptique} (largeur d'environ \SI{20}{nm} à \SI{40}{nm}).
\end{definition}

\begin{definition}[Neurotransmetteur]
Un \cle{neurotransmetteur} (ou médiateur chimique) est une \textbf{substance chimique libérée au niveau de la synapse} par l'élément présynaptique et assurant la transmission du signal nerveux vers l'élément postsynaptique, dont il modifie l'activité. Exemples : l'\emph{acétylcholine} (cœur, muscles squelettiques), l'\emph{adrénaline}, la \emph{dopamine}, la \emph{sérotonine}, le \emph{GABA}.
\end{definition}

\textbf{Classification selon les structures en contact.} Selon la partie du neurone postsynaptique sur laquelle aboutit l'axone, on distingue trois types de synapses interneuronales :
\begin{itemize}
\item la synapse \cle{axo-dendritique} : l'axone du neurone présynaptique s'articule avec une \emph{dendrite} du neurone postsynaptique (type le plus fréquent) ;
\item la synapse \cle{axo-somatique} : l'axone aboutit sur le \emph{corps cellulaire} (soma) du neurone postsynaptique ;
\item la synapse \cle{axo-axonique} : l'axone aboutit sur l'\emph{axone} (souvent le bouton synaptique) d'un autre neurone, ce qui permet de moduler la libération de neurotransmetteur.
\end{itemize}

\textbf{Classification selon la nature de la cellule postsynaptique.} On distingue :
\begin{itemize}
\item les synapses \cle{interneuronales}, entre deux neurones (dans les centres nerveux) ;
\item les synapses \cle{neuromusculaires}, entre un motoneurone et une fibre musculaire squelettique : elles constituent la \emph{plaque motrice}, dont le neurotransmetteur est l'acétylcholine ;
\item les synapses \emph{neuroglandulaires}, entre un neurone et une cellule glandulaire (ex. libération d'adrénaline par la médullosurrénale).
\end{itemize}

\begin{popfigure}
\begin{center}
\begin{tikzpicture}[scale=0.85, font=\small]
% ---------- axo-dendritique ----------
\node[popdark, font=\bfseries] at (1.4,4.6) {Axo-dendritique};
% dendrite + soma postsynaptique
\draw[thick, popblue, fill=popblueL] (0.3,0) circle (0.55);
\node[popblue] at (0.3,0) {soma};
\draw[thick, popblue] (0.55,0.35) -- (1.5,1.3) -- (1.9,2.2);
\draw[thick, popblue] (1.5,1.3) -- (2.3,1.6);
\node[popblue, right] at (2.3,1.6) {dendrite};
% axone présynaptique
\draw[thick, poporange] (1.9,4.0) -- (1.9,2.9);
\draw[thick, poporange, fill=poporangeL] (1.9,2.55) ellipse (0.35 and 0.4);
\node[poporange, left] at (1.55,3.6) {axone};
% ---------- axo-somatique ----------
\node[popdark, font=\bfseries] at (5.6,4.6) {Axo-somatique};
\draw[thick, popblue, fill=popblueL] (5.6,0) circle (0.55);
\draw[thick, popblue] (5.85,0.4) -- (6.4,1.2);
\draw[thick, poporange] (4.9,4.0) -- (4.9,1.6);
\draw[thick, poporange, fill=poporangeL] (4.9,1.15) ellipse (0.35 and 0.4);
% ---------- axo-axonique ----------
\node[popdark, font=\bfseries] at (9.9,4.6) {Axo-axonique};
\draw[thick, popblue, fill=popblueL] (8.6,0) circle (0.5);
\draw[thick, popblue] (9.1,0.1) -- (10.6,0.1);
\draw[thick, popblue, fill=popblueL] (10.9,0.1) ellipse (0.3 and 0.35);
\node[popblue] at (10.9,-0.7) {bouton};
\draw[thick, poporange] (10.9,3.6) -- (10.9,1.4);
\draw[thick, poporange, fill=poporangeL] (10.9,1.0) ellipse (0.32 and 0.38);
\node[poporange, right] at (11.25,2.6) {axone modulateur};
\draw[thick, popblue] (8.6,0.5) -- (8.6,1.5);
\node[popblue, left] at (8.3,1.2) {axone};
% légende commune
\draw[thick, poporange, fill=poporangeL] (0.3,-1.6) ellipse (0.25 and 0.18);
\node[right, popdark] at (0.7,-1.6) {élément présynaptique};
\draw[thick, popblue, fill=popblueL] (5.3,-1.6) circle (0.2);
\node[right, popdark] at (5.65,-1.6) {élément postsynaptique};
\end{tikzpicture}
\end{center}
\legende{Les trois types de synapses interneuronales selon le point de contact : axo-dendritique (sur une dendrite), axo-somatique (sur le corps cellulaire), axo-axonique (sur l'axone ou le bouton synaptique d'un autre neurone).}
\end{popfigure}

\textbf{Synapse chimique et synapse électrique.} On distingue enfin deux modalités de transmission :
\begin{itemize}
\item la \cle{synapse électrique} : les deux cellules sont reliées par des canaux de jonction (\emph{gap junctions}) ; le courant ionique passe directement, sans délai notable, et parfois dans les deux sens. Elle est rare chez les mammifères (quelques centres nerveux, muscle cardiaque) ;
\item la \cle{synapse chimique} : la transmission se fait par l'intermédiaire d'un neurotransmetteur. C'est le type \textbf{le plus répandu} dans le système nerveux des vertébrés ; c'est elle que nous étudions en détail.
\end{itemize}

\courssub{Structure de la synapse chimique}

L'observation au microscope électronique révèle une organisation très stéréotypée, commune à toutes les synapses chimiques :

\begin{itemize}
\item l'extrémité de l'axone présynaptique forme un renflement, le \cle{bouton synaptique}, qui contient de nombreuses \cle{vésicules synaptiques} remplies de neurotransmetteur, ainsi que des mitochondries (fournissant l'ATP nécessaire) ;
\item la \cle{membrane présynaptique} est la face du bouton tournée vers la cellule cible ; elle porte des \emph{canaux calciques} voltage-dépendants ;
\item la \cle{fente synaptique} (ou espace synaptique), d'environ \SI{20}{nm} à \SI{40}{nm}, sépare les deux membranes ;
\item la \cle{membrane postsynaptique} porte des \cle{récepteurs} spécifiques : des protéines capables de fixer le neurotransmetteur et couplées à des canaux ioniques.
\end{itemize}

\begin{methode}[Schématiser une synapse chimique]
Pour réaliser un schéma correct et complet d'une synapse à transmission chimique :
\begin{enumerate}
\item Représenter le \textbf{bouton synaptique} (renflement de l'axone présynaptique) contenant des \textbf{vésicules synaptiques} et des mitochondries ;
\item Dessiner la \textbf{fente synaptique}, espace étroit entre les deux membranes ;
\item Représenter la \textbf{membrane postsynaptique} porteuse de \textbf{récepteurs} (petites protéines enchâssées) ;
\item Figurer le \textbf{neurotransmetteur} (petits points) en cours de libération dans la fente ;
\item \textbf{Flécher le sens de transmission} (présynaptique $\rightarrow$ postsynaptique) et légender chaque structure.
\end{enumerate}
\end{methode}

\begin{popfigure}
\begin{center}
\begin{tikzpicture}[scale=1.0, font=\small]
% axone + bouton présynaptique
\draw[thick, poporange, fill=poporangeL] (-4.2,2.6) -- (-2.2,2.6) .. controls (-0.6,2.7) and (-0.4,2.2) .. (-0.4,1.2) .. controls (-0.4,0.35) and (-0.9,0.1) .. (-2.0,0.1) .. controls (-3.1,0.1) and (-3.6,0.35) .. (-3.6,1.2) .. controls (-3.6,2.2) and (-3.4,2.7) .. (-4.2,2.6) -- cycle;
\node[poporange, above] at (-3.2,2.65) {axone présynaptique};
% vésicules
\foreach \x/\y in {-2.9/1.9, -2.2/2.1, -1.5/1.9, -2.6/1.3, -1.8/1.2, -1.15/1.35}{
  \draw[thick, poppurple, fill=poppurpleL] (\x,\y) circle (0.22);
  \fill[poppurple] (\x,\y) circle (0.05);}
\node[poppurple, right] at (-0.55,2.15) {vésicules synaptiques};
\draw[-{Stealth[length=2mm]}, poppurple] (-0.6,2.05) -- (-1.35,1.95);
% mitochondrie
\draw[thick, popgold, fill=popgoldL] (-3.1,0.7) ellipse (0.42 and 0.22);
\draw[popgold] (-3.35,0.7) .. controls (-3.2,0.85) and (-3.05,0.55) .. (-2.9,0.7);
\node[popgold, left] at (-3.6,0.7) {mitochondrie};
% membrane présynaptique (label)
\node[poporange, below left] at (-3.6,0.05) {membrane présynaptique};
% fente synaptique
\fill[popblue!12] (-4.2,-0.35) rectangle (0.2,-0.85);
\node[popblue] at (-2.0,-0.6) {fente synaptique ($\approx$ \SI{20}{nm})};
% neurotransmetteur dans la fente
\foreach \x/\y in {-2.7/-0.55, -2.1/-0.7, -1.5/-0.5, -0.9/-0.68, -2.35/-0.45}{
  \fill[poppurple] (\x,\y) circle (0.06);}
\node[poppurple, right] at (0.3,-0.5) {neurotransmetteur};
\draw[-{Stealth[length=2mm]}, poppurple] (0.28,-0.62) -- (-0.75,-0.62);
% membrane postsynaptique
\draw[very thick, popgreen] (-4.2,-0.85) -- (0.2,-0.85);
% récepteurs
\foreach \x in {-3.3, -2.0, -0.7}{
  \draw[thick, popgreen, fill=popgreenL] (\x,-0.85) ++(-0.14,-0.42) rectangle ++(0.28,0.5);
  \draw[thick, popgreen] (\x,-0.35) ++(-0.14,0) -- ++(-0.1,0.18) (\x,-0.35) ++(0.14,0) -- ++(0.1,0.18);}
\node[popgreen, right] at (0.3,-1.35) {récepteurs postsynaptiques};
\draw[-{Stealth[length=2mm]}, popgreen] (0.28,-1.45) -- (-0.55,-1.45);
% cellule postsynaptique
\draw[thick, popgreen, fill=popgreenL!60] (-4.2,-0.85) -- (-4.2,-2.3) -- (0.2,-2.3) -- (0.2,-0.85);
\draw[very thick, popgreen] (-4.2,-0.85) -- (0.2,-0.85);
\node[popgreen!70!black] at (-2.0,-1.9) {cellule postsynaptique (neurone, muscle, glande)};
% flèche sens de transmission
\draw[-{Stealth[length=4mm]}, very thick, popdark] (-4.9,2.2) -- (-4.9,-1.8);
\node[popdark, rotate=90, above] at (-5.15,-0.2) {sens de transmission};
\end{tikzpicture}
\end{center}
\legende{Schéma d'une synapse à transmission chimique. Le bouton synaptique contient les vésicules de neurotransmetteur ; la fente synaptique sépare la membrane présynaptique de la membrane postsynaptique, qui porte les récepteurs spécifiques. La transmission se fait dans un seul sens.}
\end{popfigure}

\courssub{Fonctionnement de la synapse chimique, étape par étape}

Le déroulement de la transmission synaptique, établi par l'expérimentation (électrophysiologie, microscopie électronique, biochimie), comporte une succession d'étapes rapides :

\begin{enumerate}
\item \textbf{Arrivée du potentiel d'action} au bouton synaptique : la dépolarisation atteint l'extrémité de l'axone présynaptique ;
\item \textbf{Ouverture des canaux calciques} voltage-dépendants de la membrane présynaptique : les ions \ce{Ca^{2+}}, plus concentrés à l'extérieur, \emph{entrent} massivement dans le bouton synaptique ;
\item \textbf{Exocytose des vésicules} : l'élévation de la concentration intracellulaire en \ce{Ca^{2+}} provoque la fusion des vésicules synaptiques avec la membrane présynaptique ;
\item \textbf{Libération du neurotransmetteur} (par exemple l'acétylcholine) dans la fente synaptique, où il diffuse ;
\item \textbf{Fixation sur les récepteurs} spécifiques de la membrane postsynaptique : cette fixation ouvre des canaux ioniques, d'où une modification de la perméabilité membranaire ;
\item \textbf{Apparition d'un potentiel postsynaptique} : variation locale et graduée du potentiel de membrane de la cellule postsynaptique. S'il s'agit d'une dépolarisation (potentiel postsynaptique \emph{excitateur}, PPSE) atteignant le seuil, un potentiel d'action naît dans la cellule postsynaptique ; s'il s'agit d'une hyperpolarisation (potentiel \emph{inhibiteur}, PPSI), la cellule est au contraire freinée ;
\item \textbf{Élimination du neurotransmetteur} : pour que le signal soit bref et que la synapse redevienne fonctionnelle, le neurotransmetteur est rapidement \emph{dégradé} par une enzyme (l'acétylcholinestérase hydrolyse l'acétylcholine) ou \emph{recapturé} par l'élément présynaptique (recapture de la dopamine, de la sérotonine).
\end{enumerate}

\begin{propriete}[Deux conséquences fondamentales du fonctionnement synaptique]
\begin{itemize}
\item \textbf{Le délai synaptique} : les étapes chimiques (diffusion de \ce{Ca^{2+}}, exocytose, diffusion du neurotransmetteur, fixation sur les récepteurs) prennent du temps. Le passage d'une synapse retarde le message d'environ \SI{0,5}{ms} à \SI{1}{ms}. Ainsi, un réflexe dont l'arc comporte plusieurs synapses est plus lent qu'une simple conduction le long d'une fibre ;
\item \textbf{La transmission unidirectionnelle} : le message ne franchit la synapse que dans le sens présynaptique $\rightarrow$ postsynaptique, jamais en sens inverse.
\end{itemize}
\end{propriete}

\begin{attention}[Pourquoi la transmission est-elle unidirectionnelle ?]
Ne dis pas simplement « la synapse transmet dans un seul sens » : il faut le \emph{justifier}. La transmission est unidirectionnelle car la répartition des équipements est \textbf{asymétrique} : \emph{seule la membrane présynaptique} possède les vésicules et peut libérer le neurotransmetteur, et \emph{seule la membrane postsynaptique} porte les récepteurs capables de le reconnaître. La cellule postsynaptique ne peut donc pas exciter en retour l'élément présynaptique. Autre piège fréquent : le neurotransmetteur ne traverse \emph{pas} la membrane postsynaptique ; il se fixe à sa surface sur les récepteurs, puis est dégradé ou recapturé.
\end{attention}

\begin{exemplebox}[Le délai synaptique, un chiffre à connaître]
Le délai synaptique est d'environ \SI{0,5}{ms} à \SI{1}{ms} par synapse. Considérons un arc réflexe comportant 3 synapses et un trajet total de fibres de \SI{1,5}{m}, parcouru à la vitesse de \SI{50}{m/s} : la conduction pure prend $\dfrac{\num{1,5}}{50} = \SI{0,03}{s} = \SI{30}{ms}$, et les synapses ajoutent environ $3 \times \SI{1}{ms} = \SI{3}{ms}$. Le délai synaptique représente donc près de 10\,\% du temps de réaction : c'est pourquoi un réflexe impliquant de nombreuses synapses (réflexe polysynaptique) est nettement plus lent qu'une simple conduction.
\end{exemplebox}

\begin{exoresolu}[Interpréter une expérience sur la synapse neuromusculaire]
\textbf{Énoncé.} Sur une préparation nerf--muscle de grenouille, on place la jonction entre le nerf moteur et le muscle dans un milieu dépourvu d'ions \ce{Ca^{2+}}. On stimule le nerf : le potentiel d'action parvient normalement au bouton synaptique (vérifié par électrode), mais le muscle ne se contracte plus. Interpréter ce résultat et conclure sur le rôle des ions \ce{Ca^{2+}}.
\tcblower
\textbf{Solution détaillée.}
\begin{enumerate}
\item \textbf{Observation} : en l'absence de \ce{Ca^{2+}} extracellulaire, l'influx atteint le bouton synaptique mais la transmission au muscle est abolie.
\item \textbf{Analyse} : le blocage ne se situe ni dans la conduction (le PA arrive au bouton) ni dans le muscle (il peut se contracter si on le stimule directement), mais au niveau de la \emph{jonction}.
\item \textbf{Interprétation} : l'entrée des ions \ce{Ca^{2+}} dans le bouton synaptique, déclenchée par le potentiel d'action, est l'étape indispensable qui provoque la fusion des vésicules avec la membrane présynaptique et l'\emph{exocytose de l'acétylcholine}. Sans \ce{Ca^{2+}}, pas de libération de neurotransmetteur, donc pas de potentiel postsynaptique et pas de contraction.
\item \textbf{Conclusion} : les ions \ce{Ca^{2+}} sont le \emph{couplage} entre le signal électrique (potentiel d'action) et le signal chimique (libération du neurotransmetteur) : on parle de couplage excitation--sécrétion.
\end{enumerate}
\end{exoresolu}

\begin{savaistu}[De la grenouille de Loewi aux hôpitaux camerounais]
La découverte de la transmission chimique a des applications médicales immenses. De nombreux médicaments et poisons agissent au niveau des synapses : les organophosphorés (certains pesticides utilisés en agriculture) inhibent l'acétylcholinestérase, provoquant une accumulation toxique d'acétylcholine ; le curare, poison fléché d'Amazonie, bloque les récepteurs de l'acétylcholine et paralyse les muscles. À l'inverse, les anesthésiques et myorelaxants utilisés au bloc opératoire (par exemple au CHU de Yaoundé) exploitent ces mêmes mécanismes pour contrôler la transmission neuromusculaire. Nous détaillerons ces effets dans la section suivante.
\end{savaistu}

\begin{aretenir}[L'essentiel de la section]
\begin{itemize}
\item Une \cle{synapse} est une zone de contact fonctionnel entre un neurone et une cellule cible ; l'expérience de Loewi (1921) a prouvé que la transmission y est \emph{chimique} ;
\item On distingue les synapses axo-dendritiques, axo-somatiques et axo-axoniques ; interneuronales ou neuromusculaires (plaque motrice) ; chimiques (les plus répandues) ou électriques ;
\item La synapse chimique comprend : bouton synaptique à vésicules, membrane présynaptique, fente synaptique, membrane postsynaptique à récepteurs ;
\item Chaîne des événements : PA $\rightarrow$ entrée de \ce{Ca^{2+}} $\rightarrow$ exocytose $\rightarrow$ libération du neurotransmetteur $\rightarrow$ fixation sur les récepteurs $\rightarrow$ potentiel postsynaptique $\rightarrow$ dégradation ou recapture ;
\item Il en résulte un \emph{délai synaptique} (\SI{0,5}{ms} à \SI{1}{ms}) et une transmission \emph{unidirectionnelle}.
\end{itemize}
\end{aretenir}

\coursec{Effets de quelques substances sur la transmission synaptique}

Dans la section précédente, nous avons décortiqué le fonctionnement de la synapse à transmission chimique : arrivée du potentiel d'action au bouton présynaptique, entrée de \ce{Ca^{2+}}, exocytose du neurotransmetteur, fixation sur les récepteurs postsynaptiques, genèse d'un potentiel postsynaptique, puis élimination rapide du médiateur chimique. Chacune de ces étapes est une cible potentielle pour des substances étrangères à l'organisme : poisons, médicaments, drogues, pesticides. La question qui guide cette dernière section est donc la suivante : \emph{comment certaines substances modifient-elles la transmission synaptique, et comment le démontre-t-on expérimentalement ?}

\courssub{Problème et hypothèses : où une substance peut-elle agir ?}

\begin{experience}[Observation préalable]
Lors d'interventions chirurgicales au Centre Hospitalier d'Essos à Yaoundé, les anesthésistes injectent au patient des substances qui provoquent un relâchement complet des muscles, sans pour autant endormir la sensibilité générale par ce seul produit. Inversement, un agriculteur de la plaine des Mbo exposé sans protection à certains insecticides présente des contractions musculaires involontaires prolongées (fasciculations), une salivation abondante et des troubles cardiaques. Dans les deux cas, le système nerveux est intact, les muscles sont intacts : c'est la \cle{communication synaptique} qui est perturbée, mais en sens opposé (blocage dans un cas, hyperactivation dans l'autre).
\end{experience}

Rappelons les étapes de la transmission chimique, vues à la section 5 :
\begin{enumerate}
\item libération du neurotransmetteur par exocytose (dépendante du \ce{Ca^{2+}}) ;
\item diffusion du neurotransmetteur dans la fente synaptique ;
\item fixation du neurotransmetteur sur les récepteurs postsynaptiques ;
\item genèse du potentiel postsynaptique ;
\item élimination du neurotransmetteur (hydrolyse enzymatique, ex. acétylcholinestérase, ou recapture).
\end{enumerate}

\textbf{Hypothèses de travail :} une substance active peut agir en bloquant les récepteurs postsynaptiques (étape 3), en empêchant la libération du neurotransmetteur (étape 1), ou en inhibant son élimination (étape 5). Nous allons tester ces hypothèses à travers l'étude de quelques substances emblématiques.

\courssub{Le curare : un bloqueur des récepteurs postsynaptiques de l'acétylcholine}

\begin{definition}[Curare]
Le \cle{curare} est un poison extrait de lianes sud-américaines (genre \emph{Strychnos}, \emph{Chondrodendron}). C'est un \cle{antagoniste compétitif} de l'acétylcholine : sa molécule se fixe sur les récepteurs cholinergiques nicotiniques de la membrane postsynaptique de la \emph{jonction neuromusculaire} (plaque motrice), sans les activer, et empêche ainsi l'acétylcholine de s'y fixer.
\end{definition}

\begin{experience}[Expérience classique : le nerf sciatique et le muscle gastrocnémien de grenouille]
\textbf{Matériel :} une préparation isolée \og nerf sciatique -- muscle gastrocnémien \fg{} de grenouille, maintenue en vie dans du liquide de Ringer ; un stimulateur électrique ; un myographe (dispositif enregistrant les contractions musculaires).

\textbf{Protocole et résultats :}
\begin{enumerate}
\item \emph{Témoin :} on stimule électriquement le nerf sciatique $\Rightarrow$ le muscle se contracte (transmission neuromusculaire normale).
\item On applique une solution de curare sur la \emph{plaque motrice} (zone de contact nerf--muscle), puis on stimule à nouveau le nerf $\Rightarrow$ \textbf{le muscle ne répond plus}.
\item On stimule alors \emph{directement} le muscle curarisé (électrode posée sur le muscle) $\Rightarrow$ \textbf{le muscle se contracte normalement}.
\item Enfin, on vérifie qu'un potentiel d'action se propage toujours le long du nerf curarisé (enregistrement électrique) $\Rightarrow$ la conduction nerveuse est \emph{intacte}.
\end{enumerate}

\textbf{Interprétation :} le muscle est fonctionnel (appareil contractile intact), le nerf conduit normalement l'influx (conduction intacte) ; seule la \emph{transmission au niveau de la synapse neuromusculaire} est abolie. Le curare agit donc au niveau de la fente synaptique. Des expériences complémentaires (injection d'acétylcholine sur le muscle curarisé : pas de réponse) montrent que c'est la \emph{fixation de l'acétylcholine sur ses récepteurs} qui est bloquée (étape 3 de la liste ci-dessus).

\textbf{Conclusion :} le curare bloque les récepteurs postsynaptiques de l'acétylcholine ; le neurotransmetteur est libéré normalement mais ne peut plus agir $\Rightarrow$ \cle{paralysie musculaire} flasque, y compris des muscles respiratoires (diaphragme), d'où la mort par asphyxie en cas d'intoxication.
\end{experience}

\begin{attention}[Piège classique du Bac]
Le curare ne bloque \textbf{pas} la conduction de l'influx le long de la fibre nerveuse, et ne paralyse \textbf{pas} le muscle lui-même : il bloque uniquement la \emph{transmission} au niveau de la synapse neuromusculaire. La preuve expérimentale exigible : le muscle curarisé stimulé \emph{directement} répond encore, et le nerf curarisé conduit toujours les potentiels d'action. Écrire \og le curare empêche la propagation du message nerveux le long du nerf \fg{} est une erreur rédhibitoire.
\end{attention}

\begin{exemplebox}[Du poison de flèche au myorelaxant]
Le curare est utilisé depuis des siècles comme \emph{poison de flèche} par certains peuples d'Amazonie pour la chasse : l'animal touché est paralysé mais sa viande reste comestible (le curare n'est pas actif par voie digestive). À dose contrôlée, ses dérivés (ex. \emph{tubocurarine}) sont utilisés en médecine comme \cle{myorelaxants} lors d'interventions chirurgicales, pour obtenir le relâchement des muscles sous ventilation assistée. Même molécule, même cible : c'est la dose et le contexte qui font le poison ou le médicament.
\end{exemplebox}

\begin{popfigure}
\begin{center}
\begin{tikzpicture}[scale=0.92, every node/.style={font=\small}]
% ===== Synapse normale (gauche) =====
\begin{scope}
\fill[popblueL] (-0.5,0.5) rectangle (5.5,3.6);
\node[popdark, font=\small\bfseries] at (2.5,4.0) {Synapse normale};
\node[popmuted, font=\footnotesize] at (2.5,3.35) {bouton présynaptique};
% vésicules
\foreach \x/\y in {0.6/2.6, 1.8/2.7, 3.0/2.5, 4.2/2.7}{
  \draw[popblue, fill=white] (\x,\y) circle (0.28);
  \foreach \a in {45,135,225,315}{
    \fill[poporange] ({\x+0.13*cos(\a)},{\y+0.13*sin(\a)}) circle (0.05);}}
% ACh libérée
\foreach \x/\y in {1.0/0.9, 2.2/1.1, 3.4/0.85, 4.4/1.15}{
  \fill[poporange] (\x,\y) circle (0.09);}
% fente
\node[popmuted, font=\footnotesize] at (2.5,0.35) {fente synaptique};
% membrane postsynaptique
\fill[popgreenL] (-0.5,-1.6) rectangle (5.5,0.0);
\node[popmuted, font=\footnotesize] at (4.6,-1.35) {muscle};
% récepteurs
\foreach \x in {1.0, 2.2, 3.4, 4.4}{
  \draw[popgreen, line width=1.2pt] (\x-0.12,0) -- (\x-0.12,0.25) -- (\x+0.12,0.25) -- (\x+0.12,0);}
\node[popgreen!60!popdark, font=\footnotesize] at (2.5,-0.55) {récepteurs ACh activés};
% flèche réponse
\draw[-{Stealth}, line width=1.5pt, popgreen] (5.7,-0.8) -- (6.7,-0.8);
\node[popgreen!60!popdark, font=\footnotesize\bfseries, align=center] at (7.9,-0.8) {PA postsynaptique\\ contraction};
\end{scope}
% ===== Synapse curarisée (droite) =====
\begin{scope}[xshift=11.2cm]
\fill[popblueL] (-0.5,0.5) rectangle (5.5,3.6);
\node[popdark, font=\small\bfseries] at (2.5,4.0) {Synapse en présence de curare};
\node[popmuted, font=\footnotesize] at (2.5,3.35) {bouton présynaptique};
\foreach \x/\y in {0.6/2.6, 1.8/2.7, 3.0/2.5, 4.2/2.7}{
  \draw[popblue, fill=white] (\x,\y) circle (0.28);
  \foreach \a in {45,135,225,315}{
    \fill[poporange] ({\x+0.13*cos(\a)},{\y+0.13*sin(\a)}) circle (0.05);}}
\foreach \x/\y in {1.0/0.9, 2.2/1.1, 3.4/0.85, 4.4/1.15}{
  \fill[poporange] (\x,\y) circle (0.09);}
\node[popmuted, font=\footnotesize] at (2.5,0.35) {fente synaptique};
\fill[popgreenL] (-0.5,-1.6) rectangle (5.5,0.0);
\node[popmuted, font=\footnotesize] at (4.6,-1.35) {muscle};
% récepteurs bloqués par curare (losanges rouges)
\foreach \x in {1.0, 2.2, 3.4, 4.4}{
  \draw[popgreen, line width=1.2pt] (\x-0.12,0) -- (\x-0.12,0.25) -- (\x+0.12,0.25) -- (\x+0.12,0);
  \fill[poppink] (\x,0.42) circle (0.14);}
\node[poppink!70!popdark, font=\footnotesize] at (2.5,-0.55) {récepteurs bloqués par le curare};
\draw[-{Stealth}, line width=1.5pt, poppink] (5.7,-0.8) -- (6.7,-0.8);
\node[poppink!70!popdark, font=\footnotesize\bfseries, align=center] at (7.9,-0.8) {pas de PA\\ paralysie};
\end{scope}
% légende commune
\fill[poporange] (0.2,-2.3) circle (0.09); \node[right, font=\footnotesize] at (0.4,-2.3) {acétylcholine (ACh)};
\fill[poppink] (4.4,-2.3) circle (0.14); \node[right, font=\footnotesize] at (4.6,-2.3) {molécule de curare};
\end{tikzpicture}
\end{center}
\legende{Comparaison d'une synapse neuromusculaire normale (à gauche) et d'une synapse en présence de curare (à droite). L'acétylcholine est libérée normalement dans les deux cas, mais le curare occupe les récepteurs postsynaptiques sans les activer : aucun potentiel postsynaptique n'est généré, le muscle ne se contracte plus.}
\end{popfigure}

\courssub{L'atropine : un autre antagoniste de l'acétylcholine}

\begin{definition}[Atropine]
L'\cle{atropine} est un alcaloïde extrait de la belladone (\emph{Atropa belladonna}). C'est un antagoniste de l'acétylcholine agissant sur les récepteurs cholinergiques de type \emph{muscarinique}, présents notamment au niveau du cœur, des muscles lisses et des glandes (cibles du système nerveux parasympathique).
\end{definition}

Ses effets s'interprètent directement à partir du blocage des récepteurs :
\begin{itemize}
\item au niveau du \textbf{cœur} : le nerf vague (parasympathique) libère normalement de l'acétylcholine qui \emph{ralentit} le rythme cardiaque ; l'atropine bloquant ces récepteurs, le cœur échappe au frein vagal $\Rightarrow$ \cle{tachycardie} ;
\item au niveau de l'\textbf{iris} : blocage de la constriction pupillaire parasympathique $\Rightarrow$ \cle{mydriase} (pupille dilatée), propriété exploitée en ophtalmologie pour examiner le fond de l'œil ;
\item au niveau des \textbf{glandes} : diminution des sécrétions (sécheresse de la bouche).
\end{itemize}

\begin{attention}[Curare et atropine : deux antagonistes, deux cibles]
Curare et atropine bloquent tous deux l'acétylcholine, mais pas sur les mêmes récepteurs : le curare agit sur les récepteurs \emph{nicotiniques} de la jonction neuromusculaire (muscles squelettiques $\Rightarrow$ paralysie motrice), l'atropine sur les récepteurs \emph{muscariniques} des organes viscéraux (cœur, iris, glandes). Ne confonds pas leurs effets dans une copie.
\end{attention}

\courssub{Autres modes d'action (complément utile pour le Bac)}

\textbf{Les inhibiteurs de l'acétylcholinestérase.} De nombreux insecticides (organophosphorés, carbamates) et certains gaz de combat inhibent l'\cle{acétylcholinestérase}, l'enzyme qui hydrolyse l'acétylcholine dans la fente synaptique. Conséquence : l'acétylcholine s'accumule et stimule \emph{de façon prolongée et répétée} les récepteurs postsynaptiques $\Rightarrow$ contractions musculaires incontrôlées (fasciculations, crampes), hypersécrétion, troubles cardiaques, puis épuisement et paralysie. C'est l'effet \emph{inverse} du curare sur le plan du mécanisme (excitation prolongée au lieu d'un blocage), bien que l'issue puisse également être fatale.

\textbf{La toxine botulique.} Produite par \emph{Clostridium botulinum}, elle bloque la fusion des vésicules synaptiques avec la membrane présynaptique (clivage des protéines SNARE) $\Rightarrow$ \textbf{blocage de la libération de l'acétylcholine} (étape 1) $\Rightarrow$ paralysie flasque. Utilisée à très faible dose en médecine esthétique (Botox) et pour traiter certains spasmes musculaires.

\textbf{Les anesthésiants locaux.} La lidocaïne et ses analogues bloquent les \emph{canaux \ce{Na+} voltage-dépendants} des fibres nerveuses : la genèse et la propagation des potentiels d'action deviennent impossibles dans la zone traitée $\Rightarrow$ insensibilité locale (le dentiste peut intervenir sans douleur). Contrairement au curare, ces substances agissent sur la \emph{conduction} de l'influx, pas sur la synapse.

\begin{propriete}[Modes d'action des substances sur la transmission synaptique (admise)]
Une substance peut modifier la transmission synaptique de trois façons principales :
\begin{enumerate}
\item en \textbf{bloquant les récepteurs postsynaptiques} (ex. curare, atropine) $\Rightarrow$ transmission diminuée ou abolie ;
\item en \textbf{empêchant la libération du neurotransmetteur} présynaptique (ex. toxine botulique sur l'exocytose) $\Rightarrow$ transmission abolie ;
\item en \textbf{inhibant l'élimination du neurotransmetteur} de la fente (ex. inhibiteurs de l'acétylcholinestérase) $\Rightarrow$ transmission prolongée et excessive.
\end{enumerate}
Certaines substances agissent en amont de la synapse, en bloquant la genèse ou la propagation des potentiels d'action (anesthésiants locaux).
\end{propriete}

\begin{popfigure}
\begin{center}
\begin{tabularx}{\linewidth}{|l|X|X|X|}
\hline
\rowcolor{popblueL} \textbf{Substance} & \textbf{Site d'action} & \textbf{Mécanisme} & \textbf{Effet sur la transmission} \\
\hline
Curare & Récepteurs nicotiniques (jonction neuromusculaire) & Antagoniste compétitif de l'ACh & Bloquée $\Rightarrow$ paralysie musculaire \\
\hline
Atropine & Récepteurs muscariniques (cœur, iris, glandes) & Antagoniste de l'ACh & Bloquée $\Rightarrow$ tachycardie, mydriase \\
\hline
Insecticides organophosphorés & Acétylcholinestérase (fente synaptique) & Inhibition de l'hydrolyse de l'ACh & Prolongée, excessive $\Rightarrow$ contractions incontrôlées \\
\hline
Toxine botulique & Terminaison présynaptique & Blocage de l'exocytose de l'ACh & Abolie $\Rightarrow$ paralysie flasque \\
\hline
Anesthésiants locaux (lidocaïne) & Canaux \ce{Na+} de la fibre nerveuse & Blocage de la genèse/propagation du PA & Influx interrompu $\Rightarrow$ insensibilité locale \\
\hline
\end{tabularx}
\end{center}
\legende{Tableau de synthèse : quelques substances modifiant la transmission synaptique ou la conduction de l'influx nerveux, leur site d'action, leur mécanisme et leur effet.}
\end{popfigure}

\courssub{Méthode d'analyse et exercice d'application}

\begin{methode}[Interpréter l'effet d'une substance sur la transmission synaptique]
\begin{enumerate}
\item \textbf{Comparer} la réponse (contraction musculaire, potentiel postsynaptique enregistré) \emph{avant} et \emph{après} application de la substance, à stimulation identique.
\item \textbf{Localiser} l'étape perturbée grâce à des tests ciblés : le nerf conduit-il encore ? Le muscle stimulé directement répond-il ? Le neurotransmetteur est-il encore libéré ? S'accumule-t-il dans la fente ?
\item \textbf{Identifier} le mécanisme : blocage des récepteurs, blocage de la libération, inhibition de l'élimination, ou blocage de la conduction.
\item \textbf{Conclure} en reliant le mécanisme moléculaire à l'effet observé sur l'organisme entier.
\end{enumerate}
\end{methode}

\begin{exoresolu}[Analyse d'une expérience de curarisation]
\textbf{Énoncé.} Sur une préparation nerf--muscle de grenouille, on réalise trois séries de stimulations et on note la réponse du muscle :

\begin{center}
\begin{tabular}{|l|c|c|}
\hline
\rowcolor{popblueL} \textbf{Conditions} & \textbf{Stimulation du nerf} & \textbf{Stimulation directe du muscle} \\
\hline
Sans substance & contraction & contraction \\
\hline
Après application de la substance X sur la plaque motrice & pas de contraction & contraction \\
\hline
Après application de X + injection d'ACh sur le muscle & --- & pas de contraction \\
\hline
\end{tabular}
\end{center}

\begin{enumerate}
\item La substance X bloque-t-elle la conduction nerveuse ? Justifier.
\item Quel est le site d'action le plus probable de X ? Argumenter.
\item La substance X peut-elle être un inhibiteur de l'acétylcholinestérase ? Justifier.
\end{enumerate}

\tcblower
\textbf{Solution détaillée.}
\begin{enumerate}
\item \textbf{Non.} Si X bloquait la conduction le long du nerf, la stimulation directe du muscle resterait efficace, mais l'effet apparaît lorsque X est appliquée \emph{sur la plaque motrice} (zone synaptique). Or ici, le muscle stimulé directement répond encore : le muscle est intact. De plus, le nerf conduit toujours l'influx (sinon la stimulation nerveuse n'aurait aucun effet même sans X, ce qui n'est pas le cas au témoin). C'est donc un blocage de la \emph{transmission}, non de la conduction.
\item Le fait que l'injection d'acétylcholine sur le muscle traité ne provoque \emph{aucune} contraction montre que le muscle ne répond plus à son neurotransmetteur : les \textbf{récepteurs postsynaptiques de l'acétylcholine sont bloqués}. X est donc un antagoniste des récepteurs cholinergiques nicotiniques, de type curare.
\item \textbf{Non.} Un inhibiteur de l'acétylcholinestérase provoquerait une \emph{accumulation} d'acétylcholine dans la fente, donc une stimulation prolongée des récepteurs : on observerait des contractions répétées ou prolongées, et l'injection d'ACh resterait efficace (voire sur-efficace). Ici, la transmission est \emph{abolie} et l'ACh exogène est sans effet : c'est incompatible avec une inhibition de l'enzyme.
\end{enumerate}
\textbf{Conclusion :} X agit comme le curare, par blocage compétitif des récepteurs postsynaptiques de l'acétylcholine à la jonction neuromusculaire.
\end{exoresolu}

\courssub{Applications et enjeux de santé publique}

\begin{savaistu}[Insecticides et santé au Cameroun]
De nombreux insecticides agricoles utilisés au Cameroun (notamment dans les cacaoyères et les cultures maraîchères) appartiennent aux familles des \emph{organophosphorés} et des \emph{carbamates} : ce sont des inhibiteurs de l'acétylcholinestérase. Leur manipulation sans équipement de protection (gants, masque, combinaison) expose les agriculteurs à des intoxications aiguës (crampes, hypersalivation, troubles respiratoires) et chroniques. D'où l'importance capitale des précautions d'emploi, du respect des doses et des délais avant récolte. À l'inverse, la maîtrise de ces mécanismes a permis de développer des médicaments : myorelaxants dérivés du curare en chirurgie, atropine en ophtalmologie et comme antidote des intoxications aux organophosphorés.
\end{savaistu}

\begin{aretenir}[L'essentiel de la section]
\begin{itemize}
\item Une substance peut modifier la transmission synaptique en \textbf{bloquant les récepteurs} (curare, atropine), en \textbf{empêchant la libération} du neurotransmetteur (toxine botulique) ou en \textbf{inhibant son élimination} (insecticides organophosphorés).
\item Le \textbf{curare} bloque les récepteurs nicotiniques de l'acétylcholine à la jonction neuromusculaire $\Rightarrow$ paralysie ; la conduction nerveuse et le muscle restent fonctionnels (preuve : le muscle curarisé répond à la stimulation directe).
\item L'\textbf{atropine} bloque les récepteurs muscariniques $\Rightarrow$ tachycardie, mydriase.
\item Les \textbf{inhibiteurs de l'acétylcholinestérase} prolongent l'action de l'acétylcholine $\Rightarrow$ hyperexcitation synaptique.
\item La démarche d'interprétation repose sur la \textbf{comparaison avant/après} application et la \textbf{localisation de l'étape perturbée}.
\item Poisons, pesticides, drogues et médicaments agissent sur les mêmes cibles : seules la dose, la voie d'administration et la finalité diffèrent.
\end{itemize}
\end{aretenir}

\begin{exercice}[Pour t'entraîner]
On enregistre le potentiel postsynaptique d'une jonction neuromusculaire de grenouille dans trois conditions : (A) témoin ; (B) en présence d'une faible dose de curare ; (C) en présence d'un inhibiteur de l'acétylcholinestérase. On obtient respectivement : un potentiel d'amplitude normale et de courte durée ; un potentiel d'amplitude réduite ; un potentiel d'amplitude normale mais de durée très prolongée. Interprète chaque résultat en identifiant l'étape synaptique modifiée, puis prédis l'effet d'une forte dose de curare sur l'amplitude du potentiel postsynaptique.
\end{exercice}

\begin{corrige}[Exercice]
\begin{itemize}
\item \textbf{(A) Témoin :} libération normale d'ACh, fixation sur les récepteurs, hydrolyse rapide par l'acétylcholinestérase $\Rightarrow$ réponse brève.
\item \textbf{(B) Faible dose de curare :} une partie des récepteurs est bloquée ; moins de canaux postsynaptiques s'ouvrent $\Rightarrow$ amplitude réduite. L'étape modifiée est la \emph{fixation de l'ACh sur ses récepteurs}.
\item \textbf{(C) Inhibiteur de l'acétylcholinestérase :} l'ACh n'est plus dégradée, elle reste fixée plus longtemps sur les récepteurs $\Rightarrow$ durée prolongée. L'étape modifiée est l'\emph{élimination du neurotransmetteur}.
\item \textbf{Forte dose de curare :} la quasi-totalité des récepteurs est occupée $\Rightarrow$ le potentiel postsynaptique devient nul ou inférieur au seuil : plus de potentiel d'action postsynaptique, donc paralysie complète.
\end{itemize}
\end{corrige}

\newpage

\coursec{Activité d'intégration}

\coursec{Activité d'intégration : Le protocole anesthésique au bloc opératoire}

\courssub{Objectifs de l'activité}
Cette activité d'intégration vise à mobiliser l'ensemble des savoirs et savoir-faire de la leçon « Fonctionnement des neurones » pour résoudre une situation-problème authentique. Elle permet de :
\begin{itemize}
    \item Analyser et interpréter un enregistrement oscillographique de potentiels de repos et d'action (Document 1) : identifier les phases, mesurer les amplitudes, calculer la vitesse de conduction.
    \item Exploiter les résultats d'une expérience pharmacologique sur la transmission neuromusculaire (Document 2) : localiser le site d'action du curare.
    \item Construire une explication scientifique cohérente d'un protocole anesthésique réel (Document 3) : articuler le blocage de la genèse du potentiel d'action (anesthésique local) et le blocage de la transmission synaptique (myorelaxant).
    \item Rédiger une argumentation structurée en utilisant le vocabulaire précis de la neurophysiologie (canaux ioniques \emph{voltage-dépendants}, dépolarisation, seuil d'excitabilité, potentiel de plaque motrice, récepteurs nicotiniques).
\end{itemize}

\courssub{Situation : Pourquoi le patient ne ressent-il plus la douleur ?}
\textbf{Contexte clinique à l'Hôpital Général de Douala.}\par
Un patient de 35 ans, victime d'une fracture ouverte du tibia gauche lors d'un accident de la route sur l'autoroute Yaoundé-Douala, est pris en charge en urgence au bloc opératoire de l'Hôpital Général de Douala. L'équipe d'anesthésie-réanimation met en place une \textbf{anesthésie loco-régionale} (bloc du nerf sciatique par injection de \textbf{lidocaïne} à \SI{2}{\percent}) associée à une \textbf{curarisation} (injection intraveineuse d'\textbf{atracurium}, un myorelaxant non dépolarisant). L'intervention chirurgicale (ostéosynthèse par clou centromédullaire) dure \SI{90}{min}. Le patient reste conscient, ne signale aucune douleur, ne bouge pas le membre opéré, mais conserve ses réflexes pupillaires et sa ventilation spontanée (sous oxygénothérapie).

\textbf{Document 1 : Étude du potentiel d'action sur axone géant de calmar (\emph{Dosidicus gigas}).}\par
Un axone géant est isolé et placé dans une chambre de bain à température contrôlée (\SI{18}{\celsius}). Deux électrodes intracellulaires (E1 et E2) sont enfoncées dans l'axone, espacées d'une distance $d = \SI{5,0}{cm}$. Une électrode de référence extracellulaire est placée dans le bain. Une stimulation électrique brève (\SI{0,1}{ms}, intensité supraseuil) est appliquée au niveau de l'électrode E1. L'oscilloscope enregistre la différence de potentiel $V_m = V_{\text{intra}} - V_{\text{extra}}$ en fonction du temps pour chaque électrode.

\begin{popfigure}
\begin{tikzpicture}[scale=0.85, every node/.style={font=\small}]
    % Axes
    \draw[->] (-0.5,0) -- (13,0) node[right] {$t \text{ (ms)}$};
    \draw[->] (0,-1.2) -- (0,5) node[above] {$V_m \text{ (mV)}$};
    % Grid lines
    \foreach \y in {-70,-50,0,30} \draw[popline,very thin] (0,\y/20) -- (12.5,\y/20);
    \foreach \x in {0,1,2,3,4,5,6,7,8,9,10,11,12} \draw[popline,very thin] (\x,-1.2) -- (\x,4.8);
    % Labels y
    \node[left] at (0,-3.5) {$-70$};
    \node[left] at (0,0) {$0$};
    \node[left] at (0,1.5) {$+30$};
    % Electrode 1 trace (solid popblue)
    \draw[popblue, thick, domain=0:12, samples=500, smooth] plot (\x, {
        -3.5 + 
        (3.5+1.5)*exp(-(\x-1)^2/0.05) * (\x>0.9 && \x<1.1) + % stimulus artifact approx
        5.0*exp(-(\x-1.5)^2/0.15) * (\x>1.0 && \x<2.5) + % rising phase
        -4.0*exp(-(\x-2.5)^2/0.3) * (\x>2.0 && \x<3.5) + % falling phase
        -1.0*exp(-(\x-4)^2/1.0) * (\x>3.0) % hyperpol
    }) node[right, popblue] {E1};
    % Electrode 2 trace (dashed poporange) - delayed by 0.5 ms
    \draw[poporange, thick, dashed, domain=0:12, samples=500, smooth] plot (\x, {
        -3.5 + 
        5.0*exp(-(\x-2.0)^2/0.15) * (\x>1.5 && \x<3.0) +
        -4.0*exp(-(\x-3.0)^2/0.3) * (\x>2.5 && \x<4.0) +
        -1.0*exp(-(\x-4.5)^2/1.0) * (\x>3.5)
    }) node[right, poporange] {E2};
    % Annotations
    \draw[popmuted, <->] (1.5, 1.6) -- (2.0, 1.6) node[midway, above] {$\Delta t = \SI{0,5}{ms}$};
    \node[popblue] at (1.5, 2.5) {Pic E1 : $t_1 \approx \SI{1,5}{ms}$};
    \node[poporange] at (2.0, 3.2) {Pic E2 : $t_2 \approx \SI{2,0}{ms}$};
    % Resting potential line
    \draw[popmuted, dashed] (0,-3.5) -- (12.5,-3.5) node[right] {Potentiel de repos $V_{\text{repos}} = \SI{-70}{mV}$};
    % Threshold line
    \draw[popmuted, dashed] (0,-1.0) -- (12.5,-1.0) node[right] {Seuil $\approx \SI{-50}{mV}$};
\end{tikzpicture}
\legende{Enregistrement oscillographique simulé des potentiels d'action aux électrodes E1 (bleu continu) et E2 (orange pointillé) sur axone géant de calmar. Distance E1-E2 = \SI{5,0}{cm}. Déclenchement du stimulus à $t=0$. Amplitude du pic : $\approx \SI{+30}{mV}$. Potentiel de repos : $\SI{-70}{mV}$. Seuil d'excitabilité : $\approx \SI{-50}{mV}$.}
\end{popfigure}

\textbf{Document 2 : Effet du curare sur la préparation nerf-muscle de grenouille (\emph{Xenopus laevis}).}\par
On isole le nerf sciatique avec son muscle gastrocnémien attaché. La préparation est maintenue oxygénée dans une solution de Ringer à \SI{20}{\celsius}. On réalise trois tests successifs :
\begin{enumerate}
    \item \textbf{Témoin :} Stimulation électrique du nerf sciatique (\SI{0,5}{ms}, \SI{10}{V}) $\rightarrow$ contraction vigoureuse du gastrocnémien.
    \item \textbf{Après perfusion de curare (\SI{e-5}{mol/L}) pendant 10 min :} Stimulation électrique du nerf sciatique (mêmes paramètres) $\rightarrow$ \textbf{aucune contraction} du muscle.
    \item \textbf{Test de contrôle post-curare :} Stimulation électrique \emph{directe} du muscle gastrocnémien (électrodes posées sur le ventre musculaire, mêmes paramètres) $\rightarrow$ \textbf{contraction normale} du muscle.
\end{enumerate}

\textbf{Document 3 : Mécanismes d'action des agents anesthésiques utilisés.}\par
\begin{itemize}
    \item \textbf{Lidocaïne (anesthésique local) :} Molécule amphiphile qui se lie sur le site intracellulaire des \textbf{canaux $\ce{Na+}$ voltage-dépendants} à l'état inactivé ou ouvert, les bloquant de manière \emph{dépendante de l'utilisation} (\emph{use-dependent}). Elle n'affecte pas les canaux $\ce{K+}$ voltage-dépendants ni les canaux de fuite.
    \item \textbf{Atracurium (myorelaxant non dépolarisant) :} Antagoniste compétitif réversible des \textbf{récepteurs nicotiniques à l'acétylcholine (nAChR)} sur la plaque motrice (membrane post-synaptique de la jonction neuromusculaire). Il ne traverse pas la barrière hémato-encéphalique.
\end{itemize}

\courssub{Consignes}
\begin{exercice}[Questions de l'activité]
\begin{enumerate}
\item \textbf{Analyse du Document 1 (Potentiel d'action et conduction) :}
\begin{enumerate}
    \item Sur la courbe de l'électrode E1, identifier et nommer les quatre phases successives du potentiel d'action. Préciser les variations de perméabilité membranaire ($P_{\ce{Na+}}$, $P_{\ce{K+}}$) qui les sous-tendent.
    \item Calculer la \textbf{vitesse de conduction} de l'influx nerveux dans cet axone (en m/s). Commenter cette valeur au regard du diamètre de l'axone géant (environ \SIrange{500}{1000}{\micro\meter}) et de l'absence de myéline.
    \item Expliquer pourquoi le potentiel d'action enregistré en E2 est \emph{identique en forme et en amplitude} à celui enregistré en E1, malgré la distance de 5 cm.
\end{enumerate}

\item \textbf{Analyse du Document 2 (Transmission neuromusculaire et curare) :}
\begin{enumerate}
    \item À quelle étape précise de la transmission neuromusculaire le curare agit-il ? Justifier votre réponse en comparant les résultats des tests 1, 2 et 3.
    \item Préciser le type de récepteur ciblé et le mécanisme d'action (agoniste/antagoniste, compétitif/non compétitif).
    \item Pourquoi la stimulation directe du muscle (test 3) provoque-t-elle encore une contraction ?
\end{enumerate}

\item \textbf{Synthèse clinique (Document 3 + cours) :}
\begin{enumerate}
    \item Expliquer au niveau moléculaire et cellulaire pourquoi le patient \textbf{ne ressent pas la douleur} (blocage de la nociception) malgré l'incision cutanée et la manipulation osseuse.
    \item Expliquer pourquoi le patient \textbf{ne bouge pas le membre opéré} (blocage moteur) alors qu'il est conscient.
    \item Justifier le fait que la \textbf{ventilation spontanée} est conservée (indice : site d'action de l'atracurium vs centres respiratoires bulbo-pontins).
    \item Rédiger, en 15 lignes maximum, l'explication scientifique complète du protocole anesthésique utilisé, destinée à un interne en chirurgie orthopédique.
\end{enumerate}
\end{enumerate}
\end{exercice}

\courssub{Ressources à mobiliser}
\begin{aretenir}[Rappels de cours utiles]
\begin{itemize}
\item \textbf{Vitesse de conduction :} $v = \frac{d}{\Delta t}$ avec $d$ en m, $\Delta t$ en s $\rightarrow$ $v$ en m/s.
\item \textbf{Phases du potentiel d'action :} Repos ($P_{\ce{K+}} \gg P_{\ce{Na+}}$) $\rightarrow$ Dépolarisation (ouverture canaux $\ce{Na+}$ v-d, $P_{\ce{Na+}} \gg P_{\ce{K+}}$) $\rightarrow$ Repolarisation (inactivation canaux $\ce{Na+}$ + ouverture retardée canaux $\ce{K+}$ v-d, $P_{\ce{K+}} \gg P_{\ce{Na+}}$) $\rightarrow$ Hyperpolarisation (inertie fermeture canaux $\ce{K+}$).
\item \textbf{Propagation :} Courants locaux (circuits de décharge) entre zone active (dépolarisée) et zones adjacentes au repos $\rightarrow$ dépolarisation locale au seuil $\rightarrow$ régénération \emph{in toto} du PA (loi du tout ou rien).
\item \textbf{Synapse chimique (jonction neuromusculaire) :} Arrivée PA $\rightarrow$ ouverture canaux $\ce{Ca^2+}$ v-d (bouton) $\rightarrow$ exocytose vésicules ACh $\rightarrow$ fixation ACh sur récepteurs nAChR (canaux $\ce{Na+}/\ce{K+}$ ligand-dépendants) $\rightarrow$ dépolarisation post-synaptique (PPM) $\rightarrow$ si PPM > seuil $\rightarrow$ PA musculaire $\rightarrow$ contraction.
\item \textbf{Blocage canaux $\ce{Na+}$ (Lidocaïne) :} Empêche la phase de dépolarisation $\rightarrow$ pas de PA $\rightarrow$ pas de propagation de l'influx sensitif (douleur) ni moteur.
\item \textbf{Blocage récepteurs nAChR (Atracurium) :} Empêche l'ouverture des canaux post-synaptiques par l'ACh $\rightarrow$ pas de PPM $\rightarrow$ pas de PA musculaire $\rightarrow$ pas de contraction (paralysie flasque). Pas d'effet central (BHE).
\end{itemize}
\end{aretenir}

\begin{exoresolu}[Corrigé détaillé de l'activité d'intégration]
\tcblower
\textbf{1. Analyse du Document 1}
\begin{enumerate}
    \item \textbf{Phases du PA (électrode E1) :}
    \begin{itemize}
        \item \textbf{Phase 1 : Dépolarisation brutale (\SI{\approx 0,5}{ms} env.).} $V_m$ passe de $\SI{-70}{mV}$ à $\SI{+30}{mV}$. \textbf{Cause :} Ouverture massive et rapide des \textbf{canaux $\ce{Na+}$ voltage-dépendants} (activation $m^3$). $P_{\ce{Na+}}$ augmente $\sim 1000\times$ ; $V_m$ tend vers $E_{\ce{Na}} \approx \SI{+60}{mV}$ (équation de Nernst).
        \item \textbf{Phase 2 : Repolarisation (\SI{\approx 1,5}{ms} env.).} $V_m$ redescend vers $\SI{-70}{mV}$. \textbf{Cause :} Double mécanisme : (a) \emph{Inactivation} des canaux $\ce{Na+}$ (fermeture balle $h$) $\rightarrow P_{\ce{Na+}}$ chute ; (b) \emph{Ouverture retardée} des \textbf{canaux $\ce{K+}$ voltage-dépendants} (activation $n^4$) $\rightarrow P_{\ce{K+}}$ augmente fortement. $V_m$ tend vers $E_{\ce{K}} \approx \SI{-90}{mV}$.
        \item \textbf{Phase 3 : Hyperpolarisation (potentiel de suite, \SI{5-10}{ms}).} $V_m$ devient plus négatif que $V_{\text{repos}}$ (ex. $\SI{-80}{mV}$). \textbf{Cause :} Inertie de fermeture des canaux $\ce{K+}$ (retard de désactivation) $\rightarrow P_{\ce{K+}}$ reste temporairement $> P_{\ce{K+}}^{\text{repos}}$.
        \item \textbf{Phase 4 : Retour au repos.} Fermeture complète canaux $\ce{K+}$ v-d $\rightarrow$ restauration $P_{\ce{K+}}^{\text{repos}}$ et $P_{\ce{Na+}}^{\text{repos}}$ par la pompe $\ce{Na+/K+}$-ATPase (3 $\ce{Na+}$ sortent / 2 $\ce{K+}$ entrent / ATP).
    \end{itemize}
    \item \textbf{Vitesse de conduction :}
    \[
    v = \frac{d}{\Delta t} = \frac{\SI{5,0e-2}{m}}{\SI{0,5e-3}{s}} = \frac{0,05}{0,0005} = \SI{100}{m/s}.
    \]
    \textbf{Commentaire :} Cette vitesse (\SI{100}{m/s} = \SI{360}{km/h}) est exceptionnellement élevée pour une fibre \emph{amylinique}. Elle s'explique par le \textbf{diamètre géant} de l'axone (\SIrange{500}{1000}{\micro\meter} vs \SIrange{1}{20}{\micro\meter} pour les fibres mammifères). La résistance axoplasmique $R_a \propto 1/r^2$ est très faible, ce qui permet aux courants locaux de se propager loin (grande constante de longueur $\lambda$) et de dépolariser rapidement le segment adjacent au seuil. Chez l'homme, les fibres myélinisées (A$\alpha$, \SIrange{12}{20}{\micro\meter}) atteignent \SIrange{70}{120}{m/s} par conduction saltatoire ; les fibres C amyliniques (douleur diffuse) ne font que \SIrange{0,5}{2}{m/s}.
    \item \textbf{Identité des PA en E1 et E2 :} Le potentiel d'action est un phénomène \textbf{tout-ou-rien} et \textbf{régénératif}. En E2, ce n'est pas le signal d'E1 qui arrive atténué (conduction décrémentale), mais un \emph{nouveau} potentiel d'action déclenché \emph{in situ} par les courants locaux issus de la zone active (E1). La membrane en E2 possède la même densité de canaux $\ce{Na+}/\ce{K+}$ v-d et les mêmes gradients ioniques ; elle répond donc de manière identique (même seuil, même amplitude, même cinétique) dès que la dépolarisation locale atteint le seuil ($\approx \SI{-50}{mV}$).
\end{enumerate}

\textbf{2. Analyse du Document 2}
\begin{enumerate}
    \item \textbf{Site d'action du curare :} La \textbf{jonction neuromusculaire (synapse chimique motrice)}, plus précisément la \textbf{membrane post-synaptique (plaque motrice)}.
    \textbf{Justification :}
    \begin{itemize}
        \item Test 1 (Témoin) : Nerf stimulé $\rightarrow$ PA dans axone $\rightarrow$ libération ACh $\rightarrow$ PPM $\rightarrow$ PA muscle $\rightarrow$ Contraction. \emph{Tout fonctionne.}
        \item Test 2 (Curare + Stim. Nerf) : Nerf stimulé $\rightarrow$ PA dans axone \emph{normal} (le curare ne bloque pas les canaux $\ce{Na+}$ axoniques) $\rightarrow$ libération ACh \emph{normale} $\rightarrow$ \textbf{Pas de contraction}. Le blocage est \emph{après} la libération de l'émetteur.
        \item Test 3 (Curare + Stim. Muscle direct) : Stimulation directe $\rightarrow$ PA musculaire \emph{normal} (dépolarisation directe de la sarcolemme au-dessus du seuil) $\rightarrow$ Contraction \emph{normale}. L'excitabilité musculaire et l'appareil contractile sont intacts.
    \end{itemize}
    \textbf{Conclusion :} Le curare empêche la \textbf{transduction du signal chimique (ACh) en signal électrique (PPM/PA)} au niveau de la plaque motrice.
    \item \textbf{Récepteur et mécanisme :} \textbf{Récepteur nicotinique à l'acétylcholine (nAChR)}, canal ionique ligand-dépendant ($\ce{Na+}/\ce{K+}$). Le curare (et l'atracurium) est un \textbf{antagoniste compétitif réversible} : il se lie sur le site de fixation de l'ACh (sous-unités $\alpha$) sans ouvrir le canal, empêchant physiquement l'ACh de s'y fixer. L'augmentation de la concentration en ACh (ex. par inhibiteur de l'acétylcholinestérase comme la néostigmine) peut lever le blocage.
    \item \textbf{Stimulation directe efficace :} Elle court-circuite la synapse. Elle dépolarise directement la sarcolemme (membrane musculaire) au-delà du seuil d'excitabilité des canaux $\ce{Na+}$ voltage-dépendants \emph{musculaires}, déclenchant un PA musculaire normal qui se propage dans les tubules T et libère le $\ce{Ca^2+}$ du réticulum sarcoplasmique.
\end{enumerate}

\textbf{3. Synthèse clinique}
\begin{enumerate}
    \item \textbf{Abolition de la douleur (Analgésie) :} La lidocaïne injectée autour du nerf sciatique diffuse dans l'endoneure et atteint les axones des \textbf{fibres sensitives (nocicepteurs A$\delta$ et C)}. En bloquant les \textbf{canaux $\ce{Na+}$ voltage-dépendants} à l'état ouvert/inactivé, elle empêche la \textbf{génèse du potentiel d'action} au niveau des terminaisons nerveuses libres (récepteurs de la douleur) et sa \textbf{propagation} le long de l'axone vers la moelle épinière. L'influx nociceptif n'atteint pas les centres supérieurs (thalamus, cortex cingulaire) $\rightarrow$ \textbf{pas de perception consciente de la douleur}. Le patient est conscient car la lidocaïne n'atteint pas le SNC (barrière hémato-encéphalique, dose locale).
    \item \textbf{Immobilité du membre (Myorelaxation) :} L'atracurium, par voie IV, se distribue dans le compartiment extravasculaire et atteint les \textbf{jonctions neuromusculaires} des muscles squelettiques du membre inférieur (et de tout le corps). En bloquant \textbf{compétitivement les récepteurs nAChR} de la plaque motrice, il empêche l'ACh libérée par les motoneurones alpha (qui déchargent normalement, voire plus sous stimulation chirurgicale) de déclencher une \textbf{dépolarisation de la plaque motrice (PPM)}. Sans PPM, pas de PA musculaire, pas de relargage $\ce{Ca^2+}$, \textbf{pas de contraction} $\rightarrow$ \textbf{paralysie flasque}. Le chirurgien opère sur un membre totalement relâché.
    \item \textbf{Ventilation spontanée conservée :} L'atracurium est un \textbf{curarisant non dépolarisant quaternaire} (chargé positivement) qui \textbf{ne franchit pas la barrière hémato-encéphalique (BHE)}. Il n'agit donc \emph{pas} sur les centres respiratoires du tronc cérébral (bulbe, pont) qui commandent la ventilation automatique (\emph{drive} respiratoire central intact). \textbf{Précision rigoureuse :} L'atracurium paralyse \emph{tous} les muscles striés, y compris les muscles respiratoires (diaphragme, intercostaux). La conservation de la ventilation spontanée sous curare impose cliniquement un \textbf{monitorage du blocage neuromusculaire} (stimulation du nerf facial : TOF \emph{Train-Of-Four} à 1-2 réponses) pour maintenir une force respiratoire résiduelle suffisante, ou une ventilation assistée/contrôlée. L'énoncé simplifie la réalité clinique ; l'argument « pas d'effet central » reste valide pour expliquer l'absence de dépression centrale de la commande respiratoire.
    \item \textbf{Explication pour l'interne (15 lignes max) :}
    \begin{quote}
    \footnotesize
    « Le protocole associe deux blocs périphériques distincts. (1) \textbf{Bloc sensitif (Lidocaïne locale)} : La lidocaïne inhibe les canaux $\ce{Na+}$ voltage-dépendants des axones du nerf sciatique (fibres A$\delta$/C nociceptives et A$\alpha$ motrices). Elle empêche la dépolarisation au seuil et la propagation des potentiels d'action : l'influx douloureux n'atteint pas la moelle, l'influx moteur descendant ne quitte pas le plexus. (2) \textbf{Bloc moteur de sécurité (Atracurium IV)} : L'atracurium, antagoniste compétitif des récepteurs nAChR de la plaque motrice, empêche la transmission neuromusculaire résiduelle (échappement possible du bloc local ou réflexes rachidiens). Il garantit l'immobilité totale (paralysie flasque) sans effet sédatif (ne passe pas la BHE). L'association permet une chirurgie sans douleur ni mouvement sur patient conscient, avec récupération rapide post-opératoire (métabolisation plasmatique de l'atracurium par estérases de Hofmann, diffusion de la lidocaïne). »
    \end{quote}
\end{enumerate}
\end{exoresolu}

\courssub{Bilan de l'activité}
\begin{aretenir}[Compétences validées]
\begin{itemize}
\item \textbf{Modélisation biophysique :} Relier la courbe de potentiel d'action aux variations de perméabilité ionique ($P_{\ce{Na+}}$, $P_{\ce{K+}}$) et aux propriétés des canaux voltage-dépendants (activation/inactivation).
\item \textbf{Quantification :} Appliquer la relation $v = d/\Delta t$ sur un document oscillographique réel, manipuler les unités (cm $\rightarrow$ m, ms $\rightarrow$ s) et interpréter la valeur au regard de la structure de la fibre (diamètre, myélinisation).
\item \textbf{Raisonnement expérimental :} Utiliser la logique du « test de contournement » (stimulation nerveuse vs stimulation musculaire directe) pour localiser précisément le site d'action d'un pharmacologue (curare/atracurium) sur la chaîne de la transmission neuromusculaire.
\item \textbf{Transfert clinique :} Articuler les mécanismes cellulaires (blocage canaux $\ce{Na+}$ axoniques vs blocage récepteurs nAChR post-synaptiques) pour expliquer les effets thérapeutiques (analgésie, myorelaxation) et les effets secondaires/spécificités (conscience conservée, ventilation) d'un protocole anesthésique standard au Cameroun.
\item \textbf{Communication scientifique :} Rédiger une synthèse concise, structurée, utilisant le vocabulaire exact (antagoniste compétitif, \emph{use-dependent}, tout-ou-rien, courant local, plaque motrice) attendu au Baccalauréat.
\end{itemize}
\end{aretenir}

\begin{attention}[Erreurs fréquentes à éviter]
\begin{itemize}
\item Confondre vitesse de conduction (macroscopique, mesurable entre deux points) et vitesse de propagation du courant local (microscopique, entre segments adjacents).
\item Croire que le curare agit sur l'axone moteur (blocage canaux $\ce{Na+}$) ou sur la contractilité musculaire directe (blocage $\ce{Ca^2+}$ ou actine/myosine).
\item Penser que la lidocaïne « endort le cerveau » (effet central) : elle agit localement sur les axones périphériques.
\item Oublier les unités dans le calcul de vitesse (conversion cm $\rightarrow$ m, ms $\rightarrow$ s obligatoire).
\item Décrire les phases du PA sans citer les canaux ioniques responsables (canaux $\ce{Na+}$ v-d pour dépolarisation, canaux $\ce{K+}$ v-d pour repolarisation/hyperpolarisation).
\item Oublier que l'atracurium paralyse \emph{aussi} les muscles respiratoires : la ventilation spontanée n'est possible que si le blocage est partiel (monitorage TOF).
\end{itemize}
\end{attention}

\newpage

\coursec{Méthodes et exercices résolus}

\courssub{Fiches méthodes et exercices résolus}

\begin{methode}[Observer les centres nerveux d'un mammifère (dissection d'un cerveau de mouton ou de lapin)]
\begin{enumerate}
\item \textbf{Préparer} le matériel : cerveau frais ou fixé (formol), boîte de dissection, ciseaux, scalpel, pinces, gants, blouse, lunettes de protection.
\item \textbf{Observer la face dorsale} : repérer les deux hémisphères cérébraux séparés par la scissure interhémisphérique, les circonvolutions (gyri et sulci), le cervelet (face postérieure, aspect strié) et le tronc cérébral prolongé par la moelle épinière.
\item \textbf{Observer la face ventrale} : repérer les bulbes olfactifs, les nerfs optiques (II) et leur chiasma, l'hypophyse (souvent arrachée), le pont (Varole) et le bulbe rachidien.
\item \textbf{Réaliser une coupe sagittale médiane} au scalpel : distinguer la substance grise (cortex cérébral et cérébelleux, en périphérie) de la substance blanche (centre semi-ovale, axones myélinisés), le corps calleux (commissure interhémisphérique), le cervelet en « arbre de vie » (arbor vitae).
\item \textbf{Conclure} : la substance grise contient les corps cellulaires des neurones, les dendrites et les débuts d'axones non myélinisés ; la substance blanche contient les fibres nerveuses (axones) entourées de myéline. Toujours légender le schéma d'observation avec un titre, l'orientation (vue dorsale, coupe sagittale) et l'échelle ou le grossissement.
\end{enumerate}
\textbf{Réflexe Bac :} un schéma d'observation sans titre, sans orientation ni légendes perd systématiquement des points. Respecter le code couleur : substance grise (crayon gris), substance blanche (laissé blanc ou crayon bleu très clair).
\end{methode}

\begin{exoresolu}[Observation d'un cerveau de mouton]
Au laboratoire du lycée, on dispose d'un cerveau de mouton. En vue dorsale, on observe deux parties volumineuses plissées, séparées par un sillon profond ; en arrière, une structure striée plus petite.
\begin{enumerate}
\item Nommer les structures observées.
\item Sur une coupe sagittale, on distingue deux régions de couleurs différentes. Les identifier et expliquer cette différence.
\end{enumerate}
\tcblower
\textbf{Solution.}
\begin{enumerate}
\item Les deux parties volumineuses plissées sont les \textbf{hémisphères cérébraux} (leurs plis sont les circonvolutions), séparés par la \textbf{scissure interhémisphérique}. La structure postérieure striée est le \textbf{cervelet}.
\item Sur la coupe sagittale, on distingue :
\begin{itemize}
\item la \textbf{substance grise}, périphérique (cortex cérébral et cortex cérébelleux), formée des corps cellulaires des neurones, des dendrites et des fibres non myélinisées, d'où sa teinte grisâtre ;
\item la \textbf{substance blanche}, centrale, formée d'axones longs entourés de \textbf{myéline} (gaine lipidique blanchâtre), assurant la conduction rapide de l'influx entre aires corticales et vers la moelle.
\end{itemize}
\end{enumerate}
\textbf{Conclusion :} la couleur des régions du cerveau reflète l'organisation du tissu nerveux : corps cellulaires en périphérie (substance grise, centres d'intégration), fibres myélinisées au centre (substance blanche, voies de conduction).
\end{exoresolu}

\begin{methode}[Schématiser un neurone]
\begin{enumerate}
\item \textbf{Identifier} les trois parties obligatoires : le \cle{corps cellulaire} (pérykaryon, avec le noyau volumineux et le nucléole apparent), les \cle{dendrites} (courtes, ramifiées, nombreuses, réceptrices) et l'\cle{axone} (long, unique, émergeant au cône d'émergence, souvent myélinisé).
\item \textbf{Dessiner} : un corps cellulaire étoilé contenant un noyau central ; des dendrites arborescentes (arborisation dendritique) ; un axone long avec des segments de myéline (gaine de Schwann dans le SNP) séparés par les \textbf{nœuds de Ranvier}, se terminant par l'arborisation terminale (boutons synaptiques ou terminaisons boutonnantes).
\item \textbf{Légender} chaque structure avec des flèches fines (traits de rappel), sans croisement des traits de légende, vocabulaire précis : « Noyau », « Dendrites », « Cône d'émergence », « Gaine de myéline », « Nœud de Ranvier », « Boutons synaptiques ».
\item \textbf{Indiquer le sens de propagation} de l'influx nerveux : dendrites $\rightarrow$ corps cellulaire $\rightarrow$ cône d'émergence $\rightarrow$ axone $\rightarrow$ terminaisons synaptiques (flèches de couleur distincte).
\item \textbf{Titrer} le schéma : « Schéma d'un neurone moteur myélinisé du système nerveux périphérique ».
\end{enumerate}
\end{methode}

\begin{exoresolu}[Schéma annoté d'un neurone]
Réaliser le schéma légendé d'un neurone moteur myélinisé en indiquant le sens de propagation du message nerveux.
\tcblower
\textbf{Solution.} Le schéma attendu comporte : corps cellulaire avec noyau, dendrites, cône d'émergence, axone entouré de gaine de myéline interrompue aux nœuds de Ranvier, arborisation terminale à boutons synaptiques.

\begin{popfigure}
\begin{center}
\begin{tikzpicture}[scale=0.9, every node/.style={font=\small}]
% Corps cellulaire
\draw[fill=poporangeL, thick] (0,0) circle (0.7);
\draw[fill=popink!40] (0,0) circle (0.25);
% Dendrites
\draw[thick, popdark] (-0.5,0.5) -- (-1.4,1.3); \draw[thick, popdark] (-1.4,1.3) -- (-1.9,1.5); \draw[thick, popdark] (-1.4,1.3) -- (-1.5,1.9);
\draw[thick, popdark] (0.2,0.68) -- (0.4,1.6); \draw[thick, popdark] (0.4,1.6) -- (0.1,2.2); \draw[thick, popdark] (0.4,1.6) -- (0.9,2.1);
\draw[thick, popdark] (-0.65,-0.2) -- (-1.6,-0.6); \draw[thick, popdark] (-1.6,-0.6) -- (-2.2,-0.4);
% Cône d'émergence + Axone initial
\draw[thick, popdark] (0.7,-0.1) -- (1.6,-0.2);
% Myéline (gaine de Schwann)
\foreach \x in {1.6,2.8,4.0,5.2}{
    \draw[fill=popblueL, thick] (\x,-0.45) rectangle (\x+0.9,0.05);
    % Nœud de Ranvier (espace)
}
\draw[thick, popdark] (6.1,-0.2) -- (7.0,-0.3);
% Arborisation terminale
\draw[thick, popdark] (7.0,-0.3) -- (7.7,0.3); \draw[thick, popdark] (7.0,-0.3) -- (7.8,-0.3); \draw[thick, popdark] (7.0,-0.3) -- (7.7,-0.9);
\draw[fill=poppinkL, thick] (7.7,0.3) circle (0.12); \draw[fill=poppinkL, thick] (7.8,-0.3) circle (0.12); \draw[fill=poppinkL, thick] (7.7,-0.9) circle (0.12);
% Sens de propagation (flèches vertes)
\draw[-{Stealth}, very thick, popgreen] (-2.3,2.3) -- (-1.2,1.6);
\draw[-{Stealth}, very thick, popgreen] (0.8,0.5) -- (1.4,0.5);
\draw[-{Stealth}, very thick, popgreen] (3.0,0.6) -- (4.0,0.6);
\draw[-{Stealth}, very thick, popgreen] (6.3,0.4) -- (7.2,0.4);
% Légendes
\node[above] at (0,2.4) {\textbf{1. Dendrites}};
\node[left] at (-0.9,-0.9) {\textbf{2. Corps cellulaire} (noyau)};
\node[below] at (1.2,-0.8) {\textbf{3. Cône d'émergence}};
\node[below] at (3.2,-0.7) {\textbf{4. Gaine de myéline} (axone)};
\node[below] at (4.5,-0.7) {\textbf{5. Nœud de Ranvier}};
\node[below] at (7.6,-1.2) {\textbf{6. Boutons synaptiques}};
\end{tikzpicture}
\end{center}
\legende{Schéma d'un neurone moteur myélinisé (SNP). Les flèches vertes indiquent le sens de propagation de l'influx nerveux : dendrites $\rightarrow$ corps cellulaire $\rightarrow$ cône d'émergence $\rightarrow$ axone (conduction saltatoire) $\rightarrow$ terminaisons synaptiques.}
\end{popfigure}

\textbf{Conclusion :} le neurone est une cellule polarisée structuralement et fonctionnellement : il reçoit l'information par ses dendrites (zone réceptrice), l'intègre au cône d'émergence (zone de déclenchement) et la transmet par son axone (zone de conduction) vers les effecteurs.
\end{exoresolu}

\begin{definition}[Types de synapses]
On distingue deux grands types de synapses :
\begin{itemize}
\item \textbf{Synapse chimique} (majoritaire chez les vertébrés) : transmission par un médiateur chimique (neurotransmetteur) libéré dans une fente synaptique ($\approx 20-30$ nm). Unidirectionnelle, plasticité forte (sommation, modulation).
\item \textbf{Synapse électrique} (jonctions communicantes / \emph{gap junctions}) : continuité cytoplasmique entre pré- et postsynaptique via des connexons. Bidirectionnelle, transmission quasi instantanée (pas de délai synaptique), synchronisation (ex. : rythme cardiaque, réflexes d'échappement).
\end{itemize}
Le programme de Terminale porte essentiellement sur la \textbf{synapse à transmission chimique}.
\end{definition}

\begin{methode}[Décrire le protocole de mesure des potentiels de repos et d'action]
\begin{enumerate}
\item \textbf{Matériel} : une fibre nerveuse isolée de grand diamètre (ex. axone géant de calmar \emph{Loligo} ou nerf sciatique de grenouille), deux microélectrodes en verre tiré (pointe $< 1 \mu m$) remplies de solution KCl \SI{3}{M}, un oscilloscope (ou voltmètre à très haute impédance $> 10^9 \Omega$), un stimulateur électrique, une chambre de perfusion (sérum physiologique).
\item \textbf{Mesure du potentiel de repos} : piquer la membrane avec une microélectrode \emph{intracellulaire} (pointe à l'intérieur de l'axoplasme), la seconde électrode \emph{de référence} dans le milieu extracellulaire. L'oscilloscope indique une différence de potentiel stable : $U_r = -\SI{70}{mV}$ (intérieur négatif par rapport à l'extérieur). C'est le \cle{potentiel de repos}.
\item \textbf{Mesure du potentiel d'action} : stimuler la fibre par un choc électrique bref d'intensité croissante via des électrodes de stimulation.
\begin{itemize}
\item En dessous du \cle{seuil de dépolarisation} ($\approx -\SI{55}{mV}$) : réponse locale graduée (potentiel électrotonique), pas de propagation.
\item Au seuil et au-delà (stimulus liminaire ou supra-liminaire) : apparition d'un \cle{potentiel d'action} (PA) d'\emph{amplitude et de forme constantes}, indépendant de l'intensité du stimulus : c'est la \cle{loi du tout ou rien}.
\end{itemize}
\item \textbf{Tracer} la courbe $U_m = f(t)$ (différence de potentiel membranaire en fonction du temps) et analyser ses phases successives.
\end{enumerate}
\textbf{À retenir :} Potentiel de repos $U_r \approx -\SI{70}{mV}$ ; Seuil $U_s \approx -\SI{55}{mV}$ ; Sommet du PA $U_{max} \approx +\SI{30}{mV}$ ; Amplitude du PA $\approx \SI{100}{mV}$ ; Durée du PA $\approx \SI{1}{ms}$ à \SI{2}{ms} (fibres myélinisées mammifères).
\end{methode}

\begin{exoresolu}[Mise en évidence de la loi du tout ou rien]
On stimule un axone par des chocs électriques d'intensités croissantes et on enregistre le potentiel de membrane. Résultats :

\begin{center}
\begin{tabularx}{\linewidth}{|l|X|X|X|X|}\hline
\rowcolor{popblueL} Intensité du stimulus (u.a.) & 1 & 2 & 3 & 5 \\ \hline
Réponse enregistrée & Potentiel de repos ($-\SI{70}{mV}$) & Potentiel de repos ($-\SI{70}{mV}$) & Potentiel d'action ($+\SI{30}{mV}$) & Potentiel d'action ($+\SI{30}{mV}$) \\ \hline
\end{tabularx}
\end{center}

Interpréter ces résultats.
\tcblower
\textbf{Solution.}
\begin{itemize}
\item Pour des stimuli de 1 et 2 u.a. (stimuli \textbf{infra-liminaires}, intensité $<$ seuil), la membrane reste à son potentiel de repos de $-\SI{70}{mV}$ : \textbf{aucun potentiel d'action} n'est déclenché. Seule une dépolarisation locale non propagée (potentiel électrotonique) peut exister, non visible à l'échelle de l'oscilloscope si sous-seuil.
\item À partir de 3 u.a. (stimulus \textbf{liminaire} ou supra-liminaire, intensité $\ge$ seuil), un potentiel d'action apparaît, atteignant un sommet à $+\SI{30}{mV}$.
\item À 5 u.a., l'amplitude du PA reste \textbf{identique} ($+\SI{30}{mV}$) : elle ne dépend pas de l'intensité du stimulus.
\end{itemize}
\textbf{Interprétation :} le potentiel d'action obéit à la \textbf{loi du tout ou rien} : soit le stimulus n'atteint pas le seuil critique de dépolarisation (environ $-\SI{55}{mV}$) et aucun PA n'est émis (réponse « rien »), soit il l'atteint ou le dépasse et le PA se déclenche avec toujours la même amplitude maximale (réponse « tout »). L'intensité du stimulus est codée non par l'amplitude du PA, mais par la \textbf{fréquence} des PA (codage fréquentiel).

\textbf{Conclusion :} le message nerveux est un signal stéréotypé de type « tout ou rien », déclenché uniquement si la dépolarisation initiale atteint le seuil d'excitabilité.
\end{exoresolu}

\begin{methode}[Représenter et analyser la courbe de potentiel d'action]
\begin{enumerate}
\item \textbf{Axes} : abscisse = temps $t$ (en \textbf{ms}, millisecondes) ; ordonnée = différence de potentiel membranaire $U_m$ (en \textbf{mV}, millivolts). Toujours graduer et légender les axes avec les grandeurs et leurs unités.
\item \textbf{Phases à identifier sur la courbe} (valeurs typiques fibre myélinisée mammifère) :
\begin{itemize}
\item \textbf{Phase 0 : Repos} : $U_m = -\SI{70}{mV}$ (stable).
\item \textbf{Phase 1 : Dépolarisation} : montée brutale de $-\SI{55}{mV}$ (seuil) à $+\SI{30}{mV}$ (sommet) en $\approx \SI{0,5}{ms}$. Pente maximale $dU/dt$ maximale.
\item \textbf{Phase 2 : Repolarisation} : descente rapide de $+\SI{30}{mV}$ vers $-\SI{70}{mV}$.
\item \textbf{Phase 3 : Hyperpolarisation (potentiel de suite)} : passage transitoire sous $U_r$ (ex. $-\SI{80}{mV}$ à $-\SI{85}{mV}$).
\item \textbf{Phase 4 : Retour au repos} : rétablissement progressif de $U_r = -\SI{70}{mV}$.
\end{itemize}
\item \textbf{Relier chaque phase aux mouvements ioniques transversaux} (c'est l'interprétation exigée au Bac, pas seulement la description morphologique) :
\begin{itemize}
\item Dépolarisation $\rightarrow$ Ouverture canaux \ce{Na+} voltage-dépendants $\rightarrow$ Entrée massive \ce{Na+} (courant entrant $I_{Na}$).
\item Repolarisation $\rightarrow$ Fermeture canaux \ce{Na+} (inactivation) + Ouverture canaux \ce{K+} voltage-dépendants (lents) $\rightarrow$ Sortie massive \ce{K+} (courant sortant $I_K$).
\item Hyperpolarisation $\rightarrow$ Fermeture lente canaux \ce{K+} $\rightarrow$ Sortie résiduelle \ce{K+}.
\item Retour au repos $\rightarrow$ Pompe \ce{Na+}/\ce{K+}-ATPase (3 \ce{Na+}$_{out}$ / 2 \ce{K+}$_{in}$ par ATP) $\rightarrow$ Rétablissement gradients de concentration.
\end{itemize}
\item \textbf{Périodes de réfractarité} : \emph{Absolue} (pendant PA, canaux \ce{Na+} inactivés, impossible de déclencher 2ème PA) ; \emph{Relative} (pendant hyperpolarisation, seuil plus haut, stimulus plus fort nécessaire).
\end{enumerate}
\end{methode}

\begin{exoresolu}[Analyse d'un enregistrement de potentiel d'action]
L'enregistrement du potentiel membranaire d'un axone stimulé donne la courbe suivante : repos à $-\SI{70}{mV}$, montée rapide jusqu'à $+\SI{30}{mV}$ en \SI{0,5}{ms}, descente franchissant $-\SI{70}{mV}$ jusqu'à $-\SI{80}{mV}$, puis retour à $-\SI{70}{mV}$.
\begin{enumerate}
\item Tracer l'allure de la courbe.
\item Nommer les phases et expliquer les mouvements ioniques responsables.
\end{enumerate}
\tcblower
\textbf{Solution.}
\begin{enumerate}
\item Courbe du potentiel d'action (représentation schématique) :

\begin{popfigure}
\begin{center}
\begin{tikzpicture}
\begin{axis}[width=0.9\linewidth, height=6cm,
xlabel={Temps $t$ (ms)}, ylabel={$U_m$ (mV)},
xmin=0, xmax=7, ymin=-90, ymax=50,
xtick={0,1,2,3,4,5,6,7}, ytick={-80,-70,-55,0,30},
grid=both, grid style={popline!40},
axis lines=middle, axis line style={-Stealth},
xlabel style={anchor=north west}, ylabel style={anchor=south east}]
% Repos
\addplot[popcurve, domain=0:1, samples=10] {-70};
% Dépolarisation (montée exponentielle rapide)
\addplot[popcurve, domain=1:1.5, samples=50] {-70 + 100*(1 - exp(-15*(x-1)))};
% Repolarisation (descente exponentielle)
\addplot[popcurve, domain=1.5:3, samples=80] {30 - 110*(1 - exp(-4*(x-1.5)))};
% Hyperpolarisation + retour (lent)
\addplot[popcurve, domain=3:7, samples=100] {-80 + 10*(1 - exp(-0.8*(x-3)))};
% Seuil
\draw[dashed, popmuted, thick] (axis cs:0,-55) -- (axis cs:7,-55) node[pos=0.95, above, font=\scriptsize, popdark]{Seuil ($\approx -55$ mV)};
% Repos
\draw[dashed, popmuted, thick] (axis cs:0,-70) -- (axis cs:7,-70) node[pos=0.05, left, font=\scriptsize, popdark]{Repos ($-70$ mV)};
% Annotations phases
\node[font=\scriptsize, popdark, align=center] at (axis cs:1.25,10) {Dépolarisation\\ (Entrée \ce{Na+})};
\node[font=\scriptsize, popdark, align=center] at (axis cs:2.2,-20) {Repolarisation\\ (Sortie \ce{K+})};
\node[font=\scriptsize, popdark, align=center] at (axis cs:4.5,-85) {Hyperpolarisation\\ (Sortie \ce{K+} exc.)};
\node[font=\scriptsize, popdark, align=center] at (axis cs:6,-75) {Retour au repos\\ (Pompe \ce{Na+}/\ce{K+})};
\end{axis}
\end{tikzpicture}
\end{center}
\legende{Courbe du potentiel d'action : variations de la différence de potentiel membranaire $U_m$ (mV) en fonction du temps $t$ (ms). Les phases sont corrélées aux perméabilités ioniques membranaires.}
\end{popfigure}

\item \textbf{Interprétation ionique détaillée :}
\begin{itemize}
\item \textbf{Dépolarisation} (de $-\SI{55}{mV}$ à $+\SI{30}{mV}$) : ouverture brutale des canaux \ce{Na+} voltage-dépendants (activation), \textbf{entrée massive de \ce{Na+}} dans l'axone suivant son gradient électrochimique ($E_{Na} \approx +\SI{60}{mV}$), l'intérieur devient momentanément positif. Courant entrant $I_{Na}$.
\item \textbf{Repolarisation} : inactivation des canaux \ce{Na+} (fermeture « balle et chaîne ») et ouverture retardée des canaux \ce{K+} voltage-dépendants, \textbf{sortie de \ce{K+}} suivant son gradient ($E_K \approx -\SI{90}{mV}$), l'intérieur redevient négatif. Courant sortant $I_K$.
\item \textbf{Hyperpolarisation} ($-\SI{80}{mV}$) : les canaux \ce{K+} se ferment lentement (désactivation), la sortie de \ce{K+} se poursuit transitoirement au-delà du potentiel d'équilibre du \ce{K+}, l'intérieur devient plus négatif qu'au repos.
\item \textbf{Retour au repos} : la \textbf{pompe \ce{Na+}/\ce{K+}-ATPase} (consommant 1 ATP pour 3 \ce{Na+} sortis / 2 \ce{K+} entrés) rétablit progressivement les gradients de concentration ioniques initiaux (\ce{Na+}$_{in}$ bas, \ce{K+}$_{in}$ haut).
\end{itemize}
\end{enumerate}
\textbf{Conclusion :} le potentiel d'action résulte de variations rapides, séquentielles et réversibles de la perméabilité membranaire aux ions \ce{Na+} et \ce{K+} (canaux voltage-dépendants), la pompe \ce{Na+}/\ce{K+} assurant seulement la restauration à long terme des gradients.
\end{exoresolu}

\begin{methode}[Calculer et comparer la vitesse de conduction de l'influx nerveux]
\begin{enumerate}
\item \textbf{Formule} : $v = \dfrac{d}{\Delta t}$, avec $d$ la distance entre les deux points d'enregistrement (en \textbf{m}) et $\Delta t$ le décalage temporel entre les deux potentiels d'action (en \textbf{s}).
\item \textbf{Convertir} systématiquement les unités avant le calcul :
\begin{itemize}
\item Distance : cm $\rightarrow$ m (diviser par 100).
\item Temps : ms $\rightarrow$ s (diviser par 1000).
\end{itemize}
Rappel : \SI{1}{cm.ms^{-1}} = \SI{10}{m.s^{-1}}.
\item \textbf{Comparer} les types de fibres :
\begin{itemize}
\item Fibre \textbf{myélinisée} (conduction \emph{saltatoire} de nœud de Ranvier en nœud) : $v = \SI{10}{}$ à \SI{120}{\meter\per\second} (proportionnelle au diamètre).
\item Fibre \textbf{amyélinique} (conduction \emph{continue} de proche en proche) : $v = \SI{0,5}{}$ à \SI{2}{\meter\per\second} (proportionnelle à $\sqrt{\text{diamètre}}$).
\end{itemize}
\item \textbf{Conclure} : la myélinisation (isolation électrique) et un grand diamètre de fibre (faible résistance axoplasmique) augmentent drastiquement la vitesse de conduction.
\end{enumerate}
\end{methode}

\begin{exoresolu}[Vitesse de conduction sur deux fibres]
On enregistre le potentiel d'action en deux points A et B d'une fibre nerveuse, distants de $d = \SI{10}{cm}$. Le PA apparaît en B \SI{2}{ms} après son apparition en A (fibre 1, myélinisée). Sur une fibre 2 de même diamètre mais dépourvue de myéline, le décalage est de \SI{50}{ms}.
\begin{enumerate}
\item Calculer la vitesse de conduction dans chaque fibre.
\item Expliquer la différence observée.
\end{enumerate}
\tcblower
\textbf{Solution.}
\begin{enumerate}
\item \textbf{Fibre 1 (myélinisée) :}
Conversion : $d = \SI{10}{cm} = \SI{0,10}{m}$ ; $\Delta t_1 = \SI{2}{ms} = \SI{2e-3}{s} = \SI{0,002}{s}$.
\[ v_1 = \frac{d}{\Delta t_1} = \frac{\SI{0,10}{m}}{\SI{0,002}{s}} = \SI{50}{m.s^{-1}}. \]
\item \textbf{Fibre 2 (amyélinique) :}
$\Delta t_2 = \SI{50}{ms} = \SI{0,050}{s}$.
\[ v_2 = \frac{d}{\Delta t_2} = \frac{\SI{0,10}{m}}{\SI{0,050}{s}} = \SI{2}{m.s^{-1}}. \]
\end{enumerate}
\textbf{Interprétation :} la fibre 1 conduit l'influx 25 fois plus vite que la fibre 2 ($v_1/v_2 = 25$). La gaine de \textbf{myéline} (formée par les cellules de Schwann dans le SNP) isole électriquement l'axone : les échanges ioniques générateurs de courant transmembrane (et donc la régénération du PA) n'ont lieu qu'aux \textbf{nœuds de Ranvier} (dépourvus de myéline, riches en canaux \ce{Na+}/\ce{K+}). Le PA « saute » de nœud en nœud : c'est la \textbf{conduction saltatoire}, beaucoup plus rapide et économe en énergie (moins d'ions à pomper) que la conduction continue de proche en proche le long d'une fibre amyélinique où la membrane toute entière se dépolarise.

\textbf{Conclusion :} la myélinisation multiplie la vitesse de conduction de l'influx nerveux (ici de \SI{2}{m.s^{-1}} à \SI{50}{m.s^{-1}}), ce qui explique la gravité des maladies démyélinisantes (ex. sclérose en plaques) où la conduction est ralentie ou bloquée.
\end{exoresolu}

\begin{methode}[Schématiser une synapse à transmission chimique et expliquer son fonctionnement]
\begin{enumerate}
\item \textbf{Structures à dessiner et légender} (synapse type, ex. jonction neuromusculaire) :
\begin{itemize}
\item Élément \textbf{présynaptique} (bouton terminal) : vésicules synaptiques (neurotransmetteur), mitochondries (ATP pour recyclage), canaux \ce{Ca^{2+}} voltage-dépendants, membrane présynaptique (bouton).
\item \textbf{Fente synaptique} (espace extracellulaire $\approx \SI{20}{nm}$), matrice extracellulaire (acétylcholinestérase fixée si cholinergique).
\item Élément \textbf{postsynaptique} (membrane effectrice : plaque motrice musculaire, dendrite, corps cellulaire) : récepteurs spécifiques (canaux ioniques ligand-dépendants), canaux ioniques postsynaptiques.
\end{itemize}
\item \textbf{Étapes du fonctionnement à rédiger dans l'ordre chronologique} :
\begin{enumerate}
\item Arrivée du PA au bouton présynaptique $\rightarrow$ dépolarisation de la membrane présynaptique $\rightarrow$ ouverture des canaux \ce{Ca^{2+}} voltage-dépendants $\rightarrow$ \textbf{entrée massive de \ce{Ca^{2+}}} (gradient électrochimique favorable).
\item L'augmentation de \ce{[Ca^{2+}]}$_{in}$ déclenche la migration et la fusion des vésicules synaptiques avec la membrane présynaptique : \textbf{exocytose} du neurotransmetteur (ex. acétylcholine ACh) dans la fente synaptique.
\item \textbf{Diffusion} du neurotransmetteur à travers la fente et \textbf{fixation} sur les récepteurs spécifiques de la membrane postsynaptique (récepteurs nicotiniques pour l'ACh).
\item Ouverture des canaux ioniques postsynaptiques associés aux récepteurs $\rightarrow$ flux ioniques (\ce{Na+}$_{in}$, \ce{K+}$_{out}$ pour ACh) $\rightarrow$ variation du potentiel membranaire postsynaptique : \textbf{Potentiel PostSynaptique Excitateur (PPSE)} dépolarisant ou \textbf{Potentiel PostSynaptique Inhibiteur (PPSI)} hyperpolarisant.
\item \textbf{Inactivation rapide} du neurotransmetteur dans la fente : hydrolyse enzymatique (ex. acétylcholinestérase pour l'ACh) ou recapture par transporteurs (ex. sérotonine, dopamine, noradrénaline) $\rightarrow$ fin de la stimulation, brièveté du signal.
\end{enumerate}
\item \textbf{Propriétés fonctionnelles à citer} :
\begin{itemize}
\item Transmission \emph{unidirectionnelle} (pré $\rightarrow$ post uniquement).
\item \emph{Retard synaptique} ($\approx \SI{0,5}{ms}$ à \SI{1}{ms}, temps de la libération et diffusion).
\item \emph{Sommation} spatiale et temporelle des potentiels postsynaptiques (intégration neuronale).
\item \emph{Fatigabilité} (épuisement des vésicules lors de stimulations répétées à haute fréquence).
\item \emph{Plasticité} (modulation de l'efficacité synaptique, base de l'apprentissage/mémoire).
\end{itemize}
\end{enumerate}
\end{methode}

\begin{exoresolu}[Fonctionnement de la synapse neuromusculaire]
Chez un patient, on bloque expérimentalement les canaux calciques des terminaisons nerveuses motrices : le muscle ne se contracte plus malgré la stimulation du nerf moteur.
\begin{enumerate}
\item Schématiser la synapse concernée.
\item Expliquer pourquoi le blocage des canaux \ce{Ca^{2+}} supprime la contraction musculaire.
\end{enumerate}
\tcblower
\textbf{Solution.}
\begin{enumerate}
\item Il s'agit de la \textbf{jonction neuromusculaire} (synapse cholinergique), synapse chimique entre un motoneurone alpha et une fibre musculaire squelettique, dont le neurotransmetteur est l'\textbf{acétylcholine (ACh)}. Le schéma comporte :
\begin{itemize}
\item Côté pré : bouton terminal (vésicules d'ACh, mitochondries, canaux \ce{Ca^{2+}}), membrane présynaptique.
\item Fente synaptique (acétylcholinestérase).
\item Côté post : membrane postsynaptique (plaque motrice) avec replis de junction (augmentation surface), récepteurs nicotiniques à l'ACh (canaux \ce{Na+}/\ce{K+} ligand-dépendants).
\end{itemize}

\begin{popfigure}
\begin{center}
\begin{tikzpicture}[scale=1.1, every node/.style={font=\small}]
% Bouton présynaptique
\draw[fill=poporangeL!50, thick] (0,0) rectangle (3,1.5);
\foreach \x/\y in {0.4,0.3/ 0.8,0.5/ 1.2,0.8/ 1.8,0.3/ 2.2,1.0/ 2.6,0.6}{
    \draw[fill=popblue!60] (\x,\y) circle (0.1);
}
\draw[fill=popgreenL!50] (0.5,1.1) ellipse (0.3 and 0.15); % Mitochondrie
\node[font=\tiny] at (0.5,1.1) {Mito};
% Canaux Ca2+ présynaptiques
\foreach \x in {0.5, 1.5, 2.5}{
    \draw[thick, poppurple] (\x,0) -- (\x,-0.15);
    \draw[fill=poppurple] (\x,-0.15) rectangle (\x+0.15,0);
    \node[below, font=\tiny, poppurple] at (\x+0.075,-0.25) {Ca\textsuperscript{2+}};
}
% Fente synaptique
\draw[fill=popcream] (0,-0.5) rectangle (3,0);
\node[font=\scriptsize, popdark] at (1.5,-0.25) {Fente synaptique (\SI{20}{nm})};
\node[font=\tiny, popmuted] at (1.5,-0.4) {Acétylcholinestérase};
% Flèches diffusion ACh
\draw[-{Stealth}, poporange, thick] (1.5,0) -- (1.5,-0.5);
\draw[-{Stealth}, poporange, thick] (1.0,0) -- (1.0,-0.5);
\draw[-{Stealth}, poporange, thick] (2.0,0) -- (2.0,-0.5);
% Membrane postsynaptique (plaque motrice repliée)
\draw[thick, popdark] (0,-0.5) -- (0,-1.2) -- (0.3,-0.5) -- (0.6,-1.2) -- (0.9,-0.5) -- (1.2,-1.2) -- (1.5,-0.5) -- (1.8,-1.2) -- (2.1,-0.5) -- (2.4,-1.2) -- (2.7,-0.5) -- (3,-1.2) -- (3,-0.5);
% Récepteurs ACh
\foreach \x in {0.15, 0.45, 0.75, 1.05, 1.35, 1.65, 1.95, 2.25, 2.55, 2.85}{
    \draw[fill=poppink, thick] (\x,-0.6) rectangle (\x+0.15,-0.8);
    \node[font=\tiny, rotate=90] at (\x+0.075,-0.7) {Récepteur ACh};
}
% Légendes
\node[above] at (1.5,1.7) {\textbf{Bouton présynaptique (Motoneurone)}};
\node[below] at (1.5,-1.5) {\textbf{Membrane postsynaptique (Plaque motrice musculaire)}};
\node[left] at (-0.5,0.75) {Vésicules ACh};
\node[left] at (-0.5,0.2) {Canaux Ca\textsuperscript{2+}};
\end{tikzpicture}
\end{center}
\legende{Schéma de la jonction neuromusculaire (synapse cholinergique). 1: Vésicules d'acétylcholine. 2: Mitochondries. 3: Canaux Ca\textsuperscript{2+} voltage-dépendants. 4: Libération d'ACh par exocytose. 5: Fente synaptique avec acétylcholinestérase. 6: Récepteurs nicotiniques (canaux Na\textsuperscript{+}/K\textsuperscript{+}) sur les replis de la plaque motrice.}
\end{popfigure}

\item \textbf{Chaîne causale explicative :}
\begin{itemize}
\item En fonctionnement normal, l'arrivée du PA au bouton terminal ouvre les canaux \ce{Ca^{2+}} voltage-dépendants ; l'\textbf{entrée de calcium} (\ce{Ca^{2+}}$_{in}$) est le signal déclencheur qui provoque la fusion des vésicules synaptiques avec la membrane présynaptique (\textbf{exocytose}) et la libération d'acétylcholine dans la fente.
\item L'acétylcholine diffuse et se fixe sur les récepteurs nicotiniques de la plaque motrice, provoquant l'ouverture de canaux \ce{Na+}/\ce{K+} : entrée de \ce{Na+} $\rightarrow$ \textbf{dépolarisation postsynaptique (PPSE)} $\rightarrow$ si seuil atteint, déclenchement d'un PA musculaire $\rightarrow$ \textbf{contraction} (couplage excitation-contraction).
\item Si les canaux \ce{Ca^{2+}} sont bloqués (ex. par $\omega$-conotoxine, ou carence en \ce{Ca^{2+}}$_{ext}$), le calcium ne pénètre plus dans le bouton : \textbf{pas d'exocytose}, pas de libération d'acétylcholine, pas de PPSE, pas de PA musculaire, donc \textbf{pas de contraction} (paralysie flasque).
\end{itemize}
\end{enumerate}
\textbf{Conclusion :} l'entrée de calcium dans le bouton présynaptique est l'étape obligatoire qui couple le signal électrique (PA) au signal chimique (libération du neurotransmetteur) ; c'est le maillon « électro-chimique » de la transmission synaptique.
\end{exoresolu}

\begin{methode}[Expliquer les effets d'une substance sur la transmission synaptique]
\begin{enumerate}
\item \textbf{Identifier la cible moléculaire} de la substance :
\begin{itemize}
\item Libération du neurotransmetteur (vésicules, canaux \ce{Ca^{2+}}, protéines SNARE).
\item Récepteur postsynaptique (site orthostérique ou allostérique).
\item Enzyme de dégradation du neurotransmetteur dans la fente (ex. acétylcholinestérase, MAO, COMT).
\item Transporteur de recapture (SERT, DAT, NET, GAT).
\end{itemize}
\item \textbf{Déterminer le sens de l'effet} sur l'efficacité de la transmission :
\begin{itemize}
\item Substance \textbf{agoniste} (mime le neurotransmetteur \emph{ou} augmente sa concentration dans la fente : inhibiteur de dégradation, inhibiteur de recapture, favorise libération) $\rightarrow$ transmission \textbf{augmentée} / potentialisée.
\item Substance \textbf{antagoniste} (bloque le récepteur sans l'activer \emph{ou} diminue la concentration dans la fente : bloque libération, bloque canaux \ce{Ca^{2+}}, détruit neurotransmetteur) $\rightarrow$ transmission \textbf{diminuée ou supprimée}.
\end{itemize}
\item \textbf{Exemples classiques à connaître pour le Bac} :
\begin{itemize}
\item \textbf{Curare} (flèche empoisonnée) : antagoniste compétitif des récepteurs nicotiniques de l'ACh $\rightarrow$ blocage transmission neuromusculaire $\rightarrow$ \textbf{paralysie flasque} (mort par asphyxie).
\item \textbf{Eszérine (physostigmine) / Insecticides organophosphorés / Gaz innervants (Sarin)} : inhibiteurs \textbf{irréversibles} de l'acétylcholinestérase $\rightarrow$ accumulation d'ACh dans la fente $\rightarrow$ surstimulation cholinergique prolongée $\rightarrow$ \textbf{contractions prolongées, fasciculations, hypersécrétions, convulsions} (agoniste indirect).
\item \textbf{Toxine botulique (Botox)} : protéolyse des protéines SNARE $\rightarrow$ blocage de l'exocytose de l'ACh $\rightarrow$ \textbf{paralysie flasque} (antagoniste présynaptique).
\item \textbf{Tétanos (tétanospasmine)} : bloque la libération de GABA/glycine (neuromédiateurs inhibiteurs) au niveau spinal $\rightarrow$ désinhibition motrice $\rightarrow$ \textbf{contractions tétaniques généralisées} (trismus, opisthotonos).
\item \textbf{Caféine} : antagoniste des récepteurs à l'adénosine (A1, A2A) $\rightarrow$ levée de l'inhibition tonique $\rightarrow$ excitation centrale.
\item \textbf{Amphétamines / Cocaïne} : bloquent la recapture (DAT, NET, SERT) et/ou favorisent la libération inverse des monoamines (dopamine, noradrénaline, sérotonine) $\rightarrow$ augmentation concentration fente $\rightarrow$ \textbf{excitation, euphorie, dépendance}.
\end{itemize}
\item \textbf{Rédiger} la chaîne causale complète : Cible moléculaire $\rightarrow$ Conséquence synaptique (concentration médiateur, état récepteur) $\rightarrow$ Effet sur potentiel postsynaptique $\rightarrow$ Effet physiologique/observable.
\end{enumerate}
\end{methode}

\begin{exoresolu}[Intoxication par un insecticide organophosphoré]
Un agriculteur de la zone cacaoyère du Centre-Cameroun, intoxiqué par un insecticide organophosphoré (ex. chlorpyrifos), présente des contractions musculaires involontaires prolongées (fasciculations), une hypersalivation, des larmoiements, une bronchorrhée et des troubles respiratoires. On sait que cette substance inhibe l'acétylcholinestérase. Expliquer les symptômes observés.
\tcblower
\textbf{Solution.} Raisonnement en chaîne causale rigoureuse :
\begin{enumerate}
\item \textbf{Cible :} L'insecticide organophosphoré phosphoryle le site actif de l'\textbf{acétylcholinestérase (AChE)} de façon \textbf{irréversible} (ou quasi-irreversible « vieillissement »), inhibant totalement son activité enzymatique.
\item \textbf{Conséquence synaptique :} Normalement, l'AChE hydrolyse rapidement l'acétylcholine (ACh) en choline + acétate dans la fente synaptique, terminant le signal. L'inhibition de l'AChE entraîne une \textbf{accumulation massive et prolongée d'ACh} dans les fentes synaptiques cholinergiques.
\item \textbf{Effet sur les potentiels postsynaptiques :} Les récepteurs nicotiniques (jonction neuromusculaire, ganglions autonomes) et muscariniques (effecteurs parasympathiques : glandes, muscle lisse, cœur) sont stimulés de façon \textbf{continue et répétée} sans phase de repos. Les PPSE se succèdent sans interruption (dépolarisation maintenue).
\item \textbf{Symptômes cliniques (syndrome cholinergique aigu) :}
\begin{itemize}
\item \textbf{Nicotiniques (squelettique) :} Dépolarisation maintenue de la plaque motrice $\rightarrow$ \textbf{fasciculations} (contractions brèves involontaires), puis \textbf{paralysie par dépolarisation} (blocage de la repolarisation, canaux \ce{Na+} inactivés) $\rightarrow$ détresse respiratoire (paralysie diaphragme/intercostaux).
\item \textbf{Muscariniques (parasympathique) :} Surstimulation des effecteurs $\rightarrow$ \textbf{hypersalivation} (glandes salivaires), \textbf{larmoiements} (glandes lacrymales), \textbf{bronchorrhée} / bronchoconstriction (poumons), \textbf{bradycardie} (cœur), \textbf{diarrhée/coliques} (intestin), \textbf{myosis} (œil).
\end{itemize}
\end{enumerate}
\textbf{Conclusion :} l'inhibition de l'acétylcholinestérase transforme une transmission synaptique normalement brève et phasique en une stimulation tonique permanente : c'est un effet \emph{agoniste indirect} de l'acétylcholine, responsable du syndrome cholinergique aigu (SLUDGE : Salivation, Larmoiement, Urination, Diarrhée, Gastro-intestinal upset, Emesis) observé chez cet agriculteur. Le traitement d'urgence associe atropine (antagoniste muscarinique) et pralidoxime (réactivateur de l'AChE si administré avant « vieillissement »).
\end{exoresolu}

\begin{attention}[Les 5 erreurs qui coûtent des points au Bac]
\begin{enumerate}
\item \textbf{Confondre potentiel de repos et potentiel d'action.} Le potentiel de repos ($-\SI{70}{mV}$) est un état \textbf{permanent} maintenu par la pompe \ce{Na+}/\ce{K+} et les fuites \ce{K+} ; le potentiel d'action est une \emph{variation brève et transitoire} (quelques ms) déclenchée uniquement par un stimulus supra-liminaire. \textbf{Jamais} écrire : « le neurone produit un potentiel de repos en réponse au stimulus ».
\item \textbf{Oublier les unités et les conversions dans le calcul de vitesse.} $v = d/\Delta t$ exige des unités cohérentes (SI : m et s). Convertir : cm $\rightarrow$ m ($\div 100$), ms $\rightarrow$ s ($\div 1000$). Une vitesse de \SI{5000}{m.s^{-1}} obtenue par oubli de conversion (ex. $10/0.002 = 5000$ sans convertir cm en m) doit alerter ! Vérifier l'ordre de grandeur : myélinisé $\approx 10-100$ m/s ; amyélinique $\approx 1-2$ m/s.
\item \textbf{Décrire la courbe du PA sans l'interpréter ioniquement.} Le Bac exige de \textbf{relier chaque phase aux mouvements ioniques et aux canaux} : Dépolarisation = Ouverture canaux \ce{Na+} $\rightarrow$ Entrée \ce{Na+} ; Repolarisation = Fermeture \ce{Na+} + Ouverture \ce{K+} $\rightarrow$ Sortie \ce{K+} ; Hyperpolarisation = Fermeture lente \ce{K+} $\rightarrow$ Sortie \ce{K+} excédentaire ; Retour repos = Pompe \ce{Na+}/\ce{K+}-ATPase. Une description purement morphologique (montée, descente) vaut la moitié des points.
\item \textbf{Inverser le sens de la transmission synaptique ou confondre signal électrique/chimique.} Le neurotransmetteur est \textbf{toujours} libéré par l'élément \emph{pré}synaptique et agit sur les récepteurs de l'élément \emph{post}synaptique : la transmission est \textbf{unidirectionnelle}. Attention : le PA ne « traverse » pas la fente ; c'est un \emph{médiateur chimique} (neurotransmetteur) qui la franchit par diffusion. Ne pas écrire « le potentiel d'action passe la synapse ».
\item \textbf{Schémas non légendés, légendes imprécises ou vocabulaire approximatif.} Tout schéma (neurone, synapse, coupe nerveuse) doit porter : un \textbf{titre} informatif, des \textbf{légendes fléchées} (traits de rappel droits, non croisés), un \textbf{vocabulaire exact} : « bouton synaptique » (pas « boule », « bouton »), « fente synaptique » (pas « trou », « espace »), « gaine de myéline » (pas « enveloppe », « gaine »), « nœud de Ranvier » (pas « étranglement », « constriction »), « plaque motrice » (pas « membrane musculaire »).
\end{enumerate}
\end{attention}

\newpage

\coursec{Exercices}

\courssub{Je vérifie mes connaissances}

\begin{exercice}[QCM : Potentiels membranaires et synapse]
Pour chacune des affirmations suivantes, une seule réponse est correcte. Recopiez le numéro et la lettre correspondante.
\begin{enumerate}
    \item Le potentiel de repos d'un neurone typique est d'environ :
    \begin{itemize}
        \item[A.] $+\SI{70}{mV}$
        \item[B.] $-\SI{70}{mV}$
        \item[C.] $\SI{0}{mV}$
        \item[D.] $-\SI{55}{mV}$
    \end{itemize}
    \item Lors de la phase de dépolarisation du potentiel d'action, l'entrée massive de quel ion est responsable de l'inversion de la polarité de la membrane ?
    \begin{itemize}
        \item[A.] $\ce{K+}$
        \item[B.] $\ce{Cl-}$
        \item[C.] $\ce{Na+}$
        \item[D.] $\ce{Ca^{2+}}$
    \end{itemize}
    \item La gaine de myéline, formée par les cellules de Schwann dans le système nerveux périphérique, a pour rôle principal de :
    \begin{itemize}
        \item[A.] Augmenter la perméabilité de la membrane aux ions $\ce{Na+}$
        \item[B.] Diminuer la capacité membranaire et augmenter la résistance transversale, permettant la conduction saltatoire
        \item[C.] Synthétiser les neurotransmetteurs
        \item[D.] Bloquer la propagation de l'influx nerveux
    \end{itemize}
    \item Dans une synapse à transmission chimique, l'arrivée du potentiel d'action au niveau du bouton terminal provoque l'ouverture de canaux à :
    \begin{itemize}
        \item[A.] $\ce{Na+}$ tension-dépendants
        \item[B.] $\ce{K+}$ tension-dépendants
        \item[C.] $\ce{Ca^{2+}}$ tension-dépendants
        \item[D.] $\ce{Cl-}$ ligand-dépendants
    \end{itemize}
    \item Le curare est un poison qui se fixe sur les récepteurs nicotiniques de l'acétylcholine au niveau de la plaque motrice. Il agit en :
    \begin{itemize}
        \item[A.] Bloquant la recapture de l'acétylcholine
        \item[B.] Inhibant l'acétylcholinestérase
        \item[C.] Empêchant la fixation de l'acétylcholine sur ses récepteurs (antagoniste compétitif)
        \item[D.] Provoquant une libération massive de vésicules synaptiques
    \end{itemize}
\end{enumerate}
\end{exercice}

\begin{exercice}[Vrai ou Faux ? Justifiez]
Indiquez si chaque proposition est vraie (V) ou fausse (F) et justifiez brièvement votre réponse en mobilisant vos connaissances.
\begin{enumerate}
    \item Le potentiel de repos est maintenu uniquement par la pompe $\ce{Na+/K+}$ ATPase.
    \item La période réfractaire absolue correspond au temps pendant lequel les canaux $\ce{Na+}$ tension-dépendants sont inactivés et ne peuvent pas s'ouvrir, quel que soit l'intensité du stimulus.
    \item La vitesse de conduction de l'influx nerveux est proportionnelle au carré du diamètre de l'axone pour les fibres amyéliniques.
    \item Un potentiel postsynaptique excitateur (PPSE) correspond toujours à une dépolarisation qui déclenche systématiquement un potentiel d'action dans le neurone postsynaptique.
\end{enumerate}
\end{exercice}

\begin{exercice}[Définitions et vocabulaire précis]
Donnez une définition rigoureuse pour chacun des termes suivants, en utilisant le vocabulaire scientifique approprié (grandeurs, unités, mécanismes) :
\begin{enumerate}
    \item Potentiel de repos ($V_{\text{repos}}$)
    \item Potentiel d'action (PA)
    \item Conduction saltatoire
    \item Fente synaptique
    \item Acétylcholinestérase
\end{enumerate}
\end{exercice}

\begin{exercice}[Calcul direct : Vitesse de conduction]
Lors d'une expérience de neurophysiologie sur une fibre nerveuse de grenouille, deux électrodes d'enregistrement extracellulaires sont placées à une distance $d = \SI{12,0}{cm}$ l'une de l'autre. L'enregistrement simultané montre que le pic du potentiel d'action passe à la première électrode à l'instant $t_1 = \SI{2,50}{ms}$ et à la seconde à $t_2 = \SI{10,50}{ms}$.
\begin{enumerate}
    \item Calculer la vitesse de conduction $v$ de l'influx nerveux dans cette fibre (en $\text{m}\cdot\text{s}^{-1}$).
    \item Sachant que les fibres myéliniques de grenouille ont des vitesses comprises entre $10$ et $30\,\text{m}\cdot\text{s}^{-1}$ et les fibres amyéliniques entre $0,5$ et $2\,\text{m}\cdot\text{s}^{-1}$, cette fibre est-elle probablement myélinisée ou amyélinisée ? Justifiez.
\end{enumerate}
\end{exercice}

\courssub{Je m'entraîne}

\begin{exercice}[Analyse d'un enregistrement de potentiel d'action]
On enregistre le potentiel de membrane d'un axone géant de calmar (axone non myélinisé) stimulé par une électrode intracellulaire. La courbe obtenue (potentiel $V_m$ en fonction du temps $t$) présente les caractéristiques suivantes :
\begin{itemize}
    \item Phase 1 : Plateau à $-\SI{70}{mV}$ de $t=0$ à $t=\SI{1}{ms}$.
    \item Phase 2 : Montée brutale jusqu'à $+\SI{35}{mV}$ à $t=\SI{1,5}{ms}$.
    \item Phase 3 : Redescente rapide, franchissant $\SI{0}{mV}$ vers $t=\SI{2,0}{ms}$, atteignant $-\SI{85}{mV}$ à $t=\SI{3,0}{ms}$.
    \item Phase 4 : Remontée lente vers $-\SI{70}{mV}$ vers $t=\SI{5}{ms}$.
\end{itemize}
\begin{enumerate}
    \item Identifiez et nommez les quatre phases du potentiel d'action représentées.
    \item Quelle est la valeur du potentiel de repos ? Quelle est l'amplitude du potentiel d'action (différence entre le pic de dépolarisation et le potentiel de repos) ?
    \item Pour chaque phase (2, 3, 4), précisez les mouvements ioniques principaux (ions concernés, sens : entrée ou sortie) et l'état des canaux tension-dépendants ($\ce{Na+}$ et $\ce{K+}$ : fermés, ouverts, inactivés).
    \item Expliquez pourquoi la phase 4 (hyperpolarisation) est transitoire.
\end{enumerate}
\end{exercice}

\begin{exercice}[Détermination de la vitesse de conduction et nature de la fibre]
On étudie la conduction de l'influx nerveux dans le nerf sciatique d'une grenouille. Deux électrodes d'enregistrement $E_1$ et $E_2$ sont placées sur le nerf à une distance $D = \SI{8,0}{cm}$. On stimule le nerf en amont de $E_1$. Les enregistrements montrent un décalage temporel $\Delta t = \SI{0,80}{ms}$ entre l'arrivée du potentiel d'action à $E_1$ et son arrivée à $E_2$.
\begin{enumerate}
    \item Calculer la vitesse de conduction $v$ en $\text{m}\cdot\text{s}^{-1}$.
    \item En comparant aux valeurs de référence (fibres A$\alpha$ myélinisées : $70\text{--}120\,\text{m}\cdot\text{s}^{-1}$ ; fibres C amyéliniques : $0,5\text{--}2\,\text{m}\cdot\text{s}^{-1}$), proposer une classification pour cette fibre (type A, B ou C ; myélinisée ou non).
    \item Schématisez (schéma fonctionnel simplifié) la propagation de l'influx le long de cette fibre en indiquant le sens des courants locaux (circuits de courant) et le rôle des nœuds de Ranvier si la fibre est myélinisée.
\end{enumerate}
\end{exercice}

\begin{exercice}[Fonctionnement de la synapse chimique : de l'arrivée du PA au PPSE]
On considère une synapse cholinergique (neuromusculaire) entre un motoneurone et une fibre musculaire squelettique.
\begin{enumerate}
    \item Réalisez un schéma annoté et légendé de cette synapse en coupe, en faisant apparaître : membrane présynaptique, vésicules synaptiques, mitochondries, fente synaptique, membrane postsynaptique (plaque motrice), récepteurs nicotiniques, canaux ioniques ligand-dépendants.
    \item Décrivez chronologiquement les étapes de la transmission synaptique depuis l'arrivée du potentiel d'action au niveau du bouton terminal jusqu'à l'apparition du potentiel postsynaptique excitateur (PPSE) sur la membrane musculaire. Utilisez les termes : dépolarisation, canaux $\ce{Ca^{2+}}$ tension-dépendants, influx $\ce{Ca^{2+}}$, exocytose, diffusion, fixation récepteur, ouverture canaux $\ce{Na+}/\ce{K+}$, dépolarisation postsynaptique.
    \item Quelle est la nature ionique du PPSE à la plaque motrice ? Pourquoi est-il toujours dépolarisant ?
\end{enumerate}
\end{exercice}

\begin{exercice}[Effets de substances pharmacologiques sur la transmission neuromusculaire]
On réalise des expériences sur une préparation isolée muscle-nerf (nerf phrénique - diaphragme de rat) stimulée électriquement à une fréquence de $\SI{1}{Hz}$. On mesure la tension de contraction isométrique.
\begin{enumerate}
    \item On perfuse le bain avec une solution contenant du \textbf{curare} ($\SI{10^{-6}}{mol\cdot L^{-1}}$). La tension de contraction diminue progressivement jusqu'à disparaître en $\SI{5}{min}$. Après lavage abondant, la contraction réapparaît. Interprétez le mécanisme d'action du curare au niveau moléculaire.
    \item Dans une autre expérience, on ajoute de la \textbf{néostigmine} (inhibiteur réversible de l'acétylcholinestérase) à $\SI{10^{-5}}{mol\cdot L^{-1}}$. On observe une augmentation de l'amplitude de la contraction puis, à fortes doses, un blocage (dépolarisation permanente). Expliquez ces deux effets paradoxaux.
    \item Le \textbf{venin de botulisme} (toxine botulique) bloque la libération des vésicules synaptiques. Prédisez l'effet sur la contraction musculaire et proposez une application thérapeutique médicale (autre que cosmétique) de cette toxine.
\end{enumerate}
\end{exercice}

\courssub{Je me prépare au Bac}

\begin{exercice}[Sujet type Bac : Potentiel d'action, conduction et myélinisation]
\textbf{Contexte :} Le nerf sciatique de la grenouille (\emph{Xenopus laevis}) est un modèle classique en neurophysiologie. Il contient des fibres myélinisées (type A) et amyélinisées (type C). On souhaite comparer leurs propriétés électriques.

\textbf{Document 1 : Enregistrement intracellulaire sur une fibre A (myélinisée).}
On stimule la fibre par un courant bref. L'enregistrement du potentiel de membrane $V_m$ (en mV) en fonction du temps $t$ (en ms) donne le tableau ci-dessous :

\begin{center}
\begin{tabular}{|c|c|c|c|c|c|c|c|c|}\hline
\rowcolor{popblueL}
$t\,(\text{ms})$ & 0,0 & 0,2 & 0,4 & 0,6 & 0,8 & 1,0 & 1,2 & 1,4 \\ \hline
$V_m\,(\text{mV})$ & -70 & -65 & +20 & +35 & +10 & -60 & -80 & -72 \\ \hline
\end{tabular}
\end{center}

\textbf{Document 2 : Mesure de la vitesse de conduction.}
Deux électrodes d'enregistrement extracellulaires $E_1$ et $E_2$ sont placées sur le nerf à une distance $d = \SI{10,0}{cm}$. On stimule le nerf en amont de $E_1$. Les signaux enregistrés montrent un décalage $\Delta t_A = \SI{3,3}{ms}$ pour les fibres A et $\Delta t_C = \SI{67}{ms}$ pour les fibres C.

\textbf{Document 3 : Expérience de double stimulation (période réfractaire).}
On applique deux stimuli brefs identiques ($S_1$ et $S_2$) espacés d'un intervalle $\Delta t_{\text{stim}}$ sur une fibre A. On observe :
\begin{itemize}
    \item Pour $\Delta t_{\text{stim}} = \SI{0,5}{ms}$ : seul $S_1$ déclenche un PA.
    \item Pour $\Delta t_{\text{stim}} = \SI{2,0}{ms}$ : $S_1$ et $S_2$ déclenchent chacun un PA, mais l'amplitude du PA déclenché par $S_2$ est légèrement réduite.
    \item Pour $\Delta t_{\text{stim}} = \SI{5,0}{ms}$ : deux PA identiques sont enregistrés.
\end{itemize}

\textbf{Questions :}
\begin{enumerate}
    \item \textbf{Analyse du potentiel d'action (Doc 1).}
    \begin{enumerate}
        \item Représentez schématiquement la courbe $V_m = f(t)$ en utilisant les valeurs du tableau. Indiquez l'échelle des axes.
        \item Déterminez graphiquement ou par lecture : la valeur du potentiel de repos $V_{\text{repos}}$, l'amplitude du PA, la durée du PA (à la base), la durée de la phase de dépolarisation.
        \item Identifiez les phases de dépolarisation, repolarisation et hyperpolarisation. Pour chacune, citez l'ion principal responsable et son sens de déplacement (entrée/sortie).
    \end{enumerate}
    \item \textbf{Vitesse de conduction (Doc 2).}
    \begin{enumerate}
        \item Calculez la vitesse de conduction $v_A$ des fibres A et $v_C$ des fibres C (en $\text{m}\cdot\text{s}^{-1}$).
        \item Calculez le rapport $v_A / v_C$. Quel est l'avantage sélectif majeur de la myélinisation pour un animal prédateur comme la grenouille ?
    \end{enumerate}
    \item \textbf{Périodes réfractaires (Doc 3).}
    \begin{enumerate}
        \item Définissez la période réfractaire absolue (PRA) et la période réfractaire relative (PRR).
        \item Estimez la durée de la PRA et de la PRR pour cette fibre A à partir des résultats.
        \item Expliquez la base biophysique de la PRA (état des canaux $\ce{Na+}$) et de la PRR (état des canaux $\ce{K+}$ et potentiel de membrane).
    \end{enumerate}
    \item \textbf{Synthèse.} En vous appuyant sur les résultats, expliquez pourquoi la conduction saltatoire dans les fibres myélinisées est plus rapide et plus économe en énergie (consommation d'ATP par la pompe $\ce{Na+/K+}$) que la conduction continue dans les fibres amyélinisées.
\end{enumerate}
\end{exercice}

\begin{exercice}[Sujet type Bac : Preuve expérimentale de la transmission chimique et pharmacologie]
\textbf{Contexte :} En 1921, Otto Loewi réalise une expérience décisive sur le cœur de grenouille (\emph{Rana temporaria}) qui lui vaudra le prix Nobel. Le cœur de grenouille continue de battre in vitro s'il est perfusé par une solution saline (Ringer).

\textbf{Document 1 : Expérience de Loewi (schéma simplifié).}
\begin{itemize}
    \item \textbf{Expérience A :} Un cœur donneur (avec nerf vague intact) et un cœur receveur (dénervé) sont perfusés en série par le même fluide Ringer. On stimule électriquement le nerf vague du cœur donneur. \textbf{Résultat :} Le cœur donneur ralentit (bradycardie), puis \textbf{le cœur receveur ralentit également} après un court délai.
    \item \textbf{Expérience B :} Même montage, mais on intercale un filtre à bactéries (pore $< \SI{0,2}{\mu m}$) entre les deux cœurs. \textbf{Résultat :} Le cœur donneur ralentit, et \textbf{le cœur receveur ralentit également}.
    \item \textbf{Expérience C :} On collecte le liquide de perfusion sortant du cœur donneur \emph{après} stimulation du vague. On l'applique sur un cœur receveur isolé (non perfusé en série). \textbf{Résultat :} Le cœur receveur ralentit.
\end{itemize}

\textbf{Document 2 : Identification de la substance (années 1930).}
Dale et Loewi identifient la substance libérée : l'acétylcholine (ACh). Ils montrent que l'atropine (alcaloïde de la belladone) bloque l'effet de l'ACh sur le cœur, tandis que la physostigmine (inhibiteur de l'acétylcholinestérase) potentialise et prolonge l'effet de la stimulation vagale.

\textbf{Document 3 : Application clinique - Myasthénie grave.}
Maladie auto-immune où des anticorps bloquent les récepteurs nicotiniques de l'ACh à la plaque motrice. Traitement : inhibiteurs de l'acétylcholinestérase (pyridostigmine).

\textbf{Questions :}
\begin{enumerate}
    \item \textbf{Analyse de l'expérience de Loewi (Doc 1).}
    \begin{enumerate}
        \item Quelle conclusion fondamentale sur la nature de la transmission nerveuse (électrique vs chimique) tire-t-on de l'expérience A ?
        \item Quel est le rôle du filtre à bactéries dans l'expérience B ? Pourquoi le résultat de B (ralentissement du receveur maintenu) invalide-t-il l'hypothèse d'un messager cellulaire (bactérie/virus) et confirme-t-il la nature chimique du messager ?
        \item Que prouve l'expérience C quant à la nature chimique du messager et à sa stabilité ?
    \end{enumerate}
    \item \textbf{Mécanismes moléculaires (Doc 2).}
    \begin{enumerate}
        \item L'ACh agit sur des récepteurs muscariniques au niveau cardiaque. Ces récepteurs sont-ils des canaux ioniques ligand-dépendants ? Justifiez (voie de signalisation métabotropique via protéine G).
        \item Expliquez l'effet de l'atropine (antagoniste compétitif) et de la physostigmine (inhibiteur de l'AChE) sur la réponse cardiaque à la stimulation vagale.
    \end{enumerate}
    \item \textbf{Pathologie et traitement (Doc 3).}
    \begin{enumerate}
        \item Expliquez le mécanisme physiopathologique de la myasthénie grave au niveau de la plaque motrice (synapse neuromusculaire). Pourquoi la fatigue musculaire s'aggrave-t-elle à l'effort ?
        \item Justifiez l'utilisation d'inhibiteurs de l'acétylcholinestérase (comme la pyridostigmine) comme traitement symptomatique. Précisez l'effet sur la concentration d'ACh dans la fente synaptique et sur l'amplitude du PPSE.
    \end{enumerate}
\end{enumerate}
\end{exercice}

\begin{exercice}[Sujet type Bac : Intégration - Sclérose en plaques, conduction et potentiel d'action]
\textbf{Contexte :} La sclérose en plaques (SEP) est une maladie démyélinisante auto-immune du système nerveux central. Elle touche préférentiellement les jeunes adultes (pic vers 30 ans). Au Cameroun, sa prévalence est estimée à environ $10\text{--}15$ cas pour $100\,000$ habitants, avec une latence diagnostique souvent longue.

\textbf{Document 1 : Histologie et IRM.}
\begin{itemize}
    \item Coupe de substance blanche saine : axones entourés d'une gaine de myéline compacte (lamelles lipidiques), nœuds de Ranvier réguliers espacés de $1\text{--}2\,\text{mm}$.
    \item Coupe de plaque SEP : démyélinisation focale, infiltration lymphocytaire, gliose astrocytaire. L'axone est nu sur plusieurs millimètres.
    \item IRM cérébrale (pondération T2) : hypersignaux multiples en périventriculaire, corps calleux, tronc cérébral (plaques de démyélinisation).
\end{itemize}

\textbf{Document 2 : Électrophysiologie - Potentiels évoqués visuels (PEV).}
On stimule l'œil par un damier inversé et on enregistre l'activité électrique du cortex visuel (occipital).
\begin{itemize}
    \item Sujet témoin : latence du pic P100 = $\SI{100}{ms}$.
    \item Patient SEP (poussée aiguë névrite optique) : latence du pic P100 = $\SI{135}{ms}$, amplitude réduite.
    \item Patient SEP (rémission) : latence P100 = $\SI{115}{ms}$, amplitude quasi normale.
\end{itemize}

\textbf{Document 3 : Modélisation de la conduction.}
On modélise un axone de $\SI{1}{m}$ de long (voie optique).
\begin{itemize}
    \item Fibre saine (myélinisée) : vitesse $v_m = \SI{50}{m\cdot s^{-1}}$, capacitance membranaire au nœud $C_m = \SI{2}{nF}$, résistance axoplasmique $R_a = \SI{100}{M\Omega}$ par internœud.
    \item Fibre démyélinisée (zone de plaque) : vitesse $v_d = \SI{1}{m\cdot s^{-1}}$, capacitance membranaire $C_m' = \SI{200}{nF}$ (membrane nue sur grande surface), pas de nœuds.
\end{itemize}

\textbf{Questions :}
\begin{enumerate}
    \item \textbf{Physiopathologie (Doc 1).}
    \begin{enumerate}
        \item Rappellez le rôle isolant de la myéline (capacité membranaire $C_m$, résistance transversale $R_m$). Pourquoi la démyélinisation expose-t-elle la membrane axoplasmique à une capacitance effective massivement augmentée ?
        \item Expliquez pourquoi la démyélinisation entraîne une augmentation du seuil d'excitabilité et un risque de bloc de conduction (fatigabilité).
    \end{enumerate}
    \item \textbf{Analyse des PEV (Doc 2).}
    \begin{enumerate}
        \item Calculez le temps de conduction $t_m$ (sain) et $t_d$ (démylinisé) pour un trajet de $\SI{0,05}{m}$ (longueur approximative du nerf optique intra-crânien) en utilisant les vitesses du Doc 3.
        \item Comparez ces temps théoriques aux latences mesurées (P100). Pourquoi la latence P100 inclut-elle d'autres délais que la simple conduction axonale ?
        \item Interprétez l'évolution de la latence et de l'amplitude entre la poussée aiguë et la rémission (rémyélinisation incomplète, redistribution des canaux $\ce{Na+}$).
    \end{enumerate}
    \item \textbf{Énergétique et modélisation (Doc 3).}
    \begin{enumerate}
        \item La pompe $\ce{Na+/K+}$ ATPase consomme 1 ATP pour extruder 3 $\ce{Na+}$ et rentrer 2 $\ce{K+}$. Lors d'un PA, l'influx de $\ce{Na+}$ est proportionnel à la capacitance membranaire ($Q = C_m \Delta V$). Comparez qualitativement la charge $\ce{Na+}$ entrant par PA au niveau d'un nœud de Ranvier (sain) vs sur une zone démyélinisée de même longueur. En déduire l'impact sur la consommation d'ATP.
        \item Proposez un mécanisme explicatif pour la survenue de la fatigue centrale (asthénie majeure) chez les patients SEP, en lien avec ce coût énergétique accru de la conduction.
    \end{enumerate}
    \item \textbf{Traitement symptomatique.} La fampridine (4-aminopyridine) bloque les canaux $\ce{K+}$ tension-dépendants. Expliquez comment ce blocage peut améliorer la conduction dans une fibre démyélinisée (élargissement du PA, augmentation du courant axoplasmique) et citez une indication clinique validée (marche).
\end{enumerate}
\end{exercice}

\newpage

\coursec{Corrigés des exercices}

\coursec{Exercices d'application et corrigés détaillés}

\courssub{Entraînement : QCM, Vrai/Faux et Définitions}

\begin{corrige}[Exercice 1 — QCM : Potentiels membranaires et synapse]
\begin{enumerate}
    \item \textbf{Réponse B} : $-\SI{70}{mV}$. Le potentiel de repos d'un neurone typique est d'environ $-\SI{70}{mV}$ (intérieur négatif par rapport à l'extérieur). La valeur $-\SI{55}{mV}$ correspond au potentiel seuil, et non au potentiel de repos.
    \item \textbf{Réponse C} : $\ce{Na+}$. La dépolarisation est due à l'ouverture des canaux $\ce{Na+}$ tension-dépendants et à l'entrée massive d'ions $\ce{Na+}$ (le potentiel d'équilibre du sodium est voisin de $+\SI{60}{mV}$), ce qui inverse la polarité membranaire jusqu'à environ $+\SI{30}{mV}$ à $+\SI{40}{mV}$.
    \item \textbf{Réponse B}. La myéline est un isolant électrique : elle diminue la capacité membranaire $C_m$ et augmente la résistance transversale $R_m$, ce qui confine le courant dans l'axoplasme et impose la régénération du PA uniquement aux nœuds de Ranvier : c'est la conduction saltatoire, rapide et économe.
    \item \textbf{Réponse C} : canaux $\ce{Ca^{2+}}$ tension-dépendants. L'arrivée du PA dépolarise le bouton terminal, ce qui ouvre ces canaux ; l'influx de $\ce{Ca^{2+}}$ déclenche la fusion des vésicules synaptiques avec la membrane présynaptique (exocytose).
    \item \textbf{Réponse C}. Le curare est un antagoniste compétitif : il se fixe sur les récepteurs nicotiniques sans les activer et empêche l'ACh de s'y fixer, bloquant ainsi la transmission neuromusculaire (paralysie flasque). Il n'agit ni sur la recapture, ni sur l'acétylcholinestérase, ni sur l'exocytose.
\end{enumerate}
\end{corrige}

\begin{corrige}[Exercice 2 — Vrai ou Faux ? Justifiez]
\begin{enumerate}
    \item \textbf{FAUX.} La pompe $\ce{Na+/K+}$ ATPase (3 $\ce{Na+}$ sortis contre 2 $\ce{K+}$ rentrés par ATP hydrolysé) ne contribue directement que pour quelques millivolts (environ $-3$ à $-\SI{5}{mV}$). L'essentiel du potentiel de repos provient des \cle{fuites ioniques passives} à travers les canaux de fuite, principalement les canaux $\ce{K+}$ de fuite : la membrane au repos est beaucoup plus perméable à $\ce{K+}$ qu'à $\ce{Na+}$, d'où un potentiel proche du potentiel d'équilibre du potassium ($E_K \approx -\SI{90}{mV}$). La pompe maintient les gradients de concentration à long terme, mais ne crée pas directement l'essentiel de $V_{\text{repos}}$.
    \item \textbf{VRAI.} Pendant la PRA (environ $\SI{1}{ms}$), les canaux $\ce{Na+}$ tension-dépendants sont soit ouverts, soit \cle{inactivés} (bille d'inactivation refermée) ; or un canal inactivé ne peut pas être rouvert, quelle que soit l'intensité du stimulus. Aucun nouveau PA ne peut donc être déclenché. Cette propriété impose le sens unidirectionnel de propagation et limite la fréquence maximale des PA.
    \item \textbf{FAUX.} Pour les fibres amyéliniques, la vitesse de conduction est proportionnelle à la \textbf{racine carrée} du diamètre ($v \propto \sqrt{d}$), et non au carré. Pour les fibres myélinisées, la vitesse est approximativement proportionnelle au diamètre ($v \propto d$). L'affirmation double donc l'erreur.
    \item \textbf{FAUX.} Un PPSE est bien une dépolarisation locale (entrée nette de cations), mais il est \cle{gradué} et \textbf{décrémentiel} : il ne déclenche un PA que si la sommation (spatiale et/ou temporelle) des PPSE au niveau du cône d'émergence de l'axone atteint le potentiel seuil (environ $-\SI{55}{mV}$). Un PPSE isolé, sous-liminaire, s'éteint sans déclencher de PA.
\end{enumerate}
\end{corrige}

\begin{corrige}[Exercice 3 — Définitions et vocabulaire précis]
\begin{enumerate}
    \item \textbf{Potentiel de repos ($V_{\text{repos}}$)} : différence de potentiel électrique transmembranaire mesurée à l'état non stimulé, l'intérieur de la cellule étant négatif par rapport à l'extérieur ; sa valeur est d'environ $-\SI{70}{mV}$ chez le neurone. Il résulte de la perméabilité sélective de la membrane (dominance des canaux $\ce{K+}$ de fuite) et des gradients ioniques maintenus par la pompe $\ce{Na+/K+}$ ATPase.
    \item \textbf{Potentiel d'action (PA)} : variation brève, brutale et stéréotypée (loi du « tout ou rien ») du potentiel de membrane, déclenchée lorsque la dépolarisation atteint le seuil (environ $-\SI{55}{mV}$) ; il comporte une dépolarisation (entrée de $\ce{Na+}$ jusqu'à environ $+\SI{30}{mV}$), une repolarisation (sortie de $\ce{K+}$) et souvent une hyperpolarisation transitoire. Durée : $1$ à $\SI{2}{ms}$ ; il se propage sans décrément.
    \item \textbf{Conduction saltatoire} : mode de propagation du PA le long des fibres myélinisées, dans lequel le PA n'est régénéré qu'au niveau des nœuds de Ranvier (zones dépourvues de myéline, riches en canaux $\ce{Na+}$), le courant se propageant passivement et rapidement sous la gaine isolante entre deux nœuds. Le PA semble ainsi « sauter » de nœud en nœud, d'où des vitesses de $15$ à $120\,\text{m}\cdot\text{s}^{-1}$.
    \item \textbf{Fente synaptique} : espace extracellulaire étroit (environ $20$ à $\SI{40}{nm}$) séparant la membrane présynaptique de la membrane postsynaptique dans une synapse chimique ; c'est dans cet espace que diffuse le neurotransmetteur libéré par exocytose et qu'il est dégradé ou recapturé.
    \item \textbf{Acétylcholinestérase (AChE)} : enzyme présente dans la fente synaptique (notamment fixée sur la lame basale de la plaque motrice) qui hydrolyse l'acétylcholine en acétate et choline ; elle assure l'arrêt du signal synaptique et permet le recyclage du choline par recapture présynaptique. Son inhibition (néostigmine, organophosphorés) prolonge l'action de l'ACh.
\end{enumerate}
\end{corrige}

\courssub{Applications numériques et analyse de documents}

\begin{corrige}[Exercice 4 — Calcul direct : Vitesse de conduction]
\begin{enumerate}
    \item La vitesse de conduction est le quotient de la distance inter-électrodes par le temps de propagation :
    \[ v = \frac{d}{t_2 - t_1} = \frac{\SI{12,0}{cm}}{\SI{10,50}{ms} - \SI{2,50}{ms}} = \frac{\SI{0,120}{m}}{\SI{8,00e-3}{s}} = \SI{15,0}{m\cdot s^{-1}}. \]
    \item La valeur $v = \SI{15,0}{m\cdot s^{-1}}$ appartient à l'intervalle $10$--$30\,\text{m}\cdot\text{s}^{-1}$ caractéristique des fibres \textbf{myélinisées} de grenouille, et est très supérieure à la fourchette $0,5$--$2\,\text{m}\cdot\text{s}^{-1}$ des fibres amyéliniques. Conclusion : cette fibre est très probablement \cle{myélinisée} (conduction saltatoire).
\end{enumerate}
\end{corrige}

\begin{attention}[Point de vigilance]
Convertir systématiquement la distance en mètres et le temps en secondes avant le calcul ; une erreur fréquente consiste à diviser des cm par des ms et à oublier le facteur $10$ (ici $\SI{12}{cm}/\SI{8}{ms} = \SI{1,5}{cm\cdot ms^{-1}} = \SI{15}{m\cdot s^{-1}}$).
\end{attention}

\begin{corrige}[Exercice 5 — Analyse d'un enregistrement de potentiel d'action]
\begin{enumerate}
    \item Identification des phases :
    \begin{itemize}
        \item Phase 1 ($0$ à $\SI{1}{ms}$, $-\SI{70}{mV}$) : \textbf{potentiel de repos} (état polarisé).
        \item Phase 2 ($\SI{1}{ms}$ à $\SI{1,5}{ms}$, montée vers $+\SI{35}{mV}$) : \textbf{dépolarisation}.
        \item Phase 3 ($\SI{1,5}{ms}$ à $\SI{3,0}{ms}$, redescente vers $-\SI{85}{mV}$) : \textbf{repolarisation} puis début d'\textbf{hyperpolarisation}.
        \item Phase 4 ($\SI{3}{ms}$ à $\SI{5}{ms}$, remontée vers $-\SI{70}{mV}$) : fin de l'\textbf{hyperpolarisation} et retour au repos.
    \end{itemize}
    \item Potentiel de repos : $V_{\text{repos}} = -\SI{70}{mV}$. Amplitude du PA :
    \[ A = V_{\text{pic}} - V_{\text{repos}} = (+35) - (-70) = \SI{105}{mV}. \]
    \item Mouvements ioniques et état des canaux :
    \begin{center}
    \begin{tabularx}{\linewidth}{|l|X|X|X|}\hline
    \rowcolor{popblueL}
    \textbf{Phase} & \textbf{Ion dominant} & \textbf{Canaux $\ce{Na+}$} & \textbf{Canaux $\ce{K+}$} \\ \hline
    2 — Dépolarisation & Entrée massive de $\ce{Na+}$ & Ouverts (activation rapide) & Fermés (activation lente) \\ \hline
    3 — Repolarisation & Sortie de $\ce{K+}$ & Inactivés & Ouverts \\ \hline
    4 — Hyperpolarisation & Sortie de $\ce{K+}$ persistante & Fermés (désinactivation) & Ouverts (fermeture lente) \\ \hline
    \end{tabularx}
    \end{center}
    \item L'hyperpolarisation est transitoire car les canaux $\ce{K+}$ tension-dépendants finissent par se \textbf{refermer} (ils répondent lentement à la dépolarisation, puis se ferment lorsque la membrane se repolarise). Une fois fermés, la conductance à $\ce{K+}$ revient à sa valeur de repos ; les canaux de fuite et la pompe $\ce{Na+/K+}$ ramènent alors le potentiel à $-\SI{70}{mV}$. L'hyperpolarisation n'est donc qu'un dépassement temporaire dû à la fermeture retardée des canaux $\ce{K+}$.
\end{enumerate}
\end{corrige}

\begin{corrige}[Exercice 6 — Détermination de la vitesse de conduction et nature de la fibre]
\begin{enumerate}
    \item Calcul de la vitesse :
    \[ v = \frac{D}{\Delta t} = \frac{\SI{8,0}{cm}}{\SI{0,80}{ms}} = \frac{\SI{0,080}{m}}{\SI{8,0e-4}{s}} = \SI{100}{m\cdot s^{-1}}. \]
    \item $v = \SI{100}{m\cdot s^{-1}}$ appartient à la gamme $70$--$120\,\text{m}\cdot\text{s}^{-1}$ : il s'agit d'une fibre de \textbf{type A$\alpha$}, donc \textbf{myélinisée} (grosse fibre motrice ou sensitive). Ce n'est ni une fibre B (myélinisée fine, $3$--$15\,\text{m}\cdot\text{s}^{-1}$), ni une fibre C (amyélinique, $0,5$--$2\,\text{m}\cdot\text{s}^{-1}$).
    \item Schéma fonctionnel de la conduction saltatoire :
    \begin{popfigure}
    \begin{tikzpicture}[x=1cm,y=1cm]
        % Axone
        \draw[thick,popdark] (0,0) -- (12,0);
        \draw[thick,popdark] (0,0.7) -- (12,0.7);
        % Gaine de myéline (internœuds)
        \foreach \x in {0.4, 4.4, 8.4}{
            \draw[fill=popblueL,draw=popblue] (\x,-0.25) rectangle (\x+2.6,0.95);
        }
        % Nœuds de Ranvier
        \foreach \x in {3.2, 7.2, 11.2}{
            \draw[fill=poporange,draw=poporange] (\x,-0.25) rectangle (\x+0.6,0.95);
        }
        % Labels
        \node[popblue,font=\small] at (1.7,1.5) {gaine de myéline};
        \node[poporange,font=\small] at (3.5,1.6) {nœud de Ranvier};
        \node[font=\small] at (3.5,-0.7) {PA régénéré};
        % Flèches de propagation
        \draw[-{Stealth},very thick,poppurple] (0.3,2.1) -- (11.7,2.1) node[midway,above,font=\small]{sens de propagation};
        % Courants locaux
        \draw[-{Stealth},thick,popgreen] (3.5,0.35) -- (7.0,0.35);
        \draw[-{Stealth},thick,popgreen] (7.5,0.35) -- (11.0,0.35);
        \node[popgreen,font=\small] at (5.3,-0.35) {courant local axoplasmique};
        % Retour extracellulaire
        \draw[-{Stealth},thick,popgreen,dashed] (7.0,-1.2) -- (3.7,-1.2);
        \node[popgreen,font=\small] at (5.3,-1.55) {courant de retour extracellulaire};
    \end{tikzpicture}
    \legende{Conduction saltatoire : le PA n'est régénéré qu'aux nœuds de Ranvier (orange, riches en canaux $\ce{Na+}$) ; le courant local (vert) se propage passivement et rapidement sous la gaine de myéline (bleu) jusqu'au nœud suivant, qu'il dépolarise jusqu'au seuil.}
    \end{popfigure}
    Le courant entrant au nœud actif s'écoule dans l'axoplasme vers le nœud suivant (circuit bouclé par le liquide extracellulaire) ; la myéline empêche les fuites transmembranaires, d'où une propagation rapide de nœud en nœud.
\end{enumerate}
\end{corrige}

\begin{corrige}[Exercice 7 — Fonctionnement de la synapse chimique : de l'arrivée du PA au PPSE]
\begin{enumerate}
    \item Schéma de la synapse neuromusculaire :
    \begin{popfigure}
    \begin{tikzpicture}[x=1cm,y=1cm]
        % Bouton présynaptique
        \draw[fill=popblueL,draw=popblue,thick] (0,2.2) .. controls (0,4.4) and (6,4.4) .. (6,2.2) -- (6,1.6) .. controls (6,1.2) and (0,1.2) .. (0,1.6) -- cycle;
        % Vésicules
        \foreach \p in {(1,3),(2,3.4),(3,2.9),(4,3.5),(5,3),(2.5,2.2),(4.5,2.1)}{
            \draw[fill=poppinkL,draw=poppink] \p circle (0.28);
        }
        % Mitochondrie
        \draw[fill=popgoldL,draw=popgold,thick] (0.7,3.6) ellipse (0.55 and 0.3);
        \draw[popgold] (0.35,3.6) .. controls (0.6,3.75) and (0.8,3.45) .. (1.05,3.6);
        % Membrane présynaptique
        \draw[very thick,popblue] (0,1.35) -- (6,1.35);
        % Fente synaptique
        \draw[fill=popcream,draw=none] (0,0.55) rectangle (6,1.35);
        % Membrane postsynaptique (plissements)
        \draw[very thick,popgreen] (0,0.55) -- (1,0.55) -- (1.2,0.25) -- (1.6,0.25) -- (1.8,0.55) -- (3,0.55) -- (3.2,0.25) -- (3.6,0.25) -- (3.8,0.55) -- (6,0.55);
        \draw[fill=popgreenL,draw=popgreen] (0,0.25) rectangle (0.9,0.55);
        \draw[fill=popgreenL,draw=popgreen] (1.9,0.25) rectangle (2.9,0.55);
        \draw[fill=popgreenL,draw=popgreen] (3.9,0.25) rectangle (6,0.55);
        % Récepteurs nicotiniques
        \foreach \x in {1.4, 3.4, 5.0}{
            \draw[fill=poporange,draw=poporange,thick] (\x,0.35) rectangle (\x+0.18,1.0);
            \draw[fill=poporange,draw=poporange,thick] (\x+0.35,0.35) rectangle (\x+0.53,1.0);
        }
        % Légendes fléchées
        \node[font=\small,anchor=west] at (6.4,3.6) {1 — bouton présynaptique};
        \node[font=\small,anchor=west] at (6.4,3.0) {2 — vésicules d'ACh};
        \node[font=\small,anchor=west] at (6.4,2.4) {3 — mitochondrie};
        \node[font=\small,anchor=west] at (6.4,1.35) {4 — membrane présynaptique};
        \node[font=\small,anchor=west] at (6.4,0.9) {5 — fente synaptique};
        \node[font=\small,anchor=west] at (6.4,0.35) {6 — récepteurs nicotiniques};
        \node[font=\small,anchor=west] at (6.4,-0.25) {7 — membrane postsynaptique (plaque motrice)};
        \draw[-{Stealth},popmuted] (6.35,3.0) -- (4.6,3.1);
        \draw[-{Stealth},popmuted] (6.35,2.4) -- (1.3,3.6);
        \draw[-{Stealth},popmuted] (6.35,0.9) -- (3.5,1.05);
        \draw[-{Stealth},popmuted] (6.35,0.35) -- (5.3,0.65);
    \end{tikzpicture}
    \legende{Synapse neuromusculaire (jonction cholinergique) en coupe : le bouton terminal du motoneurone contient les vésicules d'acétylcholine et les mitochondries ; la membrane musculaire postsynaptique, plissée, porte les récepteurs nicotiniques (canaux cationiques ligand-dépendants).}
    \end{popfigure}
    \item Étapes chronologiques de la transmission :
    \begin{enumerate}
        \item L'arrivée du PA provoque la \textbf{dépolarisation} du bouton terminal.
        \item Cette dépolarisation ouvre les \textbf{canaux $\ce{Ca^{2+}}$ tension-dépendants} de la membrane présynaptique.
        \item L'\textbf{influx de $\ce{Ca^{2+}}$} (entrée de calcium, porté par son gradient électrochimique) déclenche la fusion des vésicules avec la membrane présynaptique : \textbf{exocytose} de l'acétylcholine dans la fente synaptique.
        \item L'ACh \textbf{diffuse} dans la fente (quelques dizaines de nanomètres, en moins de $\SI{0,1}{ms}$).
        \item L'ACh se \textbf{fixe sur les récepteurs nicotiniques} postsynaptiques (deux molécules d'ACh par récepteur).
        \item Cette fixation provoque l'\textbf{ouverture des canaux cationiques ligand-dépendants} perméables à $\ce{Na+}$ et $\ce{K+}$ ; l'entrée nette de $\ce{Na+}$ l'emporte.
        \item Il en résulte une \textbf{dépolarisation postsynaptique} locale : le PPSE (ou potentiel de plaque motrice), qui, s'il atteint le seuil, déclenche un PA musculaire puis la contraction.
        \item L'ACh est rapidement hydrolysée par l'acétylcholinestérase, ce qui termine le signal.
    \end{enumerate}
    \item Nature ionique du PPSE : les canaux nicotiniques ouverts laissent passer simultanément $\ce{Na+}$ (entrée) et $\ce{K+}$ (sortie), avec un potentiel d'inversion voisin de $\SI{0}{mV}$. Au potentiel de repos ($-\SI{90}{mV}$ pour la fibre musculaire), la force électromotrice sur $\ce{Na+}$ est bien plus grande que celle sur $\ce{K+}$ : le courant net est donc \textbf{entrant et dépolarisant}. C'est pourquoi le PPSE de la plaque motrice est toujours excitateur (dépolarisant), contrairement aux synapses sur récepteurs $\ce{Cl-}$ ou $\ce{K+}$ qui peuvent être inhibitrices.
\end{enumerate}
\end{corrige}

\begin{corrige}[Exercice 8 — Effets de substances pharmacologiques sur la transmission neuromusculaire]
\begin{enumerate}
    \item \textbf{Curare} : c'est un \cle{antagoniste compétitif} des récepteurs nicotiniques de l'ACh. Il se fixe de façon réversible sur le site de liaison de l'ACh sans activer le canal, et empêche l'ACh de s'y fixer. Conséquence : l'ouverture des canaux cationiques postsynaptiques diminue, le PPSE devient sous-liminaire, aucun PA musculaire n'est déclenché et la contraction s'abolit progressivement (paralysie flasque). Le caractère \textbf{réversible} de la fixation explique la récupération après lavage : le curare se dissocie des récepteurs et est éliminé du bain.
    \item \textbf{Néostigmine} : inhibiteur réversible de l'acétylcholinestérase.
    \begin{itemize}
        \item À dose modérée : l'ACh libérée n'est plus dégradée rapidement ; sa concentration dans la fente augmente et sa durée d'action se prolonge. Chaque stimulation produit un PPSE plus ample et plus durable, d'où une contraction \textbf{renforcée}.
        \item À forte dose : l'accumulation massive d'ACh maintient les canaux nicotiniques ouverts en permanence ; la membrane postsynaptique reste \textbf{dépolarisée de façon durable}. Or une dépolarisation prolongée \cle{inactive} les canaux $\ce{Na+}$ tension-dépendants voisins : plus aucun PA ne peut être régénéré. C'est le \textbf{bloc dépolarisant} (paralysie par excitation excessive), effet paradoxal d'un « excitateur ».
    \end{itemize}
    \item \textbf{Toxine botulique} : en bloquant l'exocytose des vésicules d'ACh (clivage des protéines SNARE présynaptiques), elle supprime la libération du neurotransmetteur : plus de PPSE, donc \textbf{paralysie flasque} du muscle (blocage présynaptique, irréversible jusqu'à repousse terminale). Application thérapeutique (hors cosmétique) : traitement du \textbf{blépharospasme} et des \textbf{dystonies} (contractions musculaires involontaires), du \textbf{strabisme}, de la \textbf{spasticité} post-AVC ou des migraines chroniques, par injection locale à très faible dose afin de relâcher sélectivement les muscles hyperactifs.
\end{enumerate}
\end{corrige}

\courssub{Sujets type Baccalauréat — Synthèse et argumentation}

\begin{corrige}[Exercice 9 — Sujet type Bac : Potentiel d'action, conduction et myélinisation]
\textbf{1. Analyse du potentiel d'action (Doc 1)}
\begin{enumerate}
    \item[a)] Courbe $V_m = f(t)$ :
    \begin{popfigure}
    \begin{tikzpicture}
    \begin{axis}[
        width=0.85\linewidth, height=6cm,
        xlabel={$t$ (ms)}, ylabel={$V_m$ (mV)},
        xmin=0, xmax=1.5, ymin=-100, ymax=60,
        grid=both, grid style={popline!40},
        xtick={0,0.2,0.4,0.6,0.8,1.0,1.2,1.4},
        ytick={-80,-70,0,35},
    ]
        \addplot[popcurve,smooth,mark=*,mark size=1.5pt] coordinates {
            (0,-70) (0.2,-65) (0.4,20) (0.6,35) (0.8,10) (1.0,-60) (1.2,-80) (1.4,-72)
        };
        \addplot[dashed,popmuted,domain=0:1.5] {-70};
        \node[popmuted,font=\small] at (axis cs:1.35,-62) {repos};
    \end{axis}
    \end{tikzpicture}
    \legende{Potentiel d'action enregistré sur une fibre A myélinisée de grenouille (échelle : 1 cm pour 0,2 ms en abscisse ; 1 cm pour 20 mV en ordonnée).}
    \end{popfigure}
    \item[b)] Lectures : $V_{\text{repos}} = -\SI{70}{mV}$ ; amplitude $A = 35 - (-70) = \SI{105}{mV}$ ; durée du PA à la base (du début de la montée au retour à $-\SI{70}{mV}$) : environ $\SI{1,4}{ms}$ ; durée de la dépolarisation (de $-\SI{55}{mV}$ au pic) : environ $0,3$ à $\SI{0,4}{ms}$.
    \item[c)] Phases :
    \begin{itemize}
        \item \textbf{Dépolarisation} ($0,2$--$\SI{0,6}{ms}$) : entrée massive de $\ce{Na+}$ (canaux $\ce{Na+}$ tension-dépendants ouverts).
        \item \textbf{Repolarisation} ($0,6$--$\SI{1,0}{ms}$) : sortie de $\ce{K+}$ (canaux $\ce{K+}$ ouverts, canaux $\ce{Na+}$ inactivés).
        \item \textbf{Hyperpolarisation} ($1,0$--$\SI{1,4}{ms}$, jusqu'à $-\SI{80}{mV}$) : sortie persistante de $\ce{K+}$ (fermeture lente des canaux $\ce{K+}$).
    \end{itemize}
\end{enumerate}
\textbf{2. Vitesse de conduction (Doc 2)}
\begin{enumerate}
    \item[a)] Fibres A : $v_A = \dfrac{d}{\Delta t_A} = \dfrac{\SI{0,100}{m}}{\SI{3,3e-3}{s}} \approx \SI{30}{m\cdot s^{-1}}$. \quad Fibres C : $v_C = \dfrac{\SI{0,100}{m}}{\SI{67e-3}{s}} \approx \SI{1,5}{m\cdot s^{-1}}$.
    \item[b)] Rapport : $\dfrac{v_A}{v_C} = \dfrac{30}{1,5} = 20$. La myélinisation multiplie la vitesse par un facteur 20 ici : pour la grenouille prédatrice, cela permet des réflexes de capture (détente linguale) et de fuite (saut) quasi instantanés, donc un avantage sélectif majeur en termes de survie et de prédation.
\end{enumerate}
\textbf{3. Périodes réfractaires (Doc 3)}
\begin{enumerate}
    \item[a)] \textbf{PRA} : intervalle après un PA pendant lequel aucun stimulus, même supramaximal, ne peut déclencher un second PA. \textbf{PRR} : intervalle pendant lequel seul un stimulus plus intense que le seuil peut déclencher un PA, d'amplitude souvent réduite.
    \item[b)] D'après les résultats : PRA $\approx$ entre $0,5$ et $\SI{2,0}{ms}$ (à $\SI{0,5}{ms}$ aucun second PA ; à $\SI{2,0}{ms}$ un PA apparaît) ; PRR $\approx$ de $2,0$ à $\SI{5,0}{ms}$ (PA d'amplitude réduite à $\SI{2}{ms}$, PA normal à $\SI{5}{ms}$).
    \item[c)] Base biophysique : pendant la PRA, les canaux $\ce{Na+}$ sont \textbf{inactivés} (bille d'inactivation en place) et ne peuvent pas être rouverts. Pendant la PRR, une partie des canaux $\ce{Na+}$ est désinactivée, mais la membrane est encore \textbf{hyperpolarisée} et les canaux $\ce{K+}$ restent partiellement ouverts : il faut un stimulus plus fort pour atteindre le seuil, et le PA obtenu est d'amplitude moindre.
\end{enumerate}
\textbf{4. Synthèse.} Dans une fibre myélinisée, la gaine isole l'internœud (faible capacité, forte résistance transversale) : le courant se propage passivement et rapidement jusqu'au nœud suivant, et le PA n'est régénéré qu'aux nœuds de Ranvier, espacés de $1$--$\SI{2}{mm}$. Les échanges ioniques (entrées de $\ce{Na+}$, sorties de $\ce{K+}$) ne concernent donc qu'une surface membranaire minuscule (les nœuds) ; la pompe $\ce{Na+/K+}$ n'a à rétablir les gradients que pour de petites quantités d'ions, d'où une consommation d'ATP réduite. Dans une fibre amyélinique, le PA doit être régénéré sur \emph{toute} la surface de la membrane, pas à pas : propagation lente (conduction continue) et flux ioniques massifs sur toute la longueur, donc coût énergétique élevé. La myélinisation concilie ainsi vitesse élevée et économie d'énergie.

\textbf{Barème indicatif (sur 20)} : Q1a 2 pts ; Q1b 2 pts ; Q1c 3 pts ; Q2a 2 pts ; Q2b 2 pts ; Q3a 2 pts ; Q3b 2 pts ; Q3c 2 pts ; Q4 3 pts.

\textbf{Points de vigilance} : citer explicitement les unités et les conversions ($\SI{10}{cm} = \SI{0,100}{m}$) ; nommer les ions \emph{et} leur sens de déplacement ; pour la PRA, le mot « inactivés » est exigé ; pour la synthèse, articuler vitesse (isolation, nœuds) et énergie (surface d'échange réduite, pompe $\ce{Na+/K+}$).
\end{corrige}

\begin{corrige}[Exercice 10 — Sujet type Bac : Preuve expérimentale de la transmission chimique et pharmacologie]
\textbf{1. Analyse de l'expérience de Loewi (Doc 1)}
\begin{enumerate}
    \item[a)] Le cœur receveur est \textbf{dénervé} : aucune connexion nerveuse ne peut transmettre le signal. Pourtant il ralentit après la stimulation du vague du donneur, via le seul liquide de perfusion commun. Conclusion fondamentale : la transmission du message nerveux au cœur s'effectue par l'intermédiaire d'une \textbf{substance chimique diffusible} libérée par les terminaisons nerveuses (le « Vagusstoff » de Loewi), et non par un courant électrique direct.
    \item[b)] Le filtre à pores $< \SI{0,2}{\mu m}$ retient toute cellule ou particule de taille bactérienne, mais laisse passer les petites molécules dissoutes. Le ralentissement persistant du receveur en expérience B montre que le messager est une \textbf{molécule de petite taille}, soluble, et non un microorganisme ou un fragment cellulaire : cela élimine l'hypothèse d'un transfert cellulaire et confirme la nature chimique (moléculaire) du messager.
    \item[c)] L'expérience C montre que le liquide de perfusion, \emph{prélevé puis appliqué ultérieurement} sur un cœur isolé, conserve son activité : le messager chimique est \textbf{stable} pendant un certain temps dans le milieu extracellulaire, transportable et actif à distance de son site de libération. C'est la preuve définitive de l'existence d'un médiateur chimique (l'ACh).
\end{enumerate}
\textbf{2. Mécanismes moléculaires (Doc 2)}
\begin{enumerate}
    \item[a)] \textbf{Non} : les récepteurs muscariniques sont des récepteurs \textbf{métabotropiques} couplés à une protéine G (GPCR), et non des canaux ioniques ligand-dépendants (ionotropes). La fixation de l'ACh active la protéine G, dont les sous-unités modulent en seconde intention des canaux ioniques (ouverture de canaux $\ce{K+}$, d'où l'hyperpolarisation et la bradycardie) : la réponse est plus lente et plus durable que celle des récepteurs nicotiniques (canaux) de la plaque motrice.
    \item[b)] \textbf{Atropine} : antagoniste compétitif des récepteurs muscariniques ; elle occupe le site de fixation de l'ACh sans l'activer, bloquant l'effet vagal (le cœur ne ralentit plus, voire s'accélère). \textbf{Physostigmine} : inhibiteur de l'acétylcholinestérase ; l'ACh libérée n'est plus hydrolysée, sa concentration et sa durée d'action dans la fente augmentent : l'effet bradycardisant de la stimulation vagale est \textbf{potentialisé et prolongé}.
\end{enumerate}
\textbf{3. Pathologie et traitement (Doc 3)}
\begin{enumerate}
    \item[a)] Dans la myasthénie grave, des auto-anticorps se fixent sur les récepteurs nicotiniques de la plaque motrice : ils les bloquent, accélèrent leur internalisation et détruisent les replis postsynaptiques. Le nombre de récepteurs fonctionnels diminue, donc le PPSE devient plus petit et parfois sous-liminaire : la transmission neuromusculaire est défaillante. À l'effort répété, la quantité d'ACh libérée par impulsion décroît légèrement (épuisement des vésicules immédiatement disponibles) ; avec un nombre de récepteurs déjà réduit, le PPSE tombe rapidement sous le seuil : d'où la \textbf{fatigabilité} musculaire croissante à l'effort.
    \item[b)] Les inhibiteurs de l'AChE (pyridostigmine) empêchent l'hydrolyse de l'ACh : sa \textbf{concentration dans la fente synaptique augmente} et sa durée de vie se prolonge. Chaque molécule d'ACh a ainsi plus de chances de rencontrer l'un des rares récepteurs encore fonctionnels : le PPSE retrouve une \textbf{amplitude suffisante} pour déclencher le PA musculaire. Le traitement est symptomatique (il ne supprime pas les auto-anticorps) mais restaure l'efficacité de la transmission.
\end{enumerate}
\textbf{Barème indicatif (sur 20)} : Q1a 2 pts ; Q1b 2 pts ; Q1c 2 pts ; Q2a 3 pts ; Q2b 3 pts ; Q3a 4 pts ; Q3b 4 pts.

\textbf{Points de vigilance} : insister sur le fait que le cœur receveur est \emph{dénervé} (condition essentielle du raisonnement) ; distinguer clairement récepteur ionotrope (canal) et métabotrope (protéine G) ; pour Q3a, articuler obligatoirement « diminution du nombre de récepteurs $\Rightarrow$ PPSE sous-liminaire $\Rightarrow$ fatigabilité ».
\end{corrige}

\begin{corrige}[Exercice 11 — Sujet type Bac : Intégration — Sclérose en plaques, conduction et potentiel d'action]
\textbf{1. Physiopathologie (Doc 1)}
\begin{enumerate}
    \item[a)] La myéline, empilement de lamelles membranaires lipidiques, joue le rôle d'isolant électrique : elle \textbf{diminue fortement la capacité membranaire} $C_m$ (les lamelles éloignent les plans de charges) et \textbf{augmente la résistance transversale} $R_m$ (les ions ne peuvent plus traverser la membrane dans l'internœud). Après démyélinisation, la membrane axoplasmique est exposée « à nu » sur plusieurs millimètres : toute cette surface redevant capacitive et perméable, la capacitance effective $C_m'$ est massivement augmentée (ici de $\SI{2}{nF}$ à $\SI{200}{nF}$, soit un facteur 100).
    \item[b)] Pour dépolariser la membrane jusqu'au seuil, il faut charger cette capacité ($Q = C_m \Delta V$) : avec $C_m'$ centuplée, le courant local issu du segment actif est largement « absorbé » par la charge de la membrane nue et fuit à travers elle ($R_m$ effondrée). Le courant disponible pour dépolariser le segment suivant devient insuffisant : le \textbf{seuil d'excitabilité apparent augmente}, la conduction ralentit et peut se \textbf{bloquer} (bloc de conduction), surtout lors de stimulations répétées (fatigabilité), car les pompes et gradients ioniques sont mis à rude épreuve.
\end{enumerate}
\textbf{2. Analyse des PEV (Doc 2)}
\begin{enumerate}
    \item[a)] Temps de conduction sur $\SI{0,05}{m}$ :
    \[ t_m = \frac{L}{v_m} = \frac{\SI{0,05}{m}}{\SI{50}{m\cdot s^{-1}}} = \SI{1,0e-3}{s} = \SI{1}{ms} \qquad ; \qquad t_d = \frac{L}{v_d} = \frac{\SI{0,05}{m}}{\SI{1}{m\cdot s^{-1}}} = \SI{50}{ms}. \]
    \item[b)] La différence de latence mesurée entre poussée et témoin est de $135 - 100 = \SI{35}{ms}$, du même ordre de grandeur que le retard théorique calculé ($\approx \SI{49}{ms}$) : la démyélinisation du nerf optique explique l'essentiel de l'allongement. Cependant, la latence P100 ($\SI{100}{ms}$ chez le témoin) est très supérieure au seul temps de conduction axonale ($\approx \SI{1}{ms}$) car elle inclut : la \textbf{transduction photoréceptrice} rétinienne, les \textbf{délais synaptiques} (relais rétine, corps genouillé latéral, cortex), la conduction dans les fibres amyéliniques rétiniennes et le temps d'intégration corticale.
    \item[c)] En rémission, la latence diminue ($135 \to \SI{115}{ms}$) et l'amplitude se normalise presque : cela traduit une \textbf{rémyélinisation partielle} par les oligodendrocytes et une \textbf{redistribution des canaux $\ce{Na+}$} le long de l'axone démyélinisé, qui restaure une conduction continue lente mais fonctionnelle. La latence reste supérieure à la normale car la rémyélinisation est incomplète (internœuds plus courts, gaine plus fine) ; la persistance d'une amplitude légèrement réduite reflète la perte axonale résiduelle.
\end{enumerate}
\textbf{3. Énergétique et modélisation (Doc 3)}
\begin{enumerate}
    \item[a)] La charge de sodium entrant par PA est proportionnelle à la capacité : $Q = C_m \Delta V$. Pour un même $\Delta V \approx \SI{100}{mV}$, le rapport des charges est :
    \[ \frac{Q_d}{Q_m} = \frac{C_m'}{C_m} = \frac{\SI{200}{nF}}{\SI{2}{nF}} = 100. \]
    Sur une zone démyélinisée, l'influx de $\ce{Na+}$ par unité de longueur est donc environ \textbf{100 fois plus important} qu'au niveau d'un nœud de Ranvier. Pour rétablir les gradients, la pompe $\ce{Na+/K+}$ doit extruder ces $\ce{Na+}$ au prix d'1 ATP pour 3 $\ce{Na+}$ : la consommation d'ATP est multipliée dans le même rapport (environ 100).
    \item[b)] Chez le patient SEP, chaque influx nerveux dans les voies démyélinisées coûte beaucoup plus d'ATP. Lors d'activités soutenues, la demande énergétique locale dépasse l'apport (métabolisme oxydatif, perfusion), les gradients ioniques se dissipent partiellement, les pompes ralentissent et la conduction devient défaillante : les signaux sont mal transmis, ce qui se traduit par une \textbf{fatigue centrale majeure} (asthénie), aggravée par l'effort et la chaleur (phénomène d'Uhthoff).
\end{enumerate}
\textbf{4. Traitement symptomatique.} La fampridine bloque les canaux $\ce{K+}$ tension-dépendants : la repolarisation est ralentie, le PA est \textbf{élargi en durée}. Un PA plus long injecte davantage de courant dans l'axoplasme et pendant plus longtemps : le courant longitudinal peut alors charger la forte capacité de la membrane démyélinisée et atteindre le seuil au segment suivant, restaurant une conduction jusque-là bloquée. Indication clinique validée : l'\textbf{amélioration de la marche} (vitesse de marche) chez les patients SEP avec troubles de la marche.

\textbf{Barème indicatif (sur 20)} : Q1a 2 pts ; Q1b 2 pts ; Q2a 2 pts ; Q2b 3 pts ; Q2c 2 pts ; Q3a 3 pts ; Q3b 2 pts ; Q4 4 pts.

\textbf{Points de vigilance} : exiger la relation $Q = C_m \Delta V$ écrite explicitement et le rapport calculé (facteur 100) ; ne pas confondre latence P100 (délai global incluant transduction et synapses) et temps de conduction axonale pur ; pour la fampridine, la chaîne causale « blocage canaux $\ce{K+}$ $\Rightarrow$ PA élargi $\Rightarrow$ courant axoplasmique augmenté $\Rightarrow$ franchissement du bloc » doit être complète.
\end{corrige}

\newpage

\coursec{Fiche bilan}

\begin{popfigure}
\centering
\begin{tikzpicture}[mindmap, grow cyclic,
  every node/.style={concept, font=\small},
  root concept/.append style={concept color=popblue, font=\bfseries\small, minimum size=2.6cm},
  level 1 concept/.append style={level distance=3.6cm, sibling angle=60, font=\scriptsize},
  concept/.append style={text width=2.3cm, align=center}]
\node[root concept]{Le fonctionnement des neurones}
  child[concept color=popgreen]{ node[align=center]{Tissu nerveux et neurone\\ corps cellulaire, dendrites, axone, myéline} }
  child[concept color=poporange]{ node[align=center]{Potentiel de repos\\ $-70$ mV ; pompe Na$^+$/K$^+$ ; fuites K$^+$} }
  child[concept color=poppink]{ node[align=center]{Potentiel d'action\\ seuil, dépolarisation Na$^+$, repolarisation K$^+$, tout ou rien} }
  child[concept color=poppurple]{ node[align=center]{Conduction\\ courants locaux ; sautatoire si myéline ; vitesse 1 à 120 m/s} }
  child[concept color=popgold]{ node[align=center]{Synapse chimique\\ vésicules, neurotransmetteur, récepteurs, PPS} }
  child[concept color=popblue!70]{ node[align=center]{Substances\\ curare, atropine, organophosphorés, caféine} };
\end{tikzpicture}
\legende{Carte mentale de la leçon : les six idées-forces du fonctionnement des neurones.}
\end{popfigure}

\begin{bilan}
\textbf{1. Définitions clés}
\begin{itemize}
  \item \cle{Neurone} : cellule nerveuse spécialisée comprenant un \textbf{corps cellulaire}, des \textbf{dendrites} (réception) et un \textbf{axone} (émission), parfois entouré d'une gaine de \textbf{myéline} (cellules de Schwann) interrompue aux \textbf{nœuds de Ranvier}.
  \item \cle{Potentiel de repos (PR)} : différence de potentiel transmembranaire d'un neurone au repos, d'environ $\SI{-70}{mV}$ (intérieur négatif), due à la \textbf{pompe Na$^+$/K$^+$-ATPase} (3 Na$^+$ sortent, 2 K$^+$ entrent) et aux \textbf{fuites de K$^+$}.
  \item \cle{Potentiel d'action (PA)} : inversion brève et stéréotypée du potentiel de membrane (de $-70$ mV à environ $+30$ mV) déclenchée quand la stimulation atteint le \textbf{seuil} ($\approx -55$ mV) ; il obéit à la \textbf{loi du tout ou rien}.
  \item \cle{Période réfractaire} : intervalle (1 à 2 ms) pendant lequel la fibre ne peut pas émettre un nouveau PA ; elle impose le caractère discontinu du message et sa propagation unidirectionnelle.
  \item \cle{Synapse} : zone de jonction entre un neurone et une cellule cible ; elle est \textbf{chimique} (neurotransmetteur) ou \textbf{électrique} (jonctions communicantes).
  \item \cle{Neurotransmetteur} : messager chimique (acétylcholine, glutamate, GABA\ldots) libéré par exocytose dans la fente synaptique et fixé sur des récepteurs postsynaptiques spécifiques.
\end{itemize}

\textbf{2. Valeurs et mécanismes à connaître par cœur}
\begin{center}
\begin{tabularx}{\linewidth}{|l|X|}
\hline
\rowcolor{popblueL} \textbf{Grandeur / étape} & \textbf{Valeur / mécanisme} \\
\hline
Potentiel de repos & $\approx \SI{-70}{mV}$ \\
\hline
Seuil de déclenchement du PA & $\approx \SI{-55}{mV}$ \\
\hline
Sommet du PA & $\approx \SI{+30}{mV}$ \\
\hline
Dépolarisation & Entrée massive de Na$^+$ (canaux voltage-dépendants) \\
\hline
Repolarisation & Sortie de K$^+$ ; hyperpolarisation transitoire \\
\hline
Vitesse de conduction & $\approx 1$ m/s (fibre amyélinisée) jusqu'à $\approx 120$ m/s (fibre myélinisée) ; augmente avec le diamètre et la myélinisation \\
\hline
Conduction saltatoire & Le PA « saute » de nœud de Ranvier en nœud de Ranvier \\
\hline
Transmission synaptique & PA $\rightarrow$ entrée de Ca$^{2+}$ $\rightarrow$ exocytose $\rightarrow$ fixation sur récepteurs $\rightarrow$ PPS $\rightarrow$ destruction/recapture du médiateur \\
\hline
\end{tabularx}
\end{center}

\textbf{3. Méthodes express}
\begin{itemize}
  \item \textbf{Mesure des potentiels} : microélectrode intracellulaire + électrode de référence extracellulaire reliées à un oscilloscope (ou voltmètre) ; stimuler avec des électrodes d'intensité croissante.
  \item \textbf{Lire une courbe de PA} : repérer PR $\rightarrow$ dépolarisation $\rightarrow$ sommet $\rightarrow$ repolarisation $\rightarrow$ hyperpolarisation $\rightarrow$ retour au PR ; identifier les mouvements ioniques à chaque phase.
  \item \textbf{Raisonner sur une substance} : se demander si elle bloque les récepteurs (curare), détruit l'enzyme de dégradation (organophosphorés $\Rightarrow$ accumulation d'ACh) ou module la libération du médiateur.
\end{itemize}
\end{bilan}

\begin{aretenir}[Checklist avant le Bac]
\begin{itemize}
  \item[$\square$] Je sais schématiser et légender un neurone (corps cellulaire, dendrites, axone, myéline, nœuds de Ranvier).
  \item[$\square$] Je définis le potentiel de repos et je justifie sa valeur ($-70$ mV) par la pompe Na$^+$/K$^+$ et les fuites de K$^+$.
  \item[$\square$] Je décris le protocole de mesure des potentiels de repos et d'action (microélectrodes, oscilloscope).
  \item[$\square$] Je trace et j'annote la courbe du potentiel d'action (seuil, dépolarisation, repolarisation, hyperpolarisation).
  \item[$\square$] J'associe chaque phase du PA aux mouvements de Na$^+$ et de K$^+$ à travers la membrane.
  \item[$\square$] J'énonce la loi du tout ou rien et le rôle de la période réfractaire.
  \item[$\square$] J'explique la conduction saltatoire et je cite deux facteurs de la vitesse de propagation (myéline, diamètre).
  \item[$\square$] Je distingue synapse chimique et synapse électrique.
  \item[$\square$] Je schématise une synapse chimique et je décris les étapes de la transmission (Ca$^{2+}$, exocytose, récepteurs, PPS, recapture).
  \item[$\square$] J'explique le mode d'action d'au moins deux substances (curare, organophosphorés, atropine, caféine).
  \item[$\square$] Je respecte les unités (mV, ms, m/s) et je rédige mes interprétations en reliant observation, hypothèse et conclusion.
\end{itemize}
\end{aretenir}

\findecours

\end{document}
